% SUPERSEDED: historical audit narrative; do not cite as current results.
\documentclass[11pt,journal,onecolumn]{IEEEtran}

% --- Font encoding: REQUIRED, do not remove ---------------------------------
% Without T1 + lmodern this document builds with the OT1 Computer Modern setup,
% under which XeTeX silently falls back to the regular face for every bold,
% italic and small-caps request. Every \textbf run-in lead, every \textit, the
% section headings' own styling and IEEEtran's "Abstract"/"Index Terms" labels
% then render as ordinary roman -- the file compiles clean and looks fine, so
% the loss is easy to miss. Verified 2026-08-26.
\usepackage[T1]{fontenc}
\usepackage{lmodern}

\usepackage{amsmath}
\usepackage{amssymb}
\usepackage{booktabs}
\usepackage{graphicx}
\usepackage{hyperref}
\usepackage{cite}
\usepackage{microtype}
\usepackage{algorithmic}
\usepackage{algorithm}
\usepackage{array}
\usepackage{multirow}
\usepackage{etoolbox}

% Long \texttt{} model ids, file paths and JSON field names cannot hyphenate by
% default and were producing overfull lines; allow breaks inside typewriter text
% and give TeX a last-resort stretch budget for the remaining tight paragraphs.
\appto\ttfamily{\hyphenchar\font=`\-}
\setlength{\emergencystretch}{3em}

% --- Bold headings ---------------------------------------------------------
% IEEEtran sets sections in small caps and every level below in italic, with no
% bold anywhere. These redefinitions keep IEEEtran's exact spacing, numbering
% and run-in behaviour (the parameters are copied verbatim from the class) and
% change only the font, giving a four-level ladder:
%   section        bold, centred
%   subsection     bold, flush left
%   subsubsection  bold italic, run-in
%   paragraph      italic, run-in, deeper indent
% Note: section is NOT \bfseries\scshape. Latin Modern has no bold small-caps
% shape (T1/lmr/bx/sc is undefined), so that combination silently drops one of
% the two attributes rather than compounding them.
% Delete this block to return to stock IEEEtran styling.
\makeatletter
\renewcommand\section{\@startsection{section}{1}{\z@}%
  {3.0ex plus 1.5ex minus 1.5ex}{0.7ex plus 1ex minus 0ex}%
  {\normalfont\normalsize\centering\bfseries}}
\renewcommand\subsection{\@startsection{subsection}{2}{\z@}%
  {3.5ex plus 1.5ex minus 1.5ex}{0.7ex plus .5ex minus 0ex}%
  {\normalfont\normalsize\bfseries}}
\renewcommand\subsubsection{\@startsection{subsubsection}{3}{\parindent}%
  {0ex plus 0.1ex minus 0.1ex}{0ex}%
  {\normalfont\normalsize\bfseries\itshape}}
\renewcommand\paragraph{\@startsection{paragraph}{4}{2\parindent}%
  {0ex plus 0.1ex minus 0.1ex}{0ex}%
  {\normalfont\normalsize\itshape}}
\makeatother

% --- Replication archive ---------------------------------------------------
% Set \ecnanonymousfalse and fill in \ecnarchivedoi for the camera-ready; the
% Data Availability section and any \ecnarchive reference then resolve to the
% real link automatically. Keep in sync with CITATION.cff in the archive.
\newif\ifecnanonymous
\ecnanonymoustrue
\newcommand{\ecnarchivedoi}{10.5281/zenodo.XXXXXXX}
\newcommand{\ecnarchiveurl}{https://doi.org/\ecnarchivedoi}
\newcommand{\ecnarchive}{%
  \ifecnanonymous
    the replication archive accompanying this submission%
  \else
    the replication archive (\href{\ecnarchiveurl}{\texttt{\ecnarchivedoi}})%
  \fi
}
% --- Float placement -------------------------------------------------------
% This paper carries 56 floats over 67 pages. With the default parameters (and
% the [h]-only specifiers the draft used) LaTeX deferred most of them, which
% piled the appendix tables up at the very end, pages away from their text.
\renewcommand{\topfraction}{0.90}
\renewcommand{\bottomfraction}{0.75}
\renewcommand{\textfraction}{0.08}
\renewcommand{\floatpagefraction}{0.70}
\setcounter{topnumber}{3}
\setcounter{bottomnumber}{2}
\setcounter{totalnumber}{5}

% --- Part dividers --------------------------------------------------------
% The paper is organised in two parts: the audit that relocates the central
% finding, and the one result that survives it. IEEEtran has no \part, so the
% divider is drawn explicitly and kept out of the section numbering.
\newcommand{\parttitle}[2]{%
  \par\addvspace{1.6\baselineskip}%
  \noindent\rule{\textwidth}{0.9pt}\par\vspace{0.45em}%
  \begin{center}
    {\large\bfseries PART #1}\\[0.3em]
    {\large\itshape #2}
  \end{center}%
  \vspace{0.25em}\noindent\rule{\textwidth}{0.9pt}%
  \par\addvspace{1.1\baselineskip}%
  \nopagebreak}

% --- Cross-model comparison warning ---------------------------------------
% No two campaigns share release generation, numerical precision and serving
% stack, so every between-model statement in this paper is confounded. Rather
% than repeat the caveat in prose, it is boxed at the head of each section that
% makes one, so that a reader who skims cannot miss it.
\newcommand{\confoundnote}[1]{%
  \par\addvspace{0.6\baselineskip}%
  \noindent\fbox{\parbox{\dimexpr\textwidth-2\fboxsep-2\fboxrule\relax}{%
    \footnotesize\textbf{Cross-model comparison warning.} #1}}%
  \par\addvspace{0.6\baselineskip}}

\newcommand{\preauditnote}[1]{%
  \par\addvspace{0.6\baselineskip}%
  \noindent\fbox{\parbox{\dimexpr\textwidth-2\fboxsep-2\fboxrule\relax}{%
    \footnotesize\textbf{Pre-audit data. Do not cite these numbers as results.} #1}}%
  \par\addvspace{0.6\baselineskip}}

\hypersetup{
    colorlinks=true,
    linkcolor=blue,
    filecolor=magenta,
    urlcolor=cyan,
    pdftitle={The Evaluator Is the Instrument: A Forensic Audit of Degenerate Consensus in Multi-Agent LLM Forecasting},
    pdfpagemode=FullScreen,
}

\begin{document}

% \bfseries applies inside the title group only; the running head and the PDF
% metadata title are set separately and are unaffected.
\title{\bfseries The Evaluator Is the Instrument: A Forensic Audit of Degenerate Consensus in Multi-Agent LLM Forecasting}

\author{Anonymous Authors
\thanks{Manuscript draft updated August 26, 2026. This research was conducted independently utilizing university High-Performance Computing (HPC) infrastructure. Three complete 30-event campaigns (Llama-3.1-8B-AWQ, Mistral-7B-Instruct-v0.1, Qwen2.5-7B-GGUF) were executed in July 2026 and form the primary empirical basis of the reported results. In August 2026 all 270 simulation transcripts were re-scored with a second, independent evaluator; that replication is reported in Section~\ref{sec:evaluator_replication} and both scored passes are released separately. Two analyses postdate the campaigns and were run against hosted endpoints rather than the cluster: the evaluator-comparison and truncation experiments of Section~\ref{sec:volume_mechanism}, and the single-agent baseline of Section~\ref{sec:cot_baseline} (2026-08-17). Both are marked as such wherever they enter a claim.}
}

\maketitle

\begin{abstract}
Multi-agent LLM social simulations are increasingly proposed for synthesizing calibrated probabilistic forecasts from a swarm of interacting personas---an operationalized ``Wisdom of Crowds''~\cite{surowiecki2005}. This paper is a forensic audit of one such benchmark: it documents a failure mode of that mechanism, then relocates the cause from the swarm to the instrument that measures it. \textbf{ECN-BENCH} elicits forecasts for $E=30$ resolved Polymarket events under direct elicitation and full social simulation ($N=300$ agents, $R=60$ rounds) across three open-weight models. Under simulation, $20.6\%$ of units return an exactly uniform probability vector against $3.3\%$ without it, and mean modal probability falls from $0.664$ to $0.54$--$0.56$ among the remainder. A proper scoring rule charges both effects as miscalibration~\cite{gneiting2007} while directional accuracy is unchanged, so this \textbf{loss of sharpness and degenerate consensus} alone explains the worse simulated scores---which would otherwise read as evidence that deliberation degrades reasoning.

A source-code audit traces the degeneracy to the \emph{extraction instrument} rather than the swarm: a keyword-matching heuristic silently disabled structured-output enforcement for the fixed LLM evaluator (Section~\ref{sec:epistemic_mapping}). Re-scoring the transcripts under a schema-enforced evaluator cuts the flat-forecast rate from $31/132$ to $1/132$ on the simulated units both evaluators read byte-identically, and the cross-model pattern the paper set out to explain vanishes. Degeneracy rises monotonically from $5\%$ to $65\%$ across terciles of evidence length, and re-running five further evaluators in the \emph{identical} unconstrained condition on the same transcripts localises it: four return nothing uniform at all, but a small reasoning model does ($8/58$ units), and $6$ of its $8$ failures fall on the very transcripts the original instrument failed on ($2.6$ expected by chance, $p = 0.012$). A $2 \times 4$ grid then assigns evidence length and enforcement experimentally, and neither survives: shortening the transcript does not reduce degeneracy but mildly increases it ($7/59$ at $2$k against $1/60$ untruncated, $p = 0.032$), and enforcement pooled over lengths changes nothing ($12/239$ against $11/239$, $p = 1.000$). What decides is which model reads the transcript. Neither pass is a reference measurement---three re-reads of byte-identical input reproduce only $4/30$ units exactly.

Any benchmark extracting probabilities with an LLM must therefore calibrate that extractor as an instrument: verify enforcement at runtime, repeat $k \geq 3$ times per unit, and score with two evaluators. We treat this \emph{relocation}---from a claim about deliberation to a reproducible failure of LLM-based probability extraction---as the paper's central contribution, and report flat-unit rate as its cheapest diagnostic. Remaining findings are context rather than validated results, with one exception (Section~\ref{sec:predictive_signal}): on the $20$ events resolving after every model's release date, the swarm's modal forecast is the outcome in $81.9\%$ of $80$ (model, event) pairs against a $32.6\%$ baseline (exact $p = 3.7\times10^{-22}$), and replacing the injected evidence with topically unrelated material costs $0.112$ Brier $[+0.042, +0.182]$. A single-agent baseline added after the audit bounds what that buys: one chain-of-thought call on the same evidence matches or slightly exceeds the swarm's directional accuracy ($0.834$ pooled) and its Brier, with no arm on which the swarm's advantage is distinguishable from zero (Section~\ref{sec:cot_baseline}).
\end{abstract}

\begin{IEEEkeywords}
Benchmark Forensics, Instrument Calibration, LLM-as-Judge Reliability, Structured-Output Enforcement, Multi-Agent Systems, Social Simulation, Calibration and Sharpness, Proper Scoring Rules, Forecast Degeneration, Large Language Models, Brier Score.
\end{IEEEkeywords}

\section{Introduction}
\label{sec:intro}

Multi-agent LLM simulations are increasingly proposed as a way to convert language models' latent knowledge into calibrated collective judgments---an artificial instantiation of the ``Wisdom of Crowds''~\cite{surowiecki2005}, in which a swarm of independently reasoning personas is expected to average out individual error. Whether this happens is a question about the aggregate probability distribution the swarm returns, not only about whether that distribution points in the right direction. This paper is about a specific way the expectation can fail: iterative multi-agent deliberation can converge on a distribution that is directionally correct but has lost the \emph{sharpness}---the concentration of probability mass on the realized outcome---that a calibrated forecast requires. In the most severe form, deliberation returns no distinguishable answer at all: an exactly uniform probability vector over the outcome space. We call this failure mode \textbf{loss of sharpness and degenerate consensus}. Characterizing it---how often it occurs, what it costs under a proper scoring rule, and why an aggregate calibration metric can misattribute it to a collective reasoning failure that need not have occurred---is where this paper began. It is not where it ends. The failure turns out to belong one level down, not to the deliberating swarm but to the LLM instrument that reads the swarm's transcript and compresses it into a number. This paper is the forensic audit that establishes that relocation, and the instrument-calibration protocol that follows from it.

We arrived at this question while auditing a more conventional benchmark result. \textbf{ECN-BENCH} is a multi-agent social simulation testbed (the MiroFish-Offline pipeline, built on CAMEL-AI OASIS~\cite{oasis2024}) originally designed to test three hypotheses about LLM-based crowd intelligence:
\begin{enumerate}
    \item \textbf{Prediction Lift (Claim 1):} Does structuring LLMs as a multi-agent social simulation outperform direct, single-pass elicitation? We hypothesize that social simulation yields a lower Brier Score (better calibration) than a direct query.
    \item \textbf{Convergence (Claim 2):} Does the multi-agent simulation converge systematically toward a real-world-calibrated consensus across rounds? We model this opinion dynamics process via DeGroot consensus models and quantify it using Jensen-Shannon Divergence (JSD) and variance decay.
    \item \textbf{Susceptibility (Claim 3):} Do individual agents update their beliefs in response to new information rationally (reflecting content-specific reasoning), or do they exhibit mere stimulus-reactivity? We isolate this by injecting controlled relevant breaking news versus structurally identical null-data at a fixed simulation checkpoint.
\end{enumerate}
Running this protocol across three models produced an aggregate result that, read at face value, looks like a clean rejection of the first hypothesis: mean Brier score under simulation exceeded the no-simulation baseline for every model tested. Tracing that result to the per-unit level, rather than accepting the aggregate, is what surfaced the phenomenon this paper reports. A fifth of simulated forecasts were not degraded predictions---they were uniform non-answers---and the surviving, non-degenerate forecasts were themselves measurably flatter than their no-simulation counterparts. Both effects are penalized by any proper scoring rule~\cite{gneiting2007} independently of whether the swarm's underlying judgment was correct, and together they are large enough to produce the negative aggregate result without any change in reasoning quality. That is where the investigation started rather than where it ended, and we report the three original hypotheses---along with the testbed's engineering, taxonomy, and statistical machinery---as the context in which it was found, subject to caveats detailed in Section~\ref{sec:condition_a} and Section~\ref{sec:serving_note} that limit what the original benchmark comparison can support on its own.

The audit that followed changed what this paper claims. Tracing the degeneracy through the code path that produces it shows the flat forecasts are introduced by the \emph{extraction instrument}---the fixed LLM evaluator (DeepSeek-R1-14B)---rather than by the swarm: a keyword-matching heuristic silently disabled structured-output enforcement across all three campaigns (Section~\ref{sec:epistemic_mapping}). Re-scoring all $270$ evaluated units (Llama-3.1-8B, Mistral-7B, Qwen2.5-7B; $E=30$ events each) with a second, schema-enforced evaluator drops the flat rate from $20.6\%$ to $6.1\%$ and dissolves the cross-model pattern entirely (Section~\ref{sec:evaluator_replication}). Within the single campaign where evidence length varies without a confound, degeneracy climbs from $5\%$ to $65\%$ across terciles of input length while those same transcripts return sharp forecasts under four independent re-reads (Section~\ref{sec:volume_mechanism}), identifying input volume as the trigger and the unconstrained evaluator as the locus---neither alone produces the effect. Neither evaluator pass is a reference measurement: re-reading Condition A's byte-identical input three times reproduces only $4/30$ units exactly, with a mean Brier spread of $0.128$ (max $0.715$)---large enough to reverse the sign of one model's headline lift.

We therefore report the testbed's remaining results as context for this discovery rather than as validated benchmark findings. The No-Sim baseline was shared rather than run per model (Section~\ref{sec:condition_a}); all three campaigns are statistically underpowered ($\leq 9\%$ achieved power at $E=30$, Section~\ref{sec:power_limitation}); the models differ simultaneously in generation, precision, and serving stack (Section~\ref{sec:serving_note}); a run-integrity audit finds the reference campaign's backend failed silently on 23 of 30 events while recording itself as complete, and a first attempt to repair those events was itself defective in a way that produced two spurious findings before being caught (Appendix~\ref{sec:outage}); resolution-date stratification localises the aggregate lift to memorisation in that shared reference for two of the three campaigns, but not for the reference model itself (Appendix~\ref{sec:cutoff}); round-level convergence telemetry is $27.7$--$87.7\%$ non-functional, inverting a previously reported convergence result (Section~\ref{sec:jsd_convergence}); and the reported difficulty--lift correlation of $r=0.847$ is an arithmetic artifact of regressing a difference on its own minuend (Appendix~\ref{sec:tautology}). One substantive claim survives the audit: on the $20$ events resolving after every model's release date, the swarm's modal forecast is the realised outcome in $81.9\%$ of $80$ (model, event) pairs against a $32.6\%$ chance baseline (exact $p = 3.7\times10^{-22}$), and it tracks the \emph{content} of what it deliberates over---replacing the injected evidence with topically matched but unrelated material costs $0.112$ Brier $[+0.042, +0.182]$ (Section~\ref{sec:predictive_signal}). A single-agent chain-of-thought baseline, run after the audit on the same evidence, then bounds the first half of that claim: it reproduces the directional accuracy on its own, so only the content-sensitivity contrast is evidence about deliberation rather than about what the base models already know (Section~\ref{sec:cot_baseline}).

Related literature offers competing accounts of why iterated multi-agent LLM interaction might homogenize output: multi-agent debate has been reported to improve factual accuracy through explicit disagreement~\cite{du2023debate}, while sycophantic convergence toward perceived majority or authority opinion has been documented in single-turn settings~\cite{sharma2023sycophancy}. Neither literature reports a mechanism by which prolonged deliberation would degrade a forecast's \emph{sharpness} specifically, as opposed to its accuracy---the distinction is central to what follows, since we find directional accuracy to be largely preserved while sharpness is not.

ECN-BENCH organises its event catalogue by a three-way taxonomy (Social/Electoral, Senate/Binary, Technical/Economic; see Table~\ref{tab:taxonomy}), superseding the $2\times2$ Information-Type $\times$ Resolution-Horizon grid of the original specification, which the seed corpus does not carry as per-event metadata. The benchmark is not designed to measure what LLMs ``know,'' but rather how effectively the multi-agent \emph{process architecture} translates latent knowledge, personas, and dynamic deliberation into calibrated probabilistic forecasts---and, as this paper shows, whether that process preserves the forecast's informativeness at all.

To deploy this benchmark at scale (simulating 3,000 agents per event across dozens of configurations), standard cloud APIs are economically prohibitive and methodologically unstable due to undocumented model updates. ECN-BENCH is designed specifically for university-grade, air-gapped High-Performance Computing (HPC) clusters. This paper details the systematic system engineering and surgical patches (Architecture v5.0 through v5.10) developed to resolve critical bottlenecks, including NFS database locking on shared network storage, Apptainer/Singularity container isolation, outlines dependency resolutions, and multi-GPU server routing protocols.

\textbf{How this paper is organised.} It is a forensic audit report, and its structure follows the audit rather than the original benchmark protocol. Sections~\ref{sec:related}--\ref{sec:confound} describe the testbed and its instrumentation; Section~\ref{sec:full_campaign} then presents the campaign result exactly as the benchmark originally reported it, at face value and before any audit, because that is the artifact under examination. The body then divides in two.

\textbf{Part~I} establishes the instrument failure. Section~\ref{sec:degenerate} characterises the degeneracy phenomenologically; Section~\ref{sec:epistemic_mapping} traces it to a specific defect in the scoring code, read directly from source; Section~\ref{sec:evaluator_replication} tests whether the finding survives a second, schema-enforced evaluator reading the same transcripts (it does not); Section~\ref{sec:four_model_extension} carries out the resulting calibration protocol in full---four models, five independent evaluator reads, contamination controlled; and Section~\ref{sec:volume_mechanism} isolates the interaction between input volume and instrument that produces the failure.

\textbf{Part~II} reports what is left standing once the audit is applied: a single content-sensitive predictive signal, measured on the contamination-clean subset with the calibrated protocol (Section~\ref{sec:predictive_signal}), and a single-agent baseline that bounds how much of it the deliberation is responsible for (Section~\ref{sec:cot_baseline}).

Sections~\ref{sec:discussion} and~\ref{sec:conclusion} discuss limits and conclude. The detailed lookup tables and the long-form audit logs are not in the body: the consolidated errata are in Appendix~\ref{sec:errata_appendix}, the per-event score tables in Appendix~\ref{sec:per_event_tables}, and the run-integrity, contamination and event-structure analyses---which are context for the audit rather than validated benchmark results---in Appendix~\ref{sec:full_campaign_supporting}.

\section{Related Work}
\label{sec:related}

\subsection{Static vs. Dynamic Agent Benchmarking}
Traditional LLM benchmarks evaluate models in isolation. Knowledge benchmarks such as MMLU~\cite{mmlu}, and more recently Humanity's Last Exam (HLE)~\cite{hle2024}, focus on expert-level questions where models are scored on binary accuracy. While highly useful, they represent single-pass reasoning. In contrast, SWE-bench~\cite{swebench} and WebArena~\cite{webarena2023} evaluate models as active agents within a sandbox, executing code or navigating interfaces. ECN-BENCH occupies a unique intersection: it evaluates models as a collective, measuring emergent epidemiological and forecasting dynamics within an agent society.

\subsection{Large-Scale LLM Social Simulations}
The study of multi-agent societies has transitioned from classical agent-based models (ABMs) to generative agent architectures. Foundational work by Park et al.~\cite{park2023generative} demonstrated that LLMs could maintain believable individual memories and social behaviors in virtual towns. To scale these systems, the CAMEL-AI community introduced \textbf{OASIS} (Open Agent Social Interaction Simulations)~\cite{oasis2024}, a framework capable of scaling up to one million agents to study dynamic network properties, recommendation systems, and behavioral cascades. ECN-BENCH utilizes a localized, air-gapped fork of the MiroFish platform~\cite{mirofish2024}, built underneath on CAMEL-AI OASIS, converting a generative social environment into a precise, predictive forecasting instrument.

\subsection{GraphRAG and Graph-Based Agent Memory}
Standard Retrieval-Augmented Generation (RAG) retrieves flat vector chunks, which fails to capture long-term relationship structures or temporal context. Microsoft's GraphRAG~\cite{graphrag2024} demonstrated that building a semantic knowledge graph over source documents allows for query-focused summarization over entire corpora. ECN-BENCH builds on this by implementing \textbf{Graphiti}~\cite{graphiti2024}, an open-source temporal context graph engine coupled with Neo4j. This allows individual agent memories to evolve dynamically, tracking entity associations and fact transitions with explicit temporal metadata (e.g., \texttt{valid\_at} and \texttt{invalid\_at} properties).

\subsection{Deliberative Degeneration and Forecast Sharpness}
\label{sec:related_degeneration}
Proper scoring rules such as the Brier score reward forecasts that are simultaneously calibrated and sharp; Gneiting and Raftery's ``maximize sharpness subject to calibration'' principle~\cite{gneiting2007} formalizes why a hedged-but-correct forecast is penalized relative to a confident-and-correct one, and is the theoretical basis for the mechanical account of our central finding (Section~\ref{sec:degenerate}). Separately, a growing literature examines how iterated multi-agent LLM interaction changes output diversity: multi-agent debate frameworks report that structured disagreement between model instances can improve factual accuracy relative to single-model sampling~\cite{du2023debate}, while work on sycophancy shows that models shift stated positions toward a perceived interlocutor's view even in single-turn settings~\cite{sharma2023sycophancy}. Both literatures speak to whether interaction changes what a model says; neither directly addresses whether it changes how \emph{confidently} the model says it, which is the axis on which we observe degradation. To our knowledge, the concentration of exactly-uniform, non-informative outputs under iterated social-simulation deliberation---rather than under adversarial or debate framings---has not been previously reported as a distinct failure mode with its own measurable rate and its own candidate serving-stack mechanism.


\subsection{Constrained Decoding and Structured Output}
\label{sec:related_constrained}
The defect this paper documents is the absence of grammar-constrained generation, and the repair is its presence, so the constrained-decoding literature is where our contribution sits even though we make no methodological contribution to it. Guided decoding restricts sampling to tokens permitted by a grammar or JSON schema, and is available in the serving stack this benchmark ran on through the outlines backend~\cite{vllm_guided,vllm}. The engineering trade-offs are well documented---schema compilation cost, and instability on quantised weights, which is precisely why the bypass we found was introduced in the first place (Appendix~\ref{sec:schema_bypass}).

What is less discussed, and what this paper supplies a case study for, is what an unconstrained model does \emph{instead} when the task is probability extraction rather than free text. Asked to compress a long, multi-sided, partly contradictory transcript into a distribution over $K$ outcomes, a model with no schema to satisfy has a degenerate option available that a constrained one does not: it can express its uncertainty in prose and then emit the uniform vector, which is the minimax point under a symmetric loss it is not in fact being scored by. Enforcement does not make the model better informed; it forecloses the hedge. The measurement consequence is severe because a proper scoring rule charges that hedge the full uniform loss and cannot distinguish it from a swarm that genuinely reached no consensus.

We are careful about how far this generalises, because our own data bounds it. Section~\ref{sec:volume_mechanism} runs five further evaluators in the identical unconstrained condition on the identical transcripts and finds essentially no degeneracy, so ``unconstrained decoding over long input causes uniform forecasts'' is \emph{not} what we are claiming. The hedge is available to any unconstrained model; whether a given deployed model takes it is an empirical property of that deployment. That is the practical point for benchmark designers: enforcement is cheap insurance against a failure whose incidence you cannot predict from the model card.

\subsection{LLM-as-Judge Reliability and Evaluation Instruments}
\label{sec:related_judge}
A parallel literature treats the LLM evaluator itself as an unreliable instrument rather than a neutral scoring function. Zheng et al.~\cite{zheng2023judging} establish LLM-as-a-judge as a scalable proxy for human preference while documenting systematic biases (position, verbosity, self-enhancement) and quantifying disagreement against human raters. That literature typically audits judge-versus-human agreement on a single pass; it does not, to our knowledge, report judge-versus-judge or pass-versus-pass reproducibility under byte-identical input, which is the axis Section~\ref{sec:evaluator_replication} measures directly. Our source-code-level finding that a keyword-matching heuristic silently disables structured-output enforcement for a specific evaluator model (Section~\ref{sec:epistemic_mapping}) is a concrete instance of a failure mode that judge-reliability critiques generally treat as evaluator model capability rather than evaluator configuration---a distinction that matters because the second is fixable by construction and the first is not. Section~\ref{sec:volume_mechanism} shows the distinction is not merely conceptual here: five other evaluators in the same misconfiguration do not reproduce the failure, so what we document is a property of one deployed extractor rather than a capability limit of unconstrained LLM judges.

\subsection{Benchmark Contamination and Temporal Data Leakage}
\label{sec:related_contamination}
Detecting whether a benchmark's evaluation set has leaked into a model's training data is an active research program largely separate from ours: Oren et al.~\cite{oren2024proving} give a black-box statistical test for contamination based on canonical-ordering likelihood, requiring no access to weights or pretraining data. We do not attempt anything comparable in rigor; Appendix~\ref{sec:contamination} instead uses the coarser, cheaper proxy of each event's resolution date relative to each model's public release date, which only bounds the \emph{possibility} of memorization rather than proving its occurrence. We report this as a limitation, not an equivalent method, and flag Oren et al.'s likelihood-based test as the natural next instrument to apply to ECN-BENCH's own seed corpus.

\subsection{LLM Forecasting and Calibration Benchmarks}
\label{sec:related_forecasting}
ECN-BENCH's Brier-score framing sits inside a growing body of work that evaluates LLMs directly as forecasters against real-world, resolvable questions. Halawi et al.~\cite{halawi2024approaching} build a retrieval-augmented forecasting system that approaches the accuracy of competitive human forecasters on questions drawn from live prediction platforms, and Karger et al.~\cite{karger2024forecastbench} introduce ForecastBench, a continuously refreshed benchmark designed specifically to avoid the kind of resolution-date leakage Appendix~\ref{sec:cutoff} audits after the fact. Both efforts evaluate single-model, single-pass forecasting; neither introduces a multi-agent deliberation stage between elicitation and scoring, which is the axis ECN-BENCH varies and the one on which we report loss of sharpness. Positioned against this literature, our contribution is not a new forecasting benchmark in this sense---it is closer to a finding about what happens to forecast quality, measured by the same Brier-score machinery this literature already uses, when a deliberation process is inserted before the LLM commits to a number.

The local serving stack and the HPC systems engineering (Architecture v5.0--v5.10) that make the following simulations executable on an air-gapped cluster---hybrid vLLM/Ollama routing, multi-GPU server isolation, Apptainer containerization, and NFS database concurrency fixes---are a precondition for the results below rather than part of the central finding, and are reported in full in Appendix~\ref{sec:architecture}.

\section{Experimental and Benchmark Design}
\label{sec:design}

\subsection{Taxonomy and Event Selection}
\label{sec:taxonomy}
ECN-BENCH curates $E=30$ resolved real-world categorical markets from Polymarket. Polymarket is used purely as a curation mechanism (identifying clear, multi-outcome categorical events with verified ground truth) rather than as a baseline price signal. The catalogue grew from the 25 events of specification v3.0 to the 30 events used in all campaigns reported here; references to $N=25$ in earlier drafts describe the superseded specification.

The original specification organised events on a $2\times2$ grid of Information Type $\times$ Resolution Horizon. In practice the seed corpus is labelled by a flat three-way prefix, and all analysis in this paper uses that partition, shown in Table~\ref{tab:taxonomy}. We report the three-way scheme rather than the $2\times2$ grid because it is what the data carries; horizon was never recorded as a per-event field, so the quadrant assignment cannot be reconstructed. Note that ``Senate/Binary'' is defined by question format rather than subject matter, so the partition mixes a topical axis with a structural one---an imperfection we retain for traceability with the released seed files, and which readers should bear in mind when interpreting the taxonomy breakdown of Section~\ref{sec:taxonomy_breakdown}.

\begin{table}[!htbp]
\centering
\caption{ECN-BENCH Event Taxonomy as Realised in the 30-Event Corpus}
\begin{tabular}{llcc}
\toprule
\textbf{Prefix} & \textbf{Category} & \textbf{Events} & \textbf{Option count $K$} \\
\midrule
C & Social/Electoral   & 15 & 3--9 \\
S & Senate/Binary      &  7 & 2--4 \\
T & Technical/Economic &  8 & 2--5 \\
\midrule
  & \textbf{Total}     & \textbf{30} & \textbf{2--9} \\
\bottomrule
\end{tabular}
\label{tab:taxonomy}
\end{table}

\subsection{Experimental Conditions}
To isolate the predictive value of multi-agent social interaction, each model is evaluated under three distinct conditions:
\begin{itemize}
    \item \textbf{Condition A (No-Sim Baseline):} The model is presented with the Seed Document and directly queried via the ReportAgent to produce a categorical probability forecast. This represents the raw knowledge baseline.
    \item \textbf{Condition B (With-Sim Benchmark):} A 3,000-agent MiroFish-Offline simulation is initialized. Personas read the Seed Document and participate in a 60-step social network simulation. At step 30, a \textbf{Content-Specific Relevant Injection} (breaking news related to the event) is introduced. The final forecast is extracted at step 60.
    \item \textbf{Condition C (Null Control Simulation):} Identical to Condition B, but the perturbation injected at step 30 is a \textbf{Null Injection} (structurally identical, length-matched, but entirely unrelated content). This serves as the control to measure information susceptibility against general noise susceptibility.
\end{itemize}

\subsection{Simulation Output Fidelity: Action Density and Scope Conditions}
\label{sec:action_fidelity}

As a design-transparency note, we audited social-action density by directly querying the authoritative SQLite \texttt{trace} tables across 1,380 database files on the HPC cluster. We note that the \texttt{actions.jsonl} summary files---used in earlier drafts of this analysis---systematically undercount organic social actions due to a cursor-tracking bug: during telemetry-probe rounds (12, 24, 36, 48, 60), the probe function advances the \texttt{last\_rowid} cursor past all new rows but only logs \texttt{INTERVIEW} actions to the jsonl, silently discarding organic actions written during those rounds. The \texttt{.db} files retain all actions and are the authoritative source. Table~\ref{tab:action_density_30event} reports the exact aggregate; Table~\ref{tab:action_breakdown} reports the per-type breakdown.

\begin{table}[!htbp]
\centering
\caption{Authoritative Social-Action Counts from SQLite Trace Tables (1,380 Database Files, All Conditions B+C, Both Platforms). The \texttt{actions.jsonl} summary files undercount organic actions due to a cursor-tracking bug during telemetry-probe rounds (see text); the \texttt{.db} files are the authoritative source.}
\label{tab:action_density_30event}
\small
\begin{tabular}{lcccc}
\toprule
\textbf{Model} & \textbf{DBs} & \textbf{Total Social} & \textbf{Avg/DB} & \textbf{Agent Traces} \\
 & \textbf{scanned} & \textbf{Actions} & & \\
\midrule
Llama-3.1-8B-AWQ  & 672 & \textbf{27,087} & \textbf{40.31} & 554,782 \\
Mistral-7B-Instruct-v0.1    & 264 & 9,364           & 35.47          & 257,339 \\
Qwen2.5-14B   & 180 & 2,407           & 13.38          & 113,366 \\
Qwen2.5-7B-GGUF   & 264 & 723             & 2.74           & 191,942 \\
\bottomrule
\end{tabular}
\smallskip \\
\footnotesize ``Social Actions'' = post + comment + like + comment\_like + follow + dislike + comment\_dislike (excluding telemetry probes and refresh/sign\_up). ``Agent Traces'' = all trace-table rows including telemetry. Llama-3.1-8B's 672 DBs include the main July campaign (180) plus optimization sweeps and earlier runs; per-DB averaging normalizes across unequal DB counts. Qwen2.5-14B's 180 DBs = exactly one full 30-event campaign (30 events $\times$ 3 conditions $\times$ 2 platforms). \textbf{Reproducibility caveat:} unlike every other quantity in this paper, these counts were obtained by SQL queries executed on the cluster and are not recoverable from the released \texttt{event\_results.json} artifacts, which do not carry action counts. They could not be re-verified during the audit described in Appendix~\ref{sec:data_audit} and are reported on the authority of the original scan. Exporting per-unit action counts into the released summaries is a required change for the next campaign. Note also that the Llama-3.1-8B row pools heterogeneous configurations (different $N$ and $R$ across sweeps), so its per-DB average is not directly comparable with rows drawn from a single fixed configuration.
\end{table}

\begin{table}[!htbp]
\centering
\caption{Social-Action Breakdown by Type (Authoritative .db Counts, All Conditions B+C, Both Platforms).}
\label{tab:action_breakdown}
\small
\begin{tabular}{lcccc}
\toprule
\textbf{Action} & \textbf{Llama-8B} & \textbf{Mistral-7B} & \textbf{Qwen-14B} & \textbf{Qwen-7B} \\
\midrule
post             & 8,079  & 4,476  & 1,215 & 592 \\
comment          & 5,421  & 1,978  & 585   & 108 \\
like             & 4,567  & 1,946  & 461   & 16  \\
comment\_like    & 3,010  & 167    & 124   & 1   \\
follow           & 5,993  & 708    & 11    & 1   \\
dislike          & 7      & 67     & 0     & 5   \\
comment\_dislike & 10     & 22     & 11    & 0   \\
\midrule
\textbf{Total}   & \textbf{27,087} & \textbf{9,364} & \textbf{2,407} & \textbf{723} \\
\bottomrule
\end{tabular}
\setlength{\tabcolsep}{6pt}
\end{table}

The data reveals a strong model-dependent social-action density. Llama-3.1-8B-AWQ is by far the most socially interactive model (27,087 actions, 40.31/DB), with the highest \texttt{follow} (5,993) and \texttt{comment\_like} (3,010)---indicating it builds the densest peer network and deepest peer-to-peer engagement. Mistral-7B is second (9,364, 35.47/DB) but exhibits a different interaction profile: high \texttt{post} (4,476) and \texttt{like} (1,946) but disproportionately high \texttt{dislike} (67 vs.\ Llama's 7) and low \texttt{comment\_like} (167 vs.\ Llama's 3,010), suggesting antagonistic/\allowbreak broadcast-style rather than deliberative engagement. Qwen2.5-14B---whose full 30-event campaign (180 DBs) was scanned directly---produces moderate, analysis-focused interaction (2,407, 13.38/DB): substantial \texttt{post} (1,215) and \texttt{comment} (585) but minimal \texttt{follow} (11), meaning agents write substantive contributions but build almost no social network. Qwen2.5-7B is the only near-dormant model (723, 2.74/DB). Persona generation quality is high across all models (4-axis DNA, $\sim$40\% institutional accounts, English-only rule enforced). See Appendix~\ref{sec:interaction_realism} for mechanism implications.

\section{The Seven-Dimension Grading Framework}
\label{sec:grading}

To analyze both the ultimate forecasting capability of the swarm and the internal dynamics of emergent intelligence, ECN-BENCH implements a mathematically rigorous 7-dimension scoring system. 

\subsection{D1: Prediction Accuracy (Weight: 0.25)}
The primary forecasting metric is the \textbf{Multi-Category Brier Score ($BS$)}, which evaluates the calibration of the swarm's forecast. For a given event with $K$ mutually exclusive categories, let $f_k$ be the swarm's predicted probability for category $k$, and $o_k$ be the one-hot indicator of the realized outcome ($o_k = 1$ if category $k$ occurred, $0$ otherwise). The Brier Score is formulated as:
\begin{equation}
    BS = \sum_{k=1}^K (f_k - o_k)^2
\end{equation}
This is the unnormalised multi-category (Brier) score, bounded in $[0, 2]$, where $0$ represents a perfect deterministic forecast and $2$ represents a confident forecast on the wrong category. We report this convention---rather than the $\frac{1}{2}$-scaled variant bounded in $[0,1]$---because it is what the pipeline computes: $34$ of the $360$ unit scores in the primary pass exceed $1.0$, which the $\frac{1}{2}$-scaled convention cannot produce. Across all six evaluator passes of all four campaigns ($2{,}070$ scores) the largest is $1.883$ (Mistral-7B, \texttt{C12\_C\_r1}, Evaluator~A read~3); the bound of $2.0$ is attained exactly, six times, only in the truncation grid of \S\ref{sec:truncation_grid}, where a unit places its entire mass on one wrong category. No score anywhere in the released data exceeds $2.0$. Readers comparing against $[0,1]$-normalised Brier scores elsewhere in the literature should halve all values reported here. Note that the uniform-forecast reference used below, $BS_{random} = 1 - 1/K$, is the correct expected score for $f_k = 1/K$ under \emph{this} unnormalised convention, so the Brier Skill Scores we report are internally consistent.

The aggregate Brier Score across $E$ events is $BS_c = \frac{1}{E} \sum_{e=1}^E BS_{e,c}$ for condition $c \in \{A, B, C\}$. 

We formally define \textbf{Simulation Lift ($Lift$)} as the empirical delta in Brier calibration:
\begin{equation}
    Lift = BS_A - BS_B
\end{equation}
A positive $Lift$ indicates that the social simulation workflow produced a better-calibrated forecast than direct elicitation on the events measured. It is a descriptive statistic, not a proof: establishing that the difference is not attributable to sampling variation requires the inference machinery of Section~\ref{sec:statistics}, and interpreting it as a property of the \emph{model} requires that Condition A be run per model, which was not the case in the campaigns reported here (Section~\ref{sec:condition_a}).

To normalize performance across heterogeneous question formats with varying option counts $K \in \{2, 3, 4, 5, 6, 9\}$ and prevent probability dispersion bias, ECN-BENCH mandates the \textbf{Brier Skill Score (BSS)} as the primary skill evaluation metric:
\begin{equation}
    BSS = 1 - \frac{BS_{model}}{BS_{random}} \quad \text{where} \quad BS_{random} = 1 - \frac{1}{K}
\end{equation}
Where $BS_{random}$ represents the expected multi-category Brier error under uniform random guessing ($f_k = \frac{1}{K}$). The BSS metric possesses intuitive econometric bounds:
\begin{itemize}
    \item $BSS = +1.0$: Perfect deterministic prediction ($f_{correct} = 1.0$).
    \item $BSS = 0.0$: Performance equivalent to uniform random guessing.
    \item $BSS < 0.0$: Sub-random performance (worse than uninformed guessing).
\end{itemize}
The score $D1$ is directly mapped from the Brier Skill Score: $D1_{score} = \max(0, 100 \times BSS_B)$. BSS enables fair aggregation of Boolean ($K=2$) and Multiple Choice (MCQ, $K \ge 3$) events into a unified predictability index.

\subsection{D2: Polarization (Weight: 0.15)}
To measure whether the agent swarm fragments into deeply entrenched, opposing camps, we utilize the \textbf{Esteban-Ray (ER) Index of Polarization}. Rather than measuring general dispersion, the ER Index captures both identification (internal group cohesiveness) and alienation (distance between groups). Assume the simulated population forms $G$ distinct opinion clusters regarding the event outcome. Let $p_i$ be the population proportion of cluster $i$, and $y_i$ be the mean probability vector of group $i$. The ER Index is defined as:
\begin{equation}
    ER(\alpha) = \kappa \sum_{i=1}^G \sum_{j=1}^G p_i^{1+\alpha} p_j \|y_i - y_j\|_1
\end{equation}
Where:
\begin{itemize}
    \item $\alpha$ is the polarization sensitivity parameter. We fix $\alpha = 1.0$ to weight internal identification and external alienation equally.
    \item $\kappa$ is a normalization constant scaling $ER \in [0, 1]$.
    \dots
\end{itemize}
A higher $ER$ score represents severe fragmentation. The normalized dimension score is $S_2 = 100 \times (1 - ER_{norm})$.

\subsection{D3: Herd Effect (Weight: 0.15)}
The Herd Effect measures excess conformity—the degree to which individual agents abandon their private signals to adopt the majority consensus purely due to social peer pressure, decoupled from real-world facts. For each step transition $t-1 \rightarrow t$, let $M_t = \text{argmax}_k \bar{p}_{t,k}$ be the swarm's majority opinion class. We define the individual agent conformity indicator as:
\begin{equation}
    Conform_{i,t} = \mathbf{1}\left[\text{argmax}_k p_{i,t,k} = M_t \quad \text{and} \quad \text{argmax}_k p_{i,t-1,k} \neq M_{t-1}\right]
\end{equation}
The raw herd index for a simulation is:
\begin{equation}
    Herd = \frac{1}{T \cdot N_{agents}} \sum_{t=1}^T \sum_{i=1}^{N_{agents}} Conform_{i,t}
\end{equation}
To isolate conformity driven by social mimicry rather than rational updating, ECN-BENCH calculates the excess herd score under relevant information relative to control:
\begin{equation}
    D3_{raw} = Herd_B - Herd_C
\end{equation}
Lower is better. A highly positive $D3_{raw}$ indicates that relevant injections trigger massive uncritical herd conformity, whereas a rational swarm updates beliefs diversely.

\subsection{D4: Deliberation Quality (Weight: 0.15)}
To evaluate whether agent interactions are reasoned or reactive, ECN-BENCH adapts the \textbf{Discourse Quality Index (DQI)} from political science. The fixed evaluator (Section~\ref{sec:confound}) scores a stratified random sample of $M=100$ agent interactions. Each interaction is graded across four sub-criteria (each scored $0$ to $3$ or $0$ to $4$):
\begin{enumerate}
    \item \textbf{Participation:} Equality and accessibility of speech acts.
    \item \textbf{Level of Justification:} Whether claims are accompanied by logical assertions ($0$: no justification; $1$: inferior; $2$: qualified; $3$: sophisticated).
    \item \textbf{Respect:} Tone toward peers ($0$: disrespectful; $1$: neutral; $2$: explicitly respectful).
    \item \textbf{Constructive Politics:} Collaboration toward a consensus ($0$: positional; $1$: alternative proposals; $2$: mediating consensus).
\end{enumerate}
The normalized score is $S_4 = 100 \times \frac{Score_{observed} - Score_{min}}{Score_{max} - Score_{min}}$.

\subsection{D5: Susceptibility (Weight: 0.10)}
Susceptibility quantifies whether the swarm updates its beliefs in response to meaningful information (Claim 3). We track the belief updating efficacy under meaningful perturbation (Condition B) relative to structurally identical null noise (Condition C). We define the post-injection calibration improvement as:
\begin{equation}
    PostImprove_{e,m,c} = BS_{e,m,c, (t=29)} - BS_{e,m,c, (t=60)}
\end{equation}
Where $t=29$ is the step immediately preceding the perturbation. The raw Susceptibility score is the difference in lift:
\begin{equation}
    D5_{raw} = \frac{1}{E} \sum_{e=1}^E \left(PostImprove_{e,m,B} - PostImprove_{e,m,C}\right)
\end{equation}
A positive $D5_{raw}$ mathematically demonstrates that the swarm updates its beliefs productively when presented with real signal, while remaining resistant to null noise.

\subsection{D6: Convergence (Weight: 0.10)}
We model consensus formation through the \textbf{DeGroot Consensus Model}. In the classical model, agent beliefs evolve via a transition probability matrix $W$ representing social trust, such that $x(t+1) = W x(t)$. In ECN-BENCH, we measure the mathematical distance to consensus by the final round. First, we define the dispersion trajectory of agent probability vectors at step $t$:
\begin{equation}
    Disp_t = \frac{1}{N_{agents}(N_{agents}-1)} \sum_{i < j} \sum_{k=1}^K |p_{i,t,k} - p_{j,t,k}|
\end{equation}
The opinion convergence (contraction rate) is defined as:
\begin{equation}
    Contr = 1 - \frac{Disp_{T}}{Disp_{1}}
\end{equation}
To prevent the swarm from converging on the wrong answer (i.e., highly confident but incorrect consensus), we implement an \textbf{Accuracy Guard}:
\begin{equation}
    Guard = 1 - BS_B
\end{equation}
The convergence score is formulated as the joint product:
\begin{equation}
    D6_{raw} = \text{mean}_e (Contr_e \times Guard_e)
\end{equation}
A high score requires both high consensus formation and high accuracy.

\subsection{D7: Information Diversity (Weight: 0.10)}
If all agent personas act identically, the simulation holds no architectural value. We quantify semantic diversity using \textbf{Shannon Entropy}. The fixed evaluator maps the arguments extracted from the Graphiti temporal memories into $S$ discrete semantic bins. Let $q_i$ be the relative frequency of argument bin $i$. The entropy is:
\begin{equation}
    H = - \sum_{i=1}^S q_i \ln(q_i)
\end{equation}
To normalize this across varying bin sizes, we compute Pielou's Evenness (Shannon Equitability):
\begin{equation}
    E_H = \frac{H}{\ln(S)}
\end{equation}
The final score is $S_7 = 100 \times \text{mean}_{e} \text{mean}_{t \in \{1..T-1\}} E_{H, e, t}$.

\subsection{Composite Aggregation and Grade Bands}
The final \textbf{ECN-Predictability Index (EPI)} is a weighted composite of the seven normalized dimension scores ($S_1 \dots S_7 \in [0, 100]$):
\begin{equation}
    EPI = 0.25 S_1 + 0.15 S_2 + 0.15 S_3 + 0.15 S_4 + 0.10 S_5 + 0.10 S_6 + 0.10 S_7
\end{equation}
The resulting $EPI$ maps directly to standardized academic grade bands, as defined in Table~\ref{tab:grades}.

\begin{table}[!htbp]
\centering
\caption{ECN-BENCH Score Grading Bands}
\begin{tabular}{lcl}
\toprule
\textbf{EPI Range} & \textbf{Grade} & \textbf{Interpretation} \\
\midrule
90 - 100 & A+ & Exceptional multi-agent forecasting and dynamic deliberation \\
80 - 89  & A  & Strong, well-calibrated, and resilient swarm intelligence \\
70 - 79  & B  & Solid performance, with minor mcalibrations or conformity issues \\
60 - 69  & C  & Mixed/unstable quality; susceptible to noise \\
50 - 59  & D  & Weak simulation lift; poor deliberation quality \\
$<$ 50   & F  & Fails baseline forecasting expectations \\
\bottomrule
\end{tabular}
\label{tab:grades}
\end{table}

\section{Statistical Inference and Econometric Validation}
\label{sec:statistics}

Given the high compute costs, sample sizes are severely constrained: the campaigns reported here use $E=30$ events per model, and the specification anticipates as few as $E=10$ for expensive API runs. Section~\ref{sec:power_limitation} shows that $E=30$ delivers under $9\%$ power against the effect sizes actually observed, so the tests below should be read as descriptive rather than confirmatory. Reporting simple empirical means is statistically invalid due to the non-normal distribution of squared calibration errors. ECN-BENCH mandates two strict statistical protocols:

\subsection{The Diebold-Mariano Test for Predictive Superiority}
To determine if the With-Sim Brier scores (Condition B) are significantly better than the No-Sim scores (Condition A), we utilize the \textbf{Diebold-Mariano (DM) Test}. The DM test compares the forecast accuracy of two competing series, adjusting for potential autocorrelation. 

Let the Brier loss for event $e$ under Condition A be $L_A(e)$ and under Condition B be $L_B(e)$. We define the loss differential as:
\begin{equation}
    d_e = L_A(e) - L_B(e)
\end{equation}
The null hypothesis $H_0$ is that the expected loss differential is zero: $\mathbb{E}[d_e] = 0$. The DM statistic is formulated as:
\begin{equation}
    DM = \frac{\bar{d}}{\sqrt{\hat{\sigma}_d^2 / N}}
\end{equation}
Where $\bar{d}$ is the sample mean of the differentials, and $\hat{\sigma}_d^2$ is a consistent estimate of the asymptotic variance of $\sqrt{N}\bar{d}$. Assuming independent event-level observations across Polymarket dates, the autocovariance at lags $>0$ is zero, and $\hat{\sigma}_d^2$ reduces to the sample variance of $d_e$. Under $H_0$, the DM statistic asymptotically follows a standard normal distribution. A positive DM statistic with $p < 0.05$ (critical value $> 1.64$ for a one-sided test of superiority) rigorously validates the Simulation Lift Claim.

\subsection{Bootstrap Confidence Intervals}
To compute robust bounds for the Brier score of each condition, researchers must use non-parametric bootstrapping:
\begin{enumerate}
    \item Draw $B = 1,000$ bootstrap replicates of size $N$ from the observed event loss set $\{BS_e\}_{e=1}^N$ with replacement.
    \item Compute the mean Brier score $\overline{BS}^*_b$ for each replicate $b \in \{1 \dots B\}$.
    \item Determine the $95\%$ Confidence Interval ($CI_{95}$) using the percentile method, extracting the $2.5\text{th}$ and $97.5\text{th}$ percentiles of the bootstrap distribution.
\end{enumerate}

\section{Anti-Confound Protocols and Weight Integrity}
\label{sec:confound}

To ensure that observed variations in simulation performance are strictly attributable to the benchmarked models rather than systemic variables, ECN-BENCH mandates two architectural controls:

\subsection{Methodological Constants}
\begin{itemize}
    \item \textbf{Fixed Graph-Building LLM:} Zep/Graphiti uses an internal LLM to extract entities and relations from agent transcripts. ECN-BENCH mandates that this internal model remains absolutely fixed (e.g., Qwen2.5-72B-Instruct running locally) across all benchmarks. If each benchmarked model is used to build its own memory graph, differences in memory structure become a severe confounding variable.
    \item \textbf{Fixed Evaluator LLM:} a single model is hard-coded as the MCQ ReportAgent interviewer and DQI evaluator across all runs, to prevent circular self-evaluation bias. For the July 2026 campaigns reported here this was \textbf{DeepSeek-R1-14B}, served locally, because the HPC compute nodes are air-gapped and no external API call is possible from within a batch job. The benchmark specification names GPT-4o-mini as the reference evaluator and earlier drafts of this paper described it as such throughout; that description was inherited from the specification and does not match what was executed. All evaluator-reliability results in Section~\ref{sec:evaluator_reliability} pertain to DeepSeek-R1-14B, and the low $\kappa$ values reported there should be read as a property of that evaluator at this rubric granularity, not of GPT-4o-mini.
\end{itemize}

\subsection{The Gold Standard Context Protocol}
To completely eradicate lookahead bias and ``web-leakage'' (where models bypass local prompts and retrieve resolved event dates from pre-training data), ECN-BENCH context files (\texttt{data/seeds/*/context.md}) undergo a complete sanitization:
\begin{enumerate}
    \item \textbf{Zero-URL Policy:} All URLs are completely stripped.
    \item \textbf{Strict Resolution Contracts:} Resolution criteria are written as strict legal contracts defining exact institutional sources of truth and resolving fractional edge cases.
    \item \textbf{Cutoff Enforcement:} All data is strictly frozen on the eve of the event, replacing absolute predictive language with probabilistic models to force active reasoning.
\end{enumerate}

The multi-criteria parameter sweep that selected the $N=300$, $R=60$ production configuration used for every result below, and the sub-15B model exclusion criterion derived alongside it, are engineering calibration steps rather than part of the central finding; the full sweep is reported in Appendix~\ref{sec:results}.

\section{The July 2026 Campaign at Face Value}
\label{sec:full_campaign}

\begin{figure*}[!htbp]
\centering
\includegraphics[width=\textwidth]{fig8_architecture.png}
\caption{ECN-BENCH Full System Architecture. The pipeline operates in three phases: (1) Phase 1 persona generation via LLM JSON extraction producing $N=300$ diverse agent DNA profiles; (2) concurrent Twitter and Reddit social simulation for $R=60$ rounds with Graphiti temporal memory; (3) belief aggregation and 7-dimension LLM evaluation. Condition B injects real market intelligence at round $R=30$; Condition C injects structurally identical null noise. Condition A performs direct elicitation without any simulation.}
\label{fig:architecture}
\end{figure*}

\preauditnote{This section reports the campaign \emph{as the benchmark originally scored it}, before any of the audit that follows. It is presented at face value deliberately: it is the artifact Part~I examines, and the reader needs to see what the instrument reported in order to see what the audit changes. Every quantity in Section~\ref{sec:full_campaign} is produced by the single July 2026 evaluator pass whose defect Section~\ref{sec:epistemic_mapping} documents and whose instability Section~\ref{sec:replication_noise} measures; several are additionally taken over Llama-3.1-8B transcripts that the outage of Appendix~\ref{sec:outage} later invalidated and that have since been re-simulated. None of it is a validated benchmark result. Section~\ref{sec:replication_implications} states precisely which of these numbers we withdraw, and Appendix~\ref{sec:errata_appendix} lists every correction by line. Readers looking for what the benchmark actually establishes should go to Section~\ref{sec:predictive_signal}.}

The configuration sweep that fixed $N=300$, $R=60$ using the Llama-3.1-8B reference model on two anchor events (C1, C9) is reported in Appendix~\ref{sec:results} (Appendices~\ref{sec:n_scaling}--\ref{sec:optimal_bounds}). At that configuration we ran three full 30-event campaigns covering the entire ECN-BENCH catalog for each of three model families, on NVIDIA A100 GPUs during July 2026. Figure~\ref{fig:architecture} illustrates the complete pipeline.

\subsection{Per-Event Brier Score Results}
\label{sec:full_per_event}

The per-event Brier scores for all 30 events across the three conditions and three models are lookup tables and are reported in Appendix~\ref{sec:per_event_tables} (Tables~\ref{tab:per_event_llama_8b}--\ref{tab:per_event_qwen_7b}). Here we show only the two views that carry the argument: Figure~\ref{fig:lift_heatmap} renders the per-event lift as a heatmap, and Figure~\ref{fig:cross_model_brier} summarises the three conditions per model with bootstrap intervals. A positive lift entry indicates the social simulation (Condition B) outperformed direct elicitation (Condition A), which the benchmark originally read as evidence for the Prediction Lift hypothesis (Claim 1).


\begin{figure*}[!htbp]
\centering
\includegraphics[width=\textwidth]{fig1_lift_heatmap.png}
\caption{Per-Event Simulation Lift Heatmap. Lift = $BS_A - BS_B$; green cells indicate With-Sim outperforms the shared No-Sim reference. The bimodal structure is evident: a cluster of events (T2, C11, S9, S3) consistently benefits from simulation (green band), while a second cluster (T1, C10, C15, S6) consistently degrades (orange band). X-axis labels are color-coded by taxonomy: violet = Social/Electoral (C), cyan = Senate/Binary (S), red-orange = Technical/Economic (T). Note that since Condition A is a single shared record (Section~\ref{sec:condition_a}), row-to-row variation within a column reflects $BS_B$ only.}
\label{fig:lift_heatmap}
\end{figure*}

\begin{figure*}[!htbp]
\centering
\includegraphics[width=0.9\textwidth]{fig2_cross_model_brier.png}
\caption{Cross-Model Mean Brier Score by Condition (Cond A = No-Sim, B = With-Sim, C = Null). Error bars show 95\% bootstrap confidence intervals ($B=10{,}000$). Gold bi-directional arrows mark the B-to-C susceptibility margin: the gap is positive for all three models, though significant only for Mistral-7B (Figure~\ref{fig:susceptibility}). All BS$_B$ values exceed BS$_A$; within the SBE subset alone the mean lift is $+0.363$, but see Appendix~\ref{sec:tautology} for why that subset cannot be read as evidence of a difficulty--benefit relationship.}
\label{fig:cross_model_brier}
\end{figure*}

\subsection{Condition A Is a Shared Reference, Not a Per-Model Baseline}
\label{sec:condition_a}

Before presenting the aggregate statistics we must state a structural property of the campaign that constrains every ``lift'' figure in this paper. Condition A was executed once and its output reused verbatim for all three model campaigns. The $BS_A$ values are not merely equal in the mean: the full predicted probability vectors are byte-identical across Llama-3.1-8B, Mistral-7B and Qwen2.5-7B for all 30 events. Determinism (temperature $0$, seed $42$) cannot produce this, since three different model families do not emit identical distributions; the Condition A record was generated once and copied.

The consequence is that $Lift = BS_A - BS_B$ is \emph{not} a within-model contrast between a model's own direct elicitation and its own simulated swarm. It is the distance from each model's simulated forecast to a single fixed external reference. Two things follow. First, all cross-model differences in reported lift are differences in $BS_B$ alone and carry no information about the models' direct-elicitation ability. Second, Claim 1 as originally posed---``does the simulation step add lift beyond what \emph{the model} can do alone''---is not answered by this campaign for any model; answering it requires re-running Condition A per model. We treat this as the principal design defect of the July 2026 campaign and flag every affected quantity accordingly.

\begin{figure}[!htbp]
\centering
\includegraphics[width=0.95\columnwidth]{figB_susceptibility.png}
\caption{Susceptibility margin $BS_C - BS_B$ per model over the full 30-event set, with $95\%$ percentile bootstrap confidence intervals ($B=20{,}000$). Values above zero indicate the swarm integrates the relevant round-30 injection better than the length-matched null injection. The margin is positive for all three architectures; only Mistral-7B's interval excludes zero.}
\label{fig:susceptibility}
\end{figure}

\subsection{Cross-Model Statistical Inference}
\label{sec:cross_model_stats}

\confoundnote{No two campaigns in this table share a serving stack and a
quantization level, and no two share a model generation either: Llama-3.1-8B is
AWQ-INT4 under vLLM (Jul 2024), Mistral-7B is \emph{unquantised FP16} under vLLM
(Sep 2023), and Qwen2.5-7B is GGUF~Q4\_K\_M under Ollama (Sep 2024). Model
generation, numerical precision and inference server therefore vary together
across every row, and none is held constant (Section~\ref{sec:serving_note},
Table~\ref{tab:provenance}). Differences \emph{between} rows are consequently not
attributable to architecture and are reported as descriptive only. Every quantity
here also comes from the single Evaluator~J pass whose schema-enforcement defect
is documented in Section~\ref{sec:epistemic_mapping}; switching evaluator reverses
the sign of one model's lift and flips which models' susceptibility intervals
exclude zero (Section~\ref{sec:replication_implications}). Within-row,
within-model contrasts---principally the Condition~B versus Condition~C
susceptibility margin---are the only comparisons this design supports, and even
those are subject to the shared-reference defect of
Section~\ref{sec:condition_a}.}

Table~\ref{tab:cross_model_stats} presents the aggregate statistical summary for all three models, computed over the full $E=30$ event set. Bootstrap confidence intervals use $B=10{,}000$ replicates. The Diebold-Mariano test is applied one-sided ($H_1$: $\mathbb{E}[d] > 0$, i.e., With-Sim lower loss than No-Sim).

\begin{table*}[!htbp]
\centering
\caption{Cross-Model Statistical Summary: Full 30-Event Campaign ($N=300$, $R=60$, July 2026). Bootstrap CI ($B=10{,}000$ replicates, 95\% percentile method). DM Test: one-sided. Achieved power computed for two-sided $\alpha=0.05$ at the measured effect size and $E=30$ events. Column abbreviations: $\Delta_{C-B} =$ BS$_C-$BS$_B$ (the susceptibility margin); Lift $=$ Lift$_{A \to B}$; DM $=$ Diebold--Mariano statistic; $p =$ its one-sided $p$-value; Pwr.\ $=$ achieved power.}
\label{tab:cross_model_stats}
\scriptsize
\setlength{\tabcolsep}{2pt}
\resizebox{\textwidth}{!}{%
\begin{tabular}{lccccccccccc}
\toprule
\textbf{Model} & \textbf{Serving} & \textbf{BS$_A$} & \textbf{BS$_B$ [95\% CI]} & \textbf{BS$_C$} & \textbf{$\Delta_{C-B}$ [95\% CI]} & \textbf{Lift} & \textbf{DM} & \textbf{$p$} & \textbf{$d_z$} & \textbf{Pwr.} & \textbf{Dir.Acc$_B$} \\
\midrule
\textbf{Llama-3.1-8B}  & vLLM/OpenRouter\textsuperscript{$\dagger$} & 0.3846 & 0.6850 [0.538, 0.834] & 0.5633 & $-0.122$ [$-0.232$, $-0.022$] & $-0.300$ & $-3.36$ & 1.00 & $-0.613$ & 91.9\%\textsuperscript{$\ddagger$} & 0.467 \\
\textbf{Mistral-7B}    & vLLM, \textbf{FP16} & 0.3846 & 0.5312 [0.429, 0.626] & 0.6462 & $\mathbf{+0.115}$ [$+0.009$, $+0.225$] & $-0.147$ & $-1.65$ & 0.95 & $-0.301$ & 8.8\% & 0.567 \\
\textbf{Qwen2.5-7B}    & Ollama, GGUF Q4 & 0.3846 & 0.4814 [0.372, 0.590] & 0.5345 & $+0.053$ [$-0.064$, $+0.172$] & $-0.097$ & $-1.18$ & 0.88 & $-0.215$ & 8.5\% & \textbf{0.733} \\
\bottomrule
\end{tabular}%
}
\smallskip\\
\footnotesize $BS_A$ is again identical across all three rows: Condition A is a single No-Sim record shared byte-for-byte across the campaigns, and the corrected 2026-08-16 repair (\S\ref{sec:outage}) restored it for the 24 events whose Condition A the first, defective repair had overwritten. \textsuperscript{$\dagger$}Llama's 24 repaired events were re-simulated via OpenRouter (\texttt{meta-llama/llama-3.1-8b-instruct}), not the local AWQ-INT4 vLLM serving used for its other 6 events and for both other campaigns throughout; this is a within-row serving-stack change, on top of the pre-existing cross-row confound of \S\ref{sec:serving_note}. \emph{Instrument caveat.} Those 24 events did not exist in July, so no scoring by the July local DeepSeek-R1-14B evaluator is possible for them; this row's $BS_B$ and $BS_C$ come from an August OpenRouter \texttt{deepseek-r1} pass while its $BS_A$ and the Mistral/Qwen rows remain July-scored. Llama's lift and susceptibility are therefore reported on an instrument-consistent basis (all three conditions read by one evaluator) in Table~\ref{tab:cutoff_lift} and \S\ref{sec:cutoff}, and the spread across the three August evaluators is given there; that spread is wide, and no single number should be read without it. The negative Lift$_{A \to B}$ indicates the With-Sim mean Brier score \emph{exceeded} that reference on aggregate, for Llama and Mistral; Qwen2.5-7B's CI still spans zero. The susceptibility margin $BS_C - BS_B$ is positive (favouring the hypothesis) for Mistral and Qwen2.5-7B; for Llama it is $-0.122$ under this evaluator but $+0.093$ [$-0.069$, $+0.271$] under the instrument-consistent pairing and $+0.027$ under the third, so its sign is an evaluator effect rather than a property of the data (\S\ref{sec:cutoff}). All DM $p$-values are non-significant ($p > 0.10$). We report $d_z$ (the paired standardised mean difference, $\bar{d}/s_d$) rather than the pooled two-sample $d$: Conditions A and B are evaluated on the same 30 events, so the design is within-subjects. \textsuperscript{$\ddagger$}Llama's achieved-power figure was recomputed 2026-08-16 from a normal approximation to the paired $t$ ($\text{power} \approx \Phi(|\text{DM}| - z_{0.975})$) rather than the pipeline's original (unlocated) method, so it is not exactly comparable to Mistral's and Qwen's $8.5$--$8.8\%$, which were not recomputed. It is high here only because Llama's observed effect is large; observed power restates the observed $p$-value and says nothing about the design, which remains critically underpowered---$E \geq 92$ would still be required for 80\% power at $d=0.5$ (\S\ref{sec:power_limitation}).
\end{table*}

Three key patterns emerge from the full-campaign data:

\textbf{Finding 1 --- The Susceptibility Ordering Holds for Two Models and Is Evaluator-Dependent for the Third:} The With-Sim condition (B) yields a lower mean Brier score than the Null Control condition (C) for Mistral-7B ($0.531$ vs.\ $0.646$, margin $+0.115$) and Qwen2.5-7B ($0.481$ vs.\ $0.535$, margin $+0.053$), the qualitative prediction of the Information Susceptibility hypothesis (Claim 3). For Llama-3.1-8B the sign depends on which evaluator reads the transcripts: $-0.122$ [$-0.232$, $-0.022$] under the \texttt{deepseek-r1} pass tabulated above, $+0.093$ [$-0.069$, $+0.271$] under \texttt{deepseek-v4-flash-0731}, and $+0.027$ under \texttt{gpt-5.6-luna-pro}---two of the three centred on the hypothesis' side, one excluding zero against it, on \emph{identical} simulation data. Because $BS_C - BS_B$ is a contrast between two simulated conditions and never touches Condition A, this spread is a property of the measuring instrument, not of the campaign; we report it as such rather than selecting a sign. Figure~\ref{fig:susceptibility} shows the margins with bootstrap confidence intervals. Only Mistral-7B's interval excludes zero on the hypothesis' side ($[+0.009, +0.225]$); Qwen's spans zero. At $E=30$ events the design cannot separate these effects from no effect, and \S\ref{sec:power_limitation} quantifies how little power is available.

We note a correction to earlier drafts of this analysis, which reported a sign reversal for Qwen2.5-7B ($-0.044$) and attributed it to that model's near-dormant social channel. That figure was the \texttt{content\_susceptibility.delta.B\_minus\_C} field of the pipeline summary, which measures the mean \emph{yes-probability} difference between conditions, not the Brier margin; the two are distinct metrics and were inadvertently mixed across rows of the same table. On the Brier metric used for the other two models, Qwen2.5-7B's margin is $+0.053$ and in the same direction as the rest. The dormant-channel observation (723 social actions, 2.74/DB) remains factually correct and is discussed in Appendix~\ref{sec:interaction_realism}, but it is no longer needed to explain a reversal that does not exist.

\textbf{Finding 2 --- Bimodal Per-Event Distribution (Mistral/Qwen; Not Llama):} A subset of high-information-gain events (T2, C11, S9, C8, C13, C5, S3) shows large positive lift consistently for Mistral-7B and Qwen2.5-7B, neither of which was affected by the outage or its repair (e.g.\ T2 lift $+1.26$ and $+1.275$). It does not hold for Llama-3.1-8B, whose per-event lifts were entirely re-derived from the corrected re-simulation of \S\ref{sec:outage}: twenty-four of the thirty events, including T2, C11, C8, C13, C5 and C15, now rest on transcripts produced on a different serving stack and read by a different evaluator than the other two campaigns. Per-event agreement across models is therefore no longer an apples-to-apples comparison for this campaign, and we report the cluster as a Mistral/Qwen pattern rather than a cross-model one. The only Llama events that remain directly comparable are the six the outage never touched (C3, S2--S6), which is too few to establish or refute a bimodal structure.

\textbf{Finding 3 --- Directional Accuracy Divergence (Qwen Best):} While Brier scores aggregate probabilistic calibration, directional accuracy (whether the model correctly identifies the winning outcome) reveals architectural differences. Qwen-7B achieves the highest directional accuracy in Condition B (73.3\%), matching the No-Sim baseline identically ($\Delta = 0.0$), indicating the simulation process does not degrade winner identification capacity. Mistral-7B shows significant degradation ($73.3\% \to 56.7\%$, $\Delta = -16.7\%$), confirming that over-deliberation at $R=60$ actively corrupts its directional forecasts. Llama-3.1-8B's Condition A baseline is again the shared $73.3\%$ after the corrected repair restored it; against that reference its With-Sim directional accuracy is $46.7\%$ ($\Delta = -26.7\%$) under the evaluator that scored its re-simulated events, the largest degradation of the three---but see the instrument caveat in Table~\ref{tab:cross_model_stats}'s note, since that evaluator also scores the other conditions of this campaign more harshly in absolute terms.

\subsection{Event Taxonomy Breakdown}
\label{sec:taxonomy_breakdown}

\confoundnote{This subsection compares models within taxonomy categories and is
subject to the same three-way confound as
Section~\ref{sec:cross_model_stats}: release generation, precision and serving
stack vary together across the three campaigns. Statements below of the form
``model $X$ excels on category $Y$'' therefore describe what this campaign
recorded; they do not isolate architecture, and with $n = 7$--$15$ events per
category they are additionally underpowered. The category counts are also small
enough that a single flat unit---of which Llama-3.1-8B has $8$ and Mistral-7B $7$
in Condition~B under this evaluator (Section~\ref{sec:degenerate_flat})---moves a
category mean materially.}

Table~\ref{tab:taxonomy_brier} reports mean Condition B Brier scores by taxonomy category and Figure~\ref{fig:taxonomy_breakdown} plots them alongside directional accuracy. The partition is the three-way prefix scheme of Section~\ref{sec:taxonomy}, and---as noted there---``Senate/Binary'' is a question-format label rather than a topical one, so the category contrasts below should be read as descriptive rather than as a claim about subject matter.

\begin{figure*}[!htbp]
\centering
\includegraphics[width=\textwidth]{fig3_taxonomy_breakdown.png}
\caption{Event Taxonomy Performance Breakdown. Left: Mean Brier Score (Cond B, With-Sim) by taxonomy category per model. Right: Directional accuracy percentage per taxonomy category. Qwen-7B dominates on Senate/Binary events ($BS_B = 0.361$, Dir.Acc $= 86\%$), while Llama-8B shows the most balanced cross-category performance. Senate/Binary events show highest cross-model agreement on directional accuracy, while Tech/Economic events exhibit the widest variance.}
\label{fig:taxonomy_breakdown}
\end{figure*}

To identify whether simulation benefits are event-type-specific, we decompose the mean Brier scores by the ECN-BENCH $2\times2$ taxonomy categories: Social/Electoral (prefix C, $n=15$), Senate/Binary (prefix S, $n=7$), and Technical/Economic (prefix T, $n=8$).

\begin{table}[!htbp]
\centering
\caption{Mean Brier Scores by Taxonomy Category (Condition B --- With-Sim, $N=300$, $R=60$)}
\label{tab:taxonomy_brier}
\begin{tabular}{llccc}
\toprule
\textbf{Category} & \textbf{$n$} & \textbf{Llama-8B} & \textbf{Mistral-7B} & \textbf{Qwen-7B} \\
\midrule
\textbf{Social/Electoral (C)} & 15 & 0.3451\textsuperscript{$\dagger$} & 0.6164 & 0.5082 \\
\textbf{Senate/Binary (S)}    &  7 & 0.5950\textsuperscript{$\dagger$} & 0.4129 & 0.3607 \\
\textbf{Tech/Economic (T)}    &  8 & 0.3619\textsuperscript{$\dagger$} & 0.4750 & 0.5369 \\
\bottomrule
\end{tabular}
\setlength{\tabcolsep}{6pt}
\end{table}

\textsuperscript{$\dagger$}Llama-8B's row was recomputed 2026-08-16 after the outage repair (\S\ref{sec:outage}), which redistributes events across categories differently than the earlier draft's numbers reflected (15 of Llama's 24 repaired events fall in the Social/Electoral category alone). The taxonomy decomposition no longer supports a ``most balanced generalist'' reading for Llama: it is now the \emph{best} of the three on Social/Electoral ($0.345$, against $0.616$ and $0.508$) and the \emph{worst} on Senate/Binary ($0.595$, against $0.413$ and $0.361$)---a wide, specific split rather than balance across categories, and the opposite ranking from what the pre-repair data showed on Social/Electoral. Qwen-7B still leads Senate/Binary ($0.361$) by a smaller margin than before. Mistral-7B remains competitive on Technical/Economic ($0.475$) despite its overall lower directional accuracy. With $n=7$--$15$ events per category (\S\ref{sec:taxonomy_breakdown}'s own caveat), none of these per-category contrasts should be read as more than descriptive.

\subsection{Calibration Quality Analysis}
\label{sec:calibration_analysis}

\begin{figure}[!htbp]
\centering
\includegraphics[width=0.65\textwidth]{fig4_calibration_overlay.png}
\caption{Calibration Curves: Condition B (With-Sim), all three models, $N=300$, $R=60$. All models fall above the perfect calibration diagonal in the middle range $[0.25, 0.75]$, indicating systematic \emph{under-confidence} (actual hit rates exceed predicted probabilities). The high bracket $[0.75, 1.0]$ achieves near-perfect calibration (100\% hit rate) across all models. Llama-8B achieves closest calibration at $[0.25, 0.5]$ bracket.}
\label{fig:calibration_overlay}
\end{figure}

\confoundnote{The calibration curves below are per-model and are subject to the
same three-way provenance confound as
Sections~\ref{sec:cross_model_stats}--\ref{sec:taxonomy_breakdown}; differences
between the three curves do not isolate architecture. Two further cautions apply
specifically here. The bracket occupancies are small ($n = 2$--$15$ per cell), so
several ``$100\%$ hit rate'' entries rest on two or three events. And the
probabilities being bracketed are the Evaluator~J outputs: because that
instrument returns an exactly uniform vector on $20.6\%$ of simulated units
(Section~\ref{sec:degenerate_flat}), probability mass is displaced toward the
middle brackets by the extraction defect itself, which is a competing explanation
for the ``systematic under-confidence'' read below.}

Table~\ref{tab:calibration_quality} reports the probability calibration bracket statistics for Condition B across all three models, and Figure~\ref{fig:calibration_overlay} overlays the corresponding calibration curves. An ideal calibration corresponds to predicted probability $\approx$ empirical rate within each bracket.

\begin{table*}[!htbp]
\centering
\caption{Condition B Calibration Bracket Analysis ($N=300$, $R=60$): predicted vs.\ empirical hit rates.}
\label{tab:calibration_quality}
\small
\begin{tabular}{lcccccc}
\toprule
 & \multicolumn{2}{c}{\textbf{Llama-3.1-8B}} & \multicolumn{2}{c}{\textbf{Mistral-7B}} & \multicolumn{2}{c}{\textbf{Qwen2.5-7B}} \\
\cmidrule(r){2-3} \cmidrule(lr){4-5} \cmidrule(l){6-7}
\textbf{Bracket} & \textbf{Pred.} & \textbf{Actual} & \textbf{Pred.} & \textbf{Actual} & \textbf{Pred.} & \textbf{Actual} \\
\midrule
{[0.00--0.25)} & 0.219 & 1.000 ($n=2$) & 0.200 & 1.000 ($n=2$) & --- & --- \\
{[0.25--0.50)} & 0.353 & 0.833 ($n=12$) & 0.344 & 0.615 ($n=13$) & 0.368 & 0.692 ($n=13$) \\
{[0.50--0.75)} & 0.631 & 0.846 ($n=13$) & 0.599 & 0.667 ($n=12$) & 0.643 & 0.733 ($n=15$) \\
{[0.75--1.00]} & 0.850 & 1.000 ($n=3$) & 0.833 & 1.000 ($n=3$) & 0.850 & 1.000 ($n=2$) \\
\bottomrule
\end{tabular}
\setlength{\tabcolsep}{6pt}
\end{table*}

All three models show \textbf{systematic under-confidence in the middle brackets}: predicted probabilities in the [0.25--0.50) range consistently underestimate outcome frequency (empirical rates 61.5\%--83.3\% vs.\ predicted 34.4\%--36.8\%). This under-confidence compression is a hallmark of LLM ensemble calibration under social deliberation: when agents aggregate beliefs through social interaction, probability mass accumulates in the moderate-certainty range ($0.3$--$0.7$), creating predictions that are too uncertain relative to the empirical hit rate. The high bracket [0.75--1.00] shows near-perfect calibration across all models ($100\%$ hit rate), indicating that events where the swarm achieves near-certainty are genuinely high-confidence outcomes.


\parttitle{I}{The Instrument Failure}

\section{Loss of Sharpness and Degenerate Consensus: The Central Finding}
\label{sec:degenerate}

Part~I opens with the phenomenon itself, characterised exactly as we first observed it: a mechanism by which forecast informativeness was eroded in this campaign, discovered while auditing the aggregate lift reversal reported at face value in Section~\ref{sec:cross_model_stats}. The two effects below---flat, non-informative forecasts and a milder, pervasive loss of sharpness among the remaining forecasts---account for a substantial part of that reversal and require no claim about the swarm's reasoning quality having deteriorated. This section attributes them, provisionally, to the deliberation; Sections~\ref{sec:epistemic_mapping}--\ref{sec:volume_mechanism} then take that attribution apart and relocate it to the instrument, which is the paper's actual contribution. We present the two in that order because it is the order in which the evidence arrived, and because the provisional reading is the one a reader of the aggregate would reach unaided.

\subsection{Flat Forecasts}
\label{sec:degenerate_flat}

We define a unit as \emph{flat} when every option receives probability within $0.005$ of $1/K$, i.e.\ the extracted swarm forecast is indistinguishable from a uniform guess. Across the simulated conditions, $37$ of $180$ units ($20.6\%$) are flat: Llama-3.1-8B returns $8/30$ flat under Condition B and $10/30$ under Condition C; Mistral-7B returns $7/30$ and $12/30$; Qwen2.5-7B returns none. Only $1/30$ Condition A units is flat, for each model. Those are the counts as first observed, on the July campaign files; Llama's twenty-four outage-affected events have since been re-simulated and re-scored (\S\ref{sec:outage}), which moves its counts to $7/30$ and $7/30$ and the pooled total to $33/180$ ($18.3\%$). Mistral's and Qwen's are untouched, and the argument of this section---that flat units are concentrated in the simulated conditions and inflate $BS_B$ and $BS_C$ by construction---holds identically on either version. A flat unit is not a forecast---it is an extraction failure, or a swarm that reached no usable consensus---but the Brier score charges it the full uniform loss $1 - 1/K$, which for a binary event is $0.5$. Because flat units are concentrated in Conditions B and C and nearly absent from Condition A, they push $BS_B$ and $BS_C$ upward relative to $BS_A$ by construction. This alone is enough to generate a negative aggregate lift without any claim about deliberation degrading reasoning.

\textbf{The failure is not universal.} Flat forecasts are entirely absent for Qwen2.5-7B (0/30 under both simulated conditions) and concentrated in Llama-3.1-8B and Mistral-7B. An earlier draft of this section attributed this split to the benchmarked swarm model's own serving backend (vLLM for Llama/Mistral vs.\ Ollama/GGUF for Qwen) and speculated that vLLM's guided decoding~\cite{vllm_guided} was responsible. A subsequent source-code audit (Section~\ref{sec:epistemic_mapping}) shows this attribution is not supported: the \texttt{probabilities} field is produced by the fixed evaluator (DeepSeek-R1-14B, Section~\ref{sec:confound}), not by the swarm model being benchmarked, and that evaluator's structured-output enforcement is disabled for an unrelated reason---its model identifier matches a \texttt{"reasoning model"} keyword filter in the serving client---uniformly across all three campaigns regardless of which swarm model is under test. Since that condition is constant across Llama, Mistral and Qwen, it cannot by itself explain why only two of the three degenerate. We withdraw the vLLM-versus-Ollama account and, in its place, note that the one variable we can confirm differs sharply by swarm model is the volume of \texttt{evidence\_text} the unconstrained evaluator must compress into a single distribution: up to $20{,}329$ characters for Llama-3.1-8B against a fraction of that for the near-dormant Qwen2.5-7B social channel (Appendix~\ref{sec:interaction_realism}). Section~\ref{sec:volume_mechanism} takes this from a conjecture to a measurement. Within Mistral-7B---the one campaign in which evidence length varies without also varying whether deliberation occurred---the flat rate rises monotonically from $5\%$ to $25\%$ to $65\%$ across terciles of \texttt{evidence\_text} length, while the same long transcripts, re-read by four independent Evaluator~A passes and by a second architecture, come back sharp. That was our reading for most of this work: evidence volume as the trigger, the schema-disabled evaluator as the locus. Section~\ref{sec:volume_mechanism} tests both by assignment rather than by observation and supports neither---shortening a transcript does not un-flatten it, and toggling enforcement within one model changes nothing---while confirming that the choice of evaluator dominates both. We state the volume relationship here as the association it is, and withdraw the causal reading in \S\ref{sec:truncation_grid}. The confirmed and unconfirmed parts of this account are treated in full in Sections~\ref{sec:epistemic_mapping} and~\ref{sec:volume_mechanism}; what remains open is what it is about this instrument in particular---its distillation, its quantisation, its local serving configuration---that makes it vulnerable where its relatives are not.

\subsection{Loss of Sharpness}
\label{sec:degenerate_sharpness}

Even among non-flat units, simulation systematically flattens the distribution. Restricting to units that are not exactly uniform---so that the effect cannot be an artifact of the flat units of Section~\ref{sec:degenerate_flat}, which carry a modal probability of $1/K$ by construction---the mean probability assigned to the modal option falls from $0.664$ under Condition A to $0.563$, $0.536$ and $0.538$ under Condition B for Llama-3.1-8B, Mistral-7B and Qwen2.5-7B, and falls further still under the null-injection Condition C ($0.538$, $0.475$, $0.488$). Computed over all units, as the aggregate Brier comparison implicitly does, the corresponding Condition B means are $0.658 \to 0.514$, $0.485$ and $0.538$; the gap between the two rows is the contribution of the flat units, so the two mechanisms compound rather than substitute. Since the Brier score is a strictly proper scoring rule, it rewards sharpness conditional on calibration~\cite{gneiting2007}: a forecast that hedges from $0.85$ to $0.55$ on an outcome that occurs is penalised, even though the direction is unchanged. Directional accuracy, which is insensitive to sharpness, barely moves for Llama-3.1-8B ($0.733 \to 0.717$) and not at all for Qwen2.5-7B ($0.733 \to 0.733$); Mistral-7B's larger drop ($0.733 \to 0.567$) is confounded with that model's older generation and full-precision serving (Section~\ref{sec:serving_note}) and should not be read as a general property of sharpness loss. Where directional accuracy is preserved, the divergence from the Brier trend is the signature of a sharpness loss rather than a reasoning failure: the swarm still knows which outcome is more likely, but states that knowledge less confidently than direct elicitation does.

Figure~\ref{fig:degenerate} shows both effects. Together they mean the headline quantity $BS_A - BS_B$ conflates three distinct things: genuine forecasting skill, extraction reliability, and forecast sharpness---and the second and third, not the first, are what this campaign's aggregate lift mostly measures. We recommend that future campaigns report flat-unit rate as a first-class diagnostic, exclude flat units from calibration analysis (or impute them explicitly), and decompose the Brier score into reliability, resolution and uncertainty components so that sharpness effects are separable from skill effects.

\begin{figure*}[!htbp]
\centering
\includegraphics[width=\textwidth]{figE_degenerate_forecasts.png}
\caption{Two mechanical drivers of the aggregate lift reversal. \textbf{Panel A:} percentage of events answered with an exactly uniform probability vector, by model and condition. Flat forecasts are rare under direct elicitation ($3\%$) and common under simulation (up to $40\%$) for the two vLLM-served models, and the Brier score charges each one the full uniform loss. \textbf{Panel B:} mean probability placed on the modal option. Simulation reduces forecast sharpness for every model, which a proper scoring rule penalises independently of whether the predicted direction is correct.}
\label{fig:degenerate}
\end{figure*}

\subsection{Model Provenance and Serving Heterogeneity: A Three-Way Confound}
\label{sec:serving_note}

Earlier drafts described the campaign as comparing three AWQ-quantised models, with a single documented exception for Qwen2.5-7B. An audit of the four production Slurm scripts against the run manifests shows that this description is incorrect in two further respects, and that the comparison is confounded along three axes simultaneously. Because this materially affects how the cross-model results should be read, we set out the actual configuration in Table~\ref{tab:provenance} and label models by their true weights throughout this revision.

\begin{table*}[!htbp]
\centering
\caption{Actual model provenance and serving stack, recovered from the production Slurm scripts. The ``label'' column gives the identifier written into the run manifests and used in earlier drafts; it is a \texttt{--served-model-name} string and does not describe the weights.}
\label{tab:provenance}
\footnotesize
\setlength{\tabcolsep}{4pt}
\begin{tabular}{llllll}
\toprule
\textbf{Manifest label} & \textbf{Actual weights} & \textbf{Released} & \textbf{Precision} & \textbf{Server} & \textbf{Campaign} \\
\midrule
\texttt{llama-3.1-8b-awq}  & \texttt{Meta-Llama-3.1-8B-Instruct-AWQ-INT4} & Jul 2024 & AWQ INT4 & vLLM   & 30 events, local \\
\texttt{mistral-7b-awq}    & \texttt{Mistral-7B-Instruct-\textbf{v0.1}}     & \textbf{Sep 2023} & \textbf{FP16} & vLLM & 30 events, local \\
\texttt{qwen2.5:7b}        & \texttt{qwen2.5:7b} (Ollama)                  & Sep 2024 & GGUF Q4\_K\_M & Ollama & 30 events, local \\
\texttt{qwen2.5-14b-awq}   & \texttt{qwen2.5:14b} (Ollama)                 & Sep 2024 & GGUF & Ollama & 30 events, HPC only \\
\bottomrule
\end{tabular}
\setlength{\tabcolsep}{6pt}
\end{table*}

Three corrections follow.

\textbf{(1) Mistral was not quantised.} The script named \texttt{run\_bench\_final\_mistral-7b-awq.slurm} carries the header comment \texttt{Quant: AWQ-INT4} and serves the model under the name \texttt{mistral-7b-awq}, but its \texttt{BENCH\_MODEL} path resolves to the unquantised \texttt{mistralai/Mistral-7B-Instruct-v0.1} snapshot and the vLLM invocation passes no \texttt{--quantization} flag. Mistral is therefore the only model in the campaign served at \emph{full precision}. The Llama script, by contrast, genuinely points at an AWQ-INT4 repository.

\textbf{(2) Mistral is a model generation older than the others.} \texttt{Mistral-7B-Instruct-v0.1} was released in September 2023, roughly ten months before Llama-3.1-8B-Instruct and twelve before Qwen2.5-7B-Instruct. This is by far the largest single difference between the models compared here.

\textbf{(3) Neither Qwen model was served through vLLM.} In both Qwen scripts the entire vLLM benchmark-server block is commented out and inference is provided by \texttt{ollama pull qwen2.5:7b} and \texttt{ollama pull qwen2.5:14b} respectively. The AWQ repository paths remain in the scripts as dead \texttt{BENCH\_MODEL} assignments, which is the likely origin of the ``Qwen2.5-14B-AWQ'' label used in earlier drafts. Qwen2.5-14B is GGUF via Ollama, not AWQ via vLLM.

\textbf{Consequence for interpretation.} Model generation, numerical precision and serving stack all vary together across the three campaigns, and none is held constant. This is a direct violation of the methodological-constants protocol of Section~\ref{sec:confound}, which the benchmark specification applies to the graph-building and evaluator models but which was not enforced for the benchmarked models themselves. The practical effect falls hardest on Mistral-7B, which records the worst Brier score ($0.531$), the largest directional-accuracy degradation ($-16.7$ points) and the interaction profile that earlier drafts characterised as ``antagonistic'' and explained by architecture-dependent over-deliberation. A far simpler competing explanation is now available and cannot be excluded by these data: it is a substantially older instruction-tuned model. We accordingly withdraw the architectural interpretation of Mistral's results in Appendices~\ref{sec:cross_model} and~\ref{sec:interaction_realism}, and downgrade all three-way model comparisons in this paper to \emph{descriptive}. Isolating architecture would require re-running the campaign with matched release generation and identical quantisation and serving stack across models---which is a tractable fix, since AWQ builds of all four models were already staged in the HPC cache.

\section{The Source-Code Audit: A Keyword Filter Disables Structured-Output Enforcement}
\label{sec:epistemic_mapping}

The benchmark specification defines a third measurement channel beyond the Brier score and the 7-dimension rubric: a \textbf{3-tier micro-question protocol} (``Epistemic Mapping,'' specification Architecture v3.2). For every event, three fact-grounded micro-questions are extracted from the seed context to probe whether the swarm's macro-prediction is supported by correct reasoning about the event's granular drivers---e.g., a predicted rate cut should co-occur with agent-level recognition of weak payroll data. Each micro-question is scored with a \texttt{dominant\_tag} and a \texttt{short\_rationale} grounded strictly in the simulation's discussion log, and the specification defines an \textbf{Epistemic Alignment Score} as the correlation between micro-question accuracy and macro-prediction calibration.

In the three July 2026 production campaigns that are this paper's primary evidence base, this channel produced no usable data: all $2{,}430$ micro-question instances ($3$ questions $\times$ $270$ units $\times$ $3$ models) return \texttt{dominant\_tag: "N/A"} and \texttt{short\_rationale: "Missing in LLM response"}. Taken alone, this would leave the failure's cause undetermined. It is not alone: development history on this cluster includes 63 earlier smoke-test and pilot runs that also populate this field, spanning April--July 2026, and auditing them resolves the question of \emph{where} the failure lives.

We partitioned all $63$ runs project-wide that carry the \texttt{micro\_epistemic\_mapping} field into two groups by their \texttt{benchmark\_model} string: $41$ runs whose model identifier follows OpenRouter's \texttt{provider/model} convention (e.g.\ \texttt{qwen/qwen-2.5-7b-instruct}, \texttt{google/gemma-4-26b-a4b-it}, \texttt{openai/gpt-oss-120b}, \texttt{anthropic/claude-3-haiku}, \texttt{deepseek/deepseek-v3.2}, \texttt{meta-llama/llama-3.1-8b-instruct:nitro}), consistent with cloud API routing via OpenRouter (15 of the 41 run directories are additionally named \texttt{*openrouter*} explicitly); and $22$ runs whose identifier follows the local \texttt{served-model-name} convention used by the production Slurm scripts (\texttt{llama3.1-8b-awq}, \texttt{qwen14b}, direct HPC \texttt{/scratch/.../snapshots/...} weight paths), consistent with local vLLM/Ollama serving on the HPC cluster. Manifests do not log an explicit API endpoint, so this classification is inferred from naming convention and directory context rather than read directly off a config field; we report it with that caveat.

\textbf{The result is a near-perfect split.} The $41$ OpenRouter-style runs return $237/237$ ($100\%$) usable micro-question answers across every model tested, with zero exceptions. The $22$ local-style runs return $6/966$ ($0.6\%$) usable answers; every one of the three July production campaigns is included in this group and each contributes exactly zero. The single partial exception is a $6/6$-successful run dated 8 May 2026 (\texttt{benchmark\_model: gemma27b}) launched from an HPC-resident Python environment; we cannot determine from its manifest alone whether its evaluator call was still cloud-routed, so we report it as an open exception rather than folding it into either count's interpretation. Excluding that one run, local serving is a \emph{total} failure and cloud serving is a \emph{total} success for this specific extraction task, across at least 15 distinct underlying models on the OpenRouter side and at least 4 distinct model families on the local side.

This resolves the question raised in earlier drafts of this section: the epistemic-mapping mechanism is not broken by design, and question generation, prompting, and rationale extraction all demonstrably work when the evaluator call is served through OpenRouter. \texttt{evaluator\_fallback\_used} remains \texttt{False} for every affected local unit, so the pipeline's own error detection does not distinguish this near-total local failure from a successful call---the same silent-placeholder pattern documented for the JSD telemetry (Section~\ref{sec:jsd_convergence}).

\textbf{Source-code audit: the actual mechanism, and a correction to an earlier draft of this section.} An earlier draft of this paper attributed the split above to a base-URL misconfiguration (\texttt{OPENROUTER\_BASE\_URL} pointed at the wrong local server). Having subsequently obtained access to the evaluator's source, we withdraw that explanation: \texttt{BenchmarkRoleRouter.client\_for("evaluator")} (\texttt{role\_router.py}) correctly prioritises \texttt{EVALUATOR\_LLM\_BASE\_URL} for the evaluator role, so the evaluator's requests are not misrouted. The real mechanism is different, is confirmed directly from source rather than inferred, and implicates a keyword-matching heuristic rather than the evaluator model's competence. In \texttt{LLMClient.chat\_json()} (\texttt{llm\_client.py}):
\begin{quote}
\texttt{is\_reasoning\_model = self.model and any(kw in self.model.lower() for kw in ("r1", "reasoning", "deepseek-r1"))} \\
\texttt{if is\_reasoning\_model: json\_schema = None}
\end{quote}
followed by
\begin{quote}
\texttt{if (is\_local\_vllm and is\_quantized) or is\_reasoning\_model: response\_format = None; use\_schema = False}
\end{quote}
The fixed evaluator's model identifier, \texttt{"deepseek-r1-14b"}, contains \texttt{"r1"} and matches this filter. Consequently \emph{every} evaluator call in all three production campaigns runs with \texttt{response\_format = None}---not merely without \texttt{json\_schema} structured-output enforcement, but without even the weaker \texttt{json\_object} guarantee. The model must reproduce the full nested schema (\texttt{probabilities}, \texttt{mcq\_dimensions}, \texttt{validated\_scales}, and, when applicable, \texttt{micro\_epistemic\_mapping}) from natural-language instructions alone, in a single unconstrained generation.

Two further details, also read directly from \texttt{ProbabilityEvaluator.evaluate()} (\texttt{evaluator.py}), explain why the four fields fail differently rather than uniformly. First, a retry loop (up to three attempts, with escalating temperature and error feedback) checks response completeness before accepting it---but the check tests only \texttt{probabilities} and \texttt{mcq\_dimensions}:
\begin{quote}
\texttt{if not isinstance(probs, dict) or not probs or ...: raise ValueError(...)} \\
\texttt{if not isinstance(mcq, dict) or not mcq or len(mcq) < len(MCQ\_DIMENSION\_KEYS): raise ValueError(...)}
\end{quote}
\texttt{micro\_epistemic\_mapping} is never examined by this check, so a response with an empty or incomplete mapping is accepted as successful and the loop exits without ever prompting the model to retry that specific field. Second, \texttt{\_validate\_micro\_mapping()} then silently backfills every micro-question the returned mapping does not cover with \texttt{\{"dominant\_tag": "N/A", "short\_rationale": "Missing in LLM response"\}}, which is written into the result as if it were a normal value---no exception is raised and \texttt{evaluator\_fallback\_used} is never set. \texttt{probabilities} and \texttt{mcq\_dimensions}, by contrast, are retried and are the two fields most directly and repeatedly emphasised in the prompt (\S\ref{sec:grading}), so they are structurally present in the released data even without schema enforcement; \texttt{micro\_epistemic\_mapping} additionally requires the model to reproduce \emph{event-specific, dynamically named} keys (e.g.\ \texttt{S3\_q1\_pres\_prob}) with no structural scaffolding forcing it to do so, which is a plausible reason it is the field most often dropped entirely.

This also reframes the OpenRouter/local split reported above: it does not appear to be a local-versus-cloud serving effect as such. \texttt{is\_reasoning\_model} is evaluated on \emph{whichever model string is configured for the evaluator role}, regardless of where that model is served. In the 63-run audit, the production campaigns' evaluator was fixed to \texttt{"deepseek-r1-14b"} throughout (Section~\ref{sec:confound}), while the OpenRouter-style pilot runs generally used the model under test itself as its own evaluator (\texttt{OPENROUTER\_EVALUATOR\_MODEL} set to the same value as \texttt{OPENROUTER\_BENCHMARK\_MODEL}), and none of \texttt{qwen/\allowbreak qwen-2.5-7b-instruct}, \texttt{google/\allowbreak gemma-4-26b-a4b-it}, \texttt{openai/\allowbreak gpt-oss-120b}, \texttt{anthropic/\allowbreak claude-3-haiku}, or \texttt{deepseek/\allowbreak deepseek-v3.2} matches the \texttt{"r1"/"reasoning"} filter. Serving location and schema-enforcement status were confounded in this dataset because no run in our audit paired a reasoning-named evaluator with cloud serving, or a non-reasoning-named evaluator with local serving; we cannot rule out an additional, independent local-serving effect, but the schema-disabling filter is sufficient on its own to explain the near-total local failure and does not require one.

\textbf{This same mechanism plausibly explains the low inter-rater reliability of the 7-dimension rubric (Section~\ref{sec:evaluator_reliability}), not just the missing micro-question data.} The pipeline computes that reliability by calling \texttt{evaluator.evaluate()} \emph{twice} per unit (\texttt{run1}, \texttt{run2}; \texttt{run\_ecnbench\_protocol.py}) and comparing the dominant bucket label each run assigns to each dimension---both calls use the identical evaluator client, and both therefore run with \texttt{response\_format = None} for the same reason. What Cohen's $\kappa$ measures in this campaign is consequently \emph{generation-to-generation consistency of an unconstrained model reproducing a bucket label from a natural-language instruction}, not disagreement between genuinely independent raters about an underlying construct. Unlike \texttt{micro\_epistemic\_mapping}, \texttt{mcq\_dimensions} is retried until structurally complete, so it is present in both runs, but nothing in this design constrains its \emph{content} to be stable across the two generations---which is precisely the signature low-$\kappa$, structurally-complete-but-noisy pattern reported in Section~\ref{sec:evaluator_reliability} ($\max \kappa = 0.322$, every dimension below the reliability threshold). We did not rerun the campaign with schema enforcement restored to confirm this quantitatively, and flag doing so as a priority diagnostic (Section~\ref{sec:conclusion}), but the mechanism is a direct, verified reading of the scoring code rather than a new hypothesis.

\textbf{What this audit does and does not settle, and what the rest of Part~I adds.} The defect above is established directly from source and is not in doubt. What it cannot do on its own is explain the flat-forecast finding of Section~\ref{sec:degenerate_flat}. The final \texttt{probabilities} used for Brier scoring are produced by this same schema-disabled call, so the absence of structural enforcement is a plausible contributing condition there too; but \texttt{is\_reasoning\_model} is constant across all three production campaigns (the evaluator is always \texttt{"deepseek-r1-14b"}), whereas the flat-forecast rate is not (Llama-3.1-8B and Mistral-7B are affected, Qwen2.5-7B is not), so schema-disabling alone cannot be the \emph{differentiating} cause. The audit therefore establishes the evaluator as a suspect but not as the cause, and it supersedes the vLLM-versus-Ollama framing tentatively offered in Section~\ref{sec:degenerate_flat}, which the source does not support as the operative mechanism. The remaining candidate the source does suggest is \texttt{evidence\_text} volume, which differs sharply by swarm model (up to $20{,}329$ characters for Llama-3.1-8B versus a fraction of that for the near-dormant Qwen2.5-7B social channel, Appendix~\ref{sec:interaction_realism}): an unconstrained evaluator asked to compress more, and more heterogeneous, discussion content into a single distribution may hedge toward uniform more often. Earlier drafts left that as an untested hypothesis. Section~\ref{sec:evaluator_replication} tests it by changing the instrument while holding the transcripts fixed, and Section~\ref{sec:volume_mechanism} measures the dose--response directly; together they convert it into the paper's central mechanism.

\section{Evaluator Replication: Is Degenerate Consensus a Property of the Swarm or of the Instrument?}
\label{sec:evaluator_replication}

Section~\ref{sec:degenerate_flat} left one question open and named it the highest-priority follow-up: the flat-forecast rate differs sharply between swarm models, yet the mechanism we can confirm at the level of source code---an evaluator whose structured-output enforcement is silently disabled by a keyword filter---is \emph{constant} across all three campaigns and therefore cannot explain a difference between them. That audit established the evaluator as a suspect but could not establish it as a cause, because the campaigns contain only one evaluator pass. This section closes that gap with the experiment the audit implied: we re-read the same simulation transcripts with a different evaluator and ask which of the paper's findings survive.

\subsection{Design}
\label{sec:replication_design}

The evaluator's input is the per-unit \texttt{evidence\_text} record---the seed dossier for Condition A, the simulation transcript for Conditions B and C---together with the event question and option list. Those records are fixed artifacts of the July campaigns. Re-running only the evaluator over them therefore holds the entire simulation constant and varies nothing but the measuring instrument. We denote the two passes:

\begin{itemize}
    \item \textbf{Evaluator J} (July 2026): DeepSeek-R1-14B, served locally, structured-output enforcement disabled by the keyword filter documented in Section~\ref{sec:epistemic_mapping}. This is the instrument behind every number reported in Sections~\ref{sec:full_campaign}--\ref{sec:degenerate}.
    \item \textbf{Evaluator A} (August 2026): \texttt{deepseek/deepseek-v4-flash-0731} via the OpenRouter API, JSON-object enforcement active, temperature $0.0$ and seed $42$ pinned by the benchmark's determinism harness.
\end{itemize}

All $270$ units of the three scorable campaigns were re-evaluated. Every derived quantity reported below---Brier, RPS, directional accuracy, calibration bracket---is recomputed from each pass's own \texttt{probabilities} vector using the protocol runner's own scoring functions, so no comparison can be contaminated by a metric left over from the other pass.\footnote{This is not a hypothetical hazard. The re-evaluation script updates \texttt{probabilities} but not the derived metrics, which left the working copies internally inconsistent in $82$--$88$ of $90$ units per campaign until the two passes were separated into distinct files. Both datasets are released separately and each is internally consistent in $90/90$ units; the reconciliation utility and its audit output are included in the release (Appendix~\ref{sec:data_audit}).}

Two notes on the conduct of the replication. First, the pass was initially launched against a stale copy of the event catalogue in which the identifiers C6, T1, T3, T4, T5, T6 and T8 carry entirely different events, so $21$ of $90$ units per campaign received another event's micro-questions. Those $63$ units were re-evaluated against the authoritative catalogue and every unit in the released Dataset A now carries the micro-question set its own event defines; the correction is verifiable with the released checker, which reports zero remaining mismatches. Re-running them also changed the headline degeneracy count, which we treat as evidence about the instrument rather than a footnote, and discuss in Section~\ref{sec:replication_central}. Second, Qwen2.5-14B remains unscorable under either evaluator: $87$ of its $90$ records carry no \texttt{ground\_truth} value, so it is excluded here exactly as in Section~\ref{sec:full_campaign}, and its micro-question channel was left uncorrected for the same reason.

\subsection{The Central Finding Does Not Replicate}
\label{sec:replication_central}

Table~\ref{tab:replication} and Figure~\ref{fig:replication} give the result. Because the two passes must read identical input for the contrast to isolate the instrument, we fix the scope before quoting any number: Mistral-7B and Qwen2.5-7B contribute all $30$ of their events, while Llama-3.1-8B contributes only the six the backend outage never touched (C3, S2--S6; \S\ref{sec:outage}). Its remaining $24$ events were re-simulated on 2026-08-16 and no July scoring of them exists, so they cannot enter a J-versus-A comparison at all and are reported instead in the table's third column, under the August evaluator alone.

On the strictly paired Mistral-plus-Qwen material the flat-forecast rate over simulated units falls from $19/120$ ($15.8\%$) under Evaluator J to $1/120$ ($0.8\%$) under Evaluator A, and the cross-model pattern that Section~\ref{sec:degenerate_flat} set out to explain does not survive: Condition B counts of $7/30$ and $0/30$ become $0/30$ and $0/30$. Degeneracy is eliminated outright for Mistral-7B, so the contrast between the vLLM-served Mistral and the Ollama-served Qwen---the observation the evidence-volume hypothesis was introduced to account for---is not a stable property of the transcripts. The single surviving flat unit is a Mistral Condition C unit. Condition C is the null-injection arm, in which the swarm is fed structurally identical but topically unrelated evidence, and a uniform posterior is the \emph{correct} response to evidence carrying no information about the question, so not one unconstrained-evaluator flat forecast in a condition where information \emph{was} injected survives the constrained instrument in either campaign.

Llama's six paired events make the same point far more sharply on far less data. Under Evaluator J every one of their $12$ simulated units is flat---$6/6$ in Condition B and $6/6$ in Condition C---and one of the six Condition A units is flat as well. Under Evaluator A none of the $13$ is. That is a total reversal on byte-identical text, and it is the single clearest instance of the instrument effect anywhere in this paper, but the subset it rests on is small and not randomly chosen: these are exactly the units that ran before the serving failure, so they are also Llama's only high-evidence-volume units (Appendix~\ref{sec:evidence_depth}), and flatness is confounded with evidence volume within the campaign. We therefore read the six events as a demonstration that the effect exists and can be complete, not as an estimate of how large it is.

A further, sharper lesson about the instrument came from an accident of our own analysis, which we record because it caught us mid-draft. An earlier version of this section reported $0/30$ Condition B flat units for all three models. Those numbers came from an Evaluator A pass in which $21$ units per campaign had been scored against the wrong event's micro-questions (Section~\ref{sec:replication_design}); re-running exactly those units against the correct catalogue changed Llama-3.1-8B's Condition B count over its then-30-event set from $0$ to $4$, and moved its Condition C count from $9$ to $6$. Two things differ between those passes---the micro-question channel was corrected, and the units were sampled afresh from a non-deterministic evaluator---and this design cannot separate them. Either way the conclusion is the one Section~\ref{sec:replication_noise} quantifies directly: a re-run of the same evaluator on the same transcripts moved the paper's headline degeneracy count by four units out of thirty. A single evaluator pass, of either instrument, does not measure this quantity to the precision at which the paper's original cross-model comparison was read. Those particular counts are historical: they were taken over Llama's outage-corrupted transcripts, which the repair of \S\ref{sec:outage} has since replaced, and they are quoted here as evidence about the evaluator's stability rather than as measurements of the campaign.

The milder sharpness effect attenuates in the same direction without vanishing, and it supplies its own instrument check. Under Evaluator J the Condition A mean modal probability is a single number common to all three campaigns by construction---$0.658$ over $30$ events, $0.617$ over Llama's six---because Condition A is one shared No-Sim record (Section~\ref{sec:condition_a}). Under Evaluator A that single input yields $0.713$, $0.668$ and $0.672$ on those same six events, and $0.722$ and $0.687$ for Mistral and Qwen over all $30$: the constrained instrument does not return one answer to one text, which is the instrument noise Section~\ref{sec:replication_noise} quantifies. Each within-model Condition A$\to$B drop is therefore computed against its own campaign's Condition A. Those drops shrink from $-0.200$, $-0.173$ and $-0.120$ under Evaluator J to $+0.057$, $-0.045$ and $-0.018$ under Evaluator A---for Llama the drop does not merely shrink but reverses, the swarm's simulated forecast coming out \emph{sharper} than the No-Sim reference. The aggregate lift reversal that motivated the whole investigation is no longer uniform in sign: $BS_A - BS_B$ is $-0.033$ $[-0.685, +0.544]$ for Llama-3.1-8B on its six paired events, $-0.260$ $[-0.465, -0.060]$ for Mistral-7B and $+0.035$ $[-0.103, +0.182]$ for Qwen2.5-7B, and only Mistral's interval excludes zero---the model already identified in Section~\ref{sec:serving_note} as confounded by release generation. On the separate recovery scope, Llama's $24$ re-simulated events give $-0.235$ $[-0.375, -0.100]$ under the same August evaluator; that interval excludes zero, but it is a measurement of a different quantity---a model on new transcripts, not an instrument on fixed ones---and Appendix~\ref{sec:cutoff} rather than this section is where it is interpreted.

\begin{table*}[!htbp]
\centering
\caption{The paper's headline quantities under three evaluator passes, with the comparison scope stated per column. \textbf{Evaluator J} is the July instrument (DeepSeek-R1-14B, served locally) behind every number in Sections~\ref{sec:full_campaign}--\ref{sec:degenerate}; \textbf{Evaluator A} is the August re-read of the \emph{same} transcripts by \texttt{deepseek/deepseek-v4-flash-0731}; \textbf{Evaluator OR} is that same model scoring Llama-3.1-8B's re-simulated transcripts on 2026-08-16. The J-versus-A contrast isolates the instrument only where both passes read byte-identical input. That holds for all $30$ events of Mistral-7B and Qwen2.5-7B, but for Llama-3.1-8B only on the six events the backend outage never touched (C3, S2--S6); its other $24$ were re-simulated (\S\ref{sec:outage}) and appear in the OR column instead. The two scopes answer different questions---one measures the instrument at fixed input, the other measures the model at fixed instrument---and are never pooled. Bracketed intervals are $95\%$ percentile bootstrap CIs over events ($20{,}000$ resamples, seed $7$), recomputed from the per-unit \texttt{probabilities} vectors.}
\label{tab:replication}
\scriptsize
\setlength{\tabcolsep}{2.5pt}
\begin{tabular}{llccc}
\toprule
 & & \textbf{Evaluator J} & \textbf{Evaluator A} & \textbf{Evaluator OR} \\
\textbf{Quantity} & \textbf{Model} & \textbf{(Jul, local)} & \textbf{(Aug, same input)} & \textbf{(Aug 16, re-simulated)} \\
 & & \multicolumn{2}{c}{\emph{paired instrument scope}} & \emph{recovery scope} \\
\midrule
\multirow{2}{*}{Flat units, simulated (B$+$C)} & Mistral $+$ Qwen & $19/120$ ($15.8\%$) & $\mathbf{1/120}$ ($0.8\%$) & --- \\
 & Llama-3.1-8B\textsuperscript{$\dagger$} & $12/12$ & $\mathbf{0/12}$ & $1/48$ \\
\multirow{2}{*}{Flat units, No-Sim (A)} & Mistral $+$ Qwen & $2/60$ ($3.3\%$) & $\mathbf{0/60}$ ($0.0\%$) & --- \\
 & Llama-3.1-8B\textsuperscript{$\dagger$} & $1/6$ & $\mathbf{0/6}$ & $0/24$ \\
\midrule
\multirow{3}{*}{Flat units, Condition B} & Llama-3.1-8B\textsuperscript{$\dagger$} & $6/6$ & $\mathbf{0/6}$ & $0/24$ \\
 & Mistral-7B & $7/30$ & $\mathbf{0/30}$ & --- \\
 & Qwen2.5-7B & $0/30$ & $\mathbf{0/30}$ & --- \\
\midrule
\multirow{3}{*}{Flat units, Condition C} & Llama-3.1-8B\textsuperscript{$\dagger$} & $6/6$ & $\mathbf{0/6}$ & $1/24$ \\
 & Mistral-7B & $12/30$ & $1/30$ & --- \\
 & Qwen2.5-7B & $0/30$ & $\mathbf{0/30}$ & --- \\
\midrule
\multirow{3}{*}{Mean modal probability, A} & Llama-3.1-8B\textsuperscript{$\dagger$} & $0.617$ & $0.713$ & $0.719$ \\
 & Mistral-7B & $0.658$ & $0.722$ & --- \\
 & Qwen2.5-7B & $0.658$ & $0.687$ & --- \\
\multirow{3}{*}{Mean modal probability, B} & Llama-3.1-8B\textsuperscript{$\dagger$} & $0.417$ & $0.770$ & $0.531$ \\
 & Mistral-7B & $0.485$ & $0.676$ & --- \\
 & Qwen2.5-7B & $0.538$ & $0.669$ & --- \\
\midrule
\multirow{3}{*}{Lift $BS_A - BS_B$} & Llama-3.1-8B\textsuperscript{$\dagger$} & $-0.161$ $[-0.402, +0.060]$ & $-0.033$ $[-0.685, +0.544]$ & $\mathbf{-0.235}$ $[-0.375, -0.100]$ \\
 & Mistral-7B & $-0.147$ $[-0.313, +0.028]$ & $\mathbf{-0.260}$ $[-0.465, -0.060]$ & --- \\
 & Qwen2.5-7B & $-0.097$ $[-0.249, +0.063]$ & $+0.035$ $[-0.103, +0.182]$ & --- \\
\midrule
\multirow{3}{*}{Susceptibility $BS_C - BS_B$} & Llama-3.1-8B\textsuperscript{$\dagger$} & $+0.000$ $[+0.000, +0.000]$ & $+0.472$ $[-0.056, +1.008]$ & $-0.002$ $[-0.147, +0.133]$ \\
 & Mistral-7B & $\mathbf{+0.115}$ $[+0.009, +0.225]$ & $-0.051$ $[-0.287, +0.168]$ & --- \\
 & Qwen2.5-7B & $+0.053$ $[-0.064, +0.172]$ & $\mathbf{+0.106}$ $[+0.020, +0.192]$ & --- \\
\midrule
\multirow{2}{*}{$r(BS_A,\text{lift})$ / mech. / $r(BS_A, BS_B)$} & Mistral $+$ Qwen & $+0.774$ / $+0.809$ / $+0.055$ & $+0.573$ / $+0.619$ / $+0.056$ & --- \\
 & Llama-3.1-8B\textsuperscript{$\dagger$} & --- & --- & $+0.505$ / $+0.889$ / $+0.407$ \\
\bottomrule
\end{tabular}
\smallskip\\
\footnotesize\textsuperscript{$\dagger$}Llama-3.1-8B's J and A columns cover the six outage-free events only ($n=6$: C3, S2--S6); its OR column covers the $24$ re-simulated events ($n=24$). Those six events are exactly the degenerate core described in \S\ref{sec:cutoff}, which is why Evaluator J returns a uniform vector for $6/6$ of them in both simulated conditions while Evaluator A returns a non-uniform vector for every one---the sharpest single instance of the instrument effect in the paper, but measured on a small and adversarially selected subset, so its bootstrap intervals are correspondingly wide and we draw no effect-size conclusion from them. The two Llama scopes must not be pooled or read as a time series. The OR column additionally carries the serving-stack substitution of \S\ref{sec:outage} and the raised evaluator completion budget documented in item~(5) of \S\ref{sec:outage}; round-level convergence telemetry was not extracted from the new traces, so the JSD quantities of \S\ref{sec:jsd_convergence} are \emph{unavailable} for those $24$ events rather than placeholder-valued.
\setlength{\tabcolsep}{6pt}
\end{table*}

\begin{figure*}[!htbp]
\centering
\includegraphics[width=\textwidth]{figR1_evaluator_replication.png}
\caption{Degenerate consensus under two evaluators, with the comparison scope held fixed. \textbf{Panel A:} flat-forecast counts by condition on the paired material only---all $30$ events of Mistral-7B and Qwen2.5-7B plus the six events of Llama-3.1-8B the outage never touched---where Evaluator J and Evaluator A read byte-identical text. Under the unconstrained July instrument $13/66$ Condition B units---the arm in which real information was injected---are exactly uniform, along with $18/66$ in Condition C and $3/66$ in the No-Sim Condition A. Under the constrained August instrument the counts are $0$, $1$ and $0$: the single survivor is a Condition C unit, the null-injection arm in which a uniform posterior is the warranted answer. \textbf{Panel B:} the B-versus-C susceptibility margin $BS_C - BS_B$ per event on Llama-3.1-8B's $24$ re-simulated events under \texttt{deepseek-v4-flash-0731}, the scope in which that model's benchmark result is actually measured. Thirteen of the $24$ margins are positive and the mean is $-0.002$ $[-0.147, +0.133]$: centred on zero, with an interval wider than the effect the benchmark was designed to detect and with per-event values ranging from $-0.85$ to $+0.67$. The two panels are deliberately on different scopes and the reader should not trace a model across them.}
\label{fig:replication}
\end{figure*}

\subsection{Neither Pass Is a Reference Measurement}
\label{sec:replication_noise}

It would be a mistake to read Section~\ref{sec:replication_central} as showing that Evaluator A is correct and Evaluator J was broken. Evaluator A is itself an unstable instrument, and the campaigns contain a clean way to demonstrate it.

Condition A carries no simulation: its \texttt{evidence\_text} is the seed dossier, and it is byte-identical across all three campaigns for all $30$ events. This remains true after the outage repair, which restored Condition A from the shared record rather than re-eliciting it: we verified that the Condition A rows of both the primary and the replication result files hash identically across the three campaigns on $30/30$ events, and that the repair left Llama's Condition A probability vectors bit-for-bit unchanged in both files. Because the re-evaluation was launched separately per campaign, Evaluator A therefore read each of those $30$ inputs three independent times, at temperature $0.0$ with seed $42$ pinned. If the instrument were reproducible, the three readings would agree exactly. They agree exactly for $4$ of $30$ events. The mean spread in Brier across the three re-reads is $0.129$ and the maximum is $0.715$---a range exceeding the entire between-condition effect the paper reports. Figure~\ref{fig:instrument_noise}B shows the full distribution. The corresponding count for Evaluator J is $30/30$, but this carries no information about its stability: Condition A was executed once and copied verbatim into all three campaigns (Section~\ref{sec:condition_a}), so the agreement is an artifact of file reuse rather than a replication.

Agreement between the two passes is likewise only moderate. Over the $198$ units of the paired scope, $r(BS_J, BS_A) = +0.552$, and the per-model figures are strikingly uniform: $+0.589$ for Llama-3.1-8B (on its $18$ paired units), $+0.587$ for Mistral-7B and $+0.549$ for Qwen2.5-7B. An earlier draft reported a pooled $+0.415$ with Llama at $+0.144$ and attributed the gap to something specific about that campaign; that comparison was invalid, because it scored Evaluator J's July readings against Evaluator A's readings of the \emph{re-simulated} transcripts, so it measured the repair rather than the instrument. Restricted to matched input the anomaly disappears, and the honest statement is the stronger one: the two instruments correlate at about $+0.55$ regardless of which model's transcripts they read. Mean Brier falls under Evaluator A for the two campaigns measured over all $30$ events ($0.521 \to 0.409$ for Mistral, $0.467 \to 0.334$ for Qwen) and \emph{rises} for Llama on its six paired events ($0.530 \to 0.601$), driven by Condition C, where replacing six uniform vectors with confident ones that happen to be wrong takes the mean from $0.583$ to $0.927$. Either way the choice of evaluator moves the benchmark's headline score by $0.07$ to $0.13$ Brier---larger than any of the between-condition lifts the benchmark was built to detect, and not even reliably in one direction. Agreement on the modal option, counting a tie as a disagreement, is $5/18$ for Llama, $50/90$ for Mistral and $63/90$ for Qwen, $118/198$ pooled: on a little over half the units the two instruments do not even name the same outcome as most likely.

\begin{figure*}[!htbp]
\centering
\includegraphics[width=\textwidth]{figR2_instrument_noise.png}
\caption{The evaluator as a source of variance. \textbf{Panel A:} per-unit Brier under the two evaluators, restricted to the $198$ units both passes read byte-identically---all $30$ events of Mistral-7B and Qwen2.5-7B plus Llama-3.1-8B's six outage-free events (\S\ref{sec:outage}); agreement is moderate ($r = +0.552$) and the scatter is not concentrated near the identity line. \textbf{Panel B:} Condition A's seed dossier is byte-identical across the three campaigns, so the August evaluator read each of the $30$ inputs three times at temperature $0$. Only $4$ events reproduce exactly; the largest Brier spread across three re-reads of the same text is $0.715$.}
\label{fig:instrument_noise}
\end{figure*}

\subsection{What This Means for the Paper's Claims}
\label{sec:replication_implications}

We read this replication as answering the question posed in Section~\ref{sec:degenerate_flat}, and answering it against the interpretation that section preferred.

\textbf{Degenerate consensus is substantially a property of the measurement, not of the deliberation.} The flat forecasts that motivated this paper are produced when an unconstrained evaluator is asked to compress a long transcript into a categorical distribution. On the strictly paired material a second evaluator recovers one twentieth of them ($1/120$ against $19/120$ for Mistral and Qwen, $0/12$ against $12/12$ for Llama's six outage-free events) and does not reproduce their cross-model structure at all: the Condition B counts of $6/6$, $7/30$ and $0/30$ across Llama, Mistral and Qwen all become zero. The audit in Section~\ref{sec:epistemic_mapping} identified the disabled structured-output enforcement as the confirmed source-level defect but could not attribute the cross-model difference to it. Section~\ref{sec:replication_central} supplies the missing evidence, in the negative direction the audit suspected: with enforcement active the difference the evidence-volume hypothesis was built to explain is no longer there to explain. We accordingly restate the paper's central finding as a finding about instrumentation: \emph{iterative multi-agent deliberation produces transcripts that an unconstrained LLM evaluator frequently fails to reduce to an informative distribution, and that failure is indistinguishable, in the aggregate Brier score, from the swarm having reached no consensus.} The phenomenon is real, reproducible from the released artifacts, and consequential for anyone scoring a multi-agent simulation this way. We stop short of locating it wholly in the evaluator, but the evidence for a residual is now weaker than in earlier drafts: no Condition B unit anywhere---not in the $180$ paired units, not in the $24$ re-simulated Llama events---returns a uniform vector under the constrained instrument, and the two flat units that do survive are both in Condition C, where uniformity is the warranted answer. A contribution from the deliberation itself is not excluded; it is simply no longer visible above an instrument with this reproducibility.

\textbf{The benchmark's three original hypotheses are not identifiable at this evaluator noise level.} Section~\ref{sec:cross_model_stats} already reported $\leq 9\%$ achieved power at $E = 30$. Table~\ref{tab:replication} adds a second, independent obstacle: switching evaluator on fixed transcripts moves Qwen2.5-7B's lift from $-0.097$ to $+0.035$---a sign reversal---and moves the set of models whose susceptibility interval excludes zero from $\{$Mistral$\}$ to $\{$Qwen$\}$, i.e.\ to a disjoint set. Llama's susceptibility over the same six events goes from exactly $+0.000$ (both simulated conditions uniform, so the contrast is identically zero by construction) to $+0.472$ with an interval that still spans zero. Both passes are single measurements with an instrument whose reproducibility on byte-identical input is $4/30$. Neither the lift claim, the convergence claim nor the susceptibility claim can be settled by either dataset, and we withdraw any reading of Section~\ref{sec:cross_model_stats} that treats the sign of a single-pass lift as evidence about deliberation.

\textbf{What survives unchanged.} Three results in this paper do not depend on which evaluator read the transcripts. The JSD convergence telemetry defect is simulation-side: the constant $0.7583$ placeholder occupies $87.7\%$, $40.7\%$ and $27.7\%$ of Llama, Mistral and Qwen checkpoints in the raw records, exactly as reported in Section~\ref{sec:jsd_convergence}. (Llama's figure describes the archived pre-repair campaign, whose $24$ outage-affected events \S\ref{sec:outage} has since replaced; the re-simulation did not extract round-level telemetry, so those events now carry no JSD series at all rather than a placeholder one, and the defect is neither reproduced nor refuted by the new data.) The difficulty--lift correlation remains an arithmetic artifact under every pass: pooled over the $60$ strictly paired Mistral and Qwen events, $r(BS_A, \text{lift})$ is $+0.774$ against a mechanical expectation of $+0.809$ under Evaluator J and $+0.573$ against $+0.619$ under Evaluator A, while the non-tautological association $r(BS_A, BS_B)$ stays negligible ($+0.055$ and $+0.056$). One qualification belongs here rather than in a footnote: on Llama's $24$ re-simulated events the observed $r(BS_A, \text{lift})$ of $+0.505$ falls well \emph{below} its mechanical expectation of $+0.889$, because $r(BS_A, BS_B) = +0.407$ there---the first campaign in this paper in which the two conditions' difficulties are substantially associated. At $n=24$ on a single evaluator pass we do not read that as a finding, but it is the one place where the tautology argument does not fully account for the observed correlation, and it should be checked rather than assumed in any follow-up. And the provenance confound of Section~\ref{sec:serving_note}---three models differing simultaneously in release generation, precision and serving stack---is a property of the experimental design, not of any measurement taken from it.

\textbf{Recommendation for future campaigns.} A benchmark that extracts probabilities with an LLM must treat that extractor as an instrument requiring calibration, not as an oracle. Table~\ref{tab:benchmark_checklist} states this as a checklist: four gates we would now block a campaign on, and three quantities we would now refuse to interpret a score without. Each row is stated together with the specific failure in this paper it would have caught, because every one of them is a lesson we paid for rather than a generic recommendation. Four points deserve emphasis in prose. Verify at runtime that structured-output enforcement is actually in force rather than assuming it from a request parameter, since the defect documented in Section~\ref{sec:epistemic_mapping} was a silent keyword-match bypass. Derive completion flags from the artifacts rather than from the orchestrator, since a pipeline that sets its own success fields will report a dead backend as a successful experiment. Run the evaluator $k \geq 3$ times per unit and report the between-run spread alongside every score, so that a $0.715$ Brier range on identical input cannot be mistaken for an effect. And score every campaign with at least two independent evaluators, treating any finding that does not survive both as a property of the instrument until shown otherwise. We single out the flat-unit rate as the cheapest useful diagnostic of the set: it separated a measurement artifact from a substantive result here, and it costs nothing to compute from data every such benchmark already has.

\begin{table*}[!htbp]
\centering
\caption{A checklist for benchmarks that extract probabilities with an LLM. Each row states a check, the failure it would have caught \emph{in this paper}, and what it costs to run. Every one of these was learned by having the corresponding failure go undetected here; the ``caught it here'' column cites the section where the damage is quantified, so the checklist is an account of this campaign rather than a set of generic best practices. We recommend treating the first four as blocking gates on any campaign that reports a score, and the last three as reporting requirements.}
\label{tab:benchmark_checklist}
\footnotesize
\setlength{\tabcolsep}{4pt}
% Narrow p{} columns cannot be justified without loose lines; ragged right.
\begin{tabular}{p{0.06\textwidth}>{\raggedright\arraybackslash}p{0.29\textwidth}>{\raggedright\arraybackslash}p{0.40\textwidth}>{\raggedright\arraybackslash}p{0.17\textwidth}}
\toprule
& \textbf{Check} & \textbf{What it would have caught here} & \textbf{Cost} \\
\midrule
\multicolumn{4}{l}{\emph{Blocking gates --- run before any number leaves the pipeline}} \\
\midrule
G1 & \textbf{Verify structured-output enforcement at runtime}, by asserting on the request actually sent, not on the parameter you passed. & A keyword filter matching \texttt{"r1"} in the evaluator's model id silently dropped \texttt{response\_format} for every unit of every July campaign, and nothing in the logs or the result records recorded it (\S\ref{sec:epistemic_mapping}). We list this as a gate because a pipeline must know what it sent, \emph{not} because we can show enforcement prevents degeneracy: isolated within one evaluator it does not (\S\ref{sec:truncation_grid}). & One assertion per request. \\
\addlinespace
G2 & \textbf{Derive completion flags from the artifacts}, never from the orchestrator. Require a minimum count of successful telemetry probes before a unit is scored. & $69$ units recorded \texttt{simulation\_status: "completed"}, \texttt{error: null} and \texttt{simulation\_executed: true} while their serving backend had been dead for $23$ events (\S\ref{sec:outage}). The failure was invisible in every field a consumer would check. & One count per unit. \\
\addlinespace
G3 & \textbf{Assert that the scored record's question and options match the catalogue entry} its event id resolves to, by value and not by id. & Seven event ids denote different events in two catalogue copies; two of those collisions (T1, T5) share option set \emph{and} ground truth with the intended events, so they pass every id-level and schema-level check while simulating an unrelated question (\S\ref{sec:outage}). & One string comparison per unit. \\
\addlinespace
G4 & \textbf{Assert that the evidence field is non-empty and of plausible length} before scoring, and log its length with the score. & A repair script read \texttt{question} and \texttt{evidence\_text} from a config file that stores neither, so $72$ units were scored on an empty evidence string---the three conditions differing by one character---and produced two publishable-looking findings before being caught (\S\ref{sec:outage}). & One length check per unit. \\
\midrule
\multicolumn{4}{l}{\emph{Reporting requirements --- run before any number is interpreted}} \\
\midrule
R1 & \textbf{Report the flat-unit rate} (share of forecasts within $\epsilon$ of uniform) per condition, always, alongside the aggregate score. & It is the single diagnostic that separated a measurement artifact from a substantive result here, and the only one that made the artifact legible as an artifact rather than as a null effect (\S\ref{sec:degenerate_flat}). & Free; it is a function of data you already have. \\
\addlinespace
R2 & \textbf{Run the evaluator $k \geq 3$ times per unit} and report the between-run spread with every score. & Three re-reads of byte-identical input at temperature $0$ reproduce exactly $4$ of $30$ units, with a maximum Brier spread of $0.715$---larger than any between-condition effect the benchmark was built to detect (\S\ref{sec:replication_noise}). & $k\times$ evaluator cost. \\
\addlinespace
R3 & \textbf{Score with at least two independent evaluators} and treat any finding that does not survive both as a property of the instrument until shown otherwise. & Switching evaluator on fixed transcripts reverses the sign of one model's lift and moves the set of models with a susceptibility interval excluding zero to a disjoint set (\S\ref{sec:replication_central}). & $2\times$ evaluator cost. \\
\bottomrule
\end{tabular}
\setlength{\tabcolsep}{6pt}
\end{table*}

\section{Carrying Out the Calibration Protocol: Four Models, Five Reads, Contamination Controlled}
\label{sec:four_model_extension}

\confoundnote{This is the best-controlled comparison in the paper---five
evaluator reads per unit, contamination-clean events only---and it is still not a
controlled model comparison. The four campaigns differ in release generation,
precision and serving stack (Section~\ref{sec:serving_note}): Llama-3.1-8B is
AWQ-INT4 under vLLM, Mistral-7B unquantised FP16 under vLLM, and both Qwen models
GGUF under Ollama. Qwen2.5-7B and Qwen2.5-14B are the only pair sharing family,
release date and inference server, and even there the $14$B campaign's exact
quantization level is not recorded in its manifest, so a like-for-like reading is
available for that pair alone and only approximately. Rankings below are reported
as descriptive; $20$ events also leave every per-model difference underpowered,
which is why the analysis of Section~\ref{sec:predictive_signal} rests on the
pooled figure rather than on any single row.}

We carried out the preceding recommendation in full. Every unit of the three originally-scorable campaigns was re-read three further times with Evaluator A (\texttt{deepseek/deepseek-v4-flash-0731}), giving four independent Evaluator~A reads per unit---the single pass of Section~\ref{sec:replication_central} plus three more---and once with a fourth, architecturally unrelated model, \texttt{openai/gpt-5.6-luna-pro} (Evaluator C). We also recovered Qwen2.5-14B in two stages. First, its $87$ empty-\texttt{ground\_truth} records (Appendix~\ref{sec:fourth_model_status}) were backfilled from the authoritative catalogue, but $6$ of its $30$ events were found to have been dropped because the same stale-catalogue collision documented in Section~\ref{sec:replication_design} had, for this campaign, reached further upstream than a mislabelled micro-question: the swarm itself deliberated on the wrong topic entirely (event~T3's transcript debated a Bitcoin price threshold, not the 2024 Taiwanese presidential election its correct dossier specifies), leaving a $24$-event campaign. Second, 2026-08-21: those six events (T3, T4, T5, T6, C6, T8) were re-simulated end-to-end---300 agents, 60 rounds, all three conditions---against the correct root catalogue, and rescored under all five evaluator passes (primary \texttt{deepseek/deepseek-r1}, three Evaluator~A repeats, Evaluator~C). A first re-simulation attempt reproduced the identical stale-catalogue bug (verified: T3 again resolved to the Bitcoin question) because the launch script still resolved \texttt{--events-raw} to the wrong catalogue by a relative path; the corrected launch fixed only that path, leaving the injection bank deliberately unchanged (\S\ref{sec:fourth_model_status}). Qwen2.5-14B is now a full $30$-event, $90$-unit campaign, scored the same way as the other three models.

Because Condition~A's shared dossier makes raw cross-model Brier comparisons vulnerable to the contamination effect of Appendix~\ref{sec:cutoff}, we restrict the four-way comparison below to the $20$ events resolving after every model's public release date that survive in all four campaigns---the strictest contamination control applied to any comparison in this paper. With Qwen2.5-14B now complete, this is the full globally-clean set of Appendix~\ref{sec:cutoff} (T6, C6 and T8 are back in it; an earlier draft restricted this comparison to $17$ events because those three were still missing for Qwen2.5-14B). Per-unit Brier and directional accuracy are computed from the mean of the (up to four) Evaluator~A reads, not a single pass, following Section~\ref{sec:replication_noise}'s own diagnosis that a single read is not a reference measurement.

\begin{table*}[!htbp]
\centering
\caption{Four-model comparison on the $20$ events every campaign shares that resolve after every model's release date (no model could have memorised these outcomes). Lift and susceptibility use the per-unit mean Brier across all available Evaluator~A reads; accuracy uses the per-unit averaged probability vector. Flat rate is the fraction of individual Evaluator~A reads in Condition C returning an exactly uniform vector. Recomputed 2026-08-21 on the full $20$-event common set: Qwen2.5-14B's six wrong-topic exclusions (T3, T4, T5, T6, C6, T8) were re-simulated against the correct catalogue and rescored (\S\ref{sec:fourth_model_status}), restoring T6, C6 and T8 to the events common to all four campaigns; the set grows from $17$ to the full $20$.}
\label{tab:four_model_clean}
\small
\begin{tabular}{lccccc}
\toprule
\textbf{Model} & \textbf{Lift $A\!\to\!B$} & \textbf{Dir.\ acc., B} & \textbf{Dir.\ acc., C} & \textbf{Susceptibility} & \textbf{Flat rate, Cond.\ C} \\
\midrule
Llama-3.1-8B    & $-0.114$          & $0.684$          & $0.500$          & $+0.059$          & $3.8\%$ \\
Mistral-7B      & $-0.044$          & $0.850$          & $0.650$          & $+0.070$          & $\mathbf{2.5\%}$ \\
Qwen2.5-7B      & $+0.004$          & $0.850$          & $0.550$          & $\mathbf{+0.173}$ & $2.5\%$ \\
Qwen2.5-14B     & $\mathbf{+0.069}$ & $\mathbf{0.900}$ & $\mathbf{0.800}$ & $+0.153$          & $5.0\%$ \\
\bottomrule
\end{tabular}
\setlength{\tabcolsep}{6pt}
\end{table*}

\textbf{Contamination control changes which model looks worst, not just the aggregate sign.} Mistral-7B's lift on the full, uncontrolled $30$-event set, using the per-unit mean across all available Evaluator~A reads rather than the single pass of Section~\ref{sec:cross_model_stats}, is $-0.224$---the paper's clearest case of simulation harming a model. On the $17$ clean events it is $-0.020$, indistinguishable from the noise floor established in Section~\ref{sec:replication_noise}. Most of what looked like Mistral-specific harm from deliberation was memorisation in the shared Condition~A reference, consistent with Appendix~\ref{sec:cutoff}'s finding for the aggregate, now shown to hold at the individual-model level as well. Llama-3.1-8B's numbers in this section were recomputed 2026-08-16 after the outage repair of \S\ref{sec:outage}; the reading that motivated this paragraph in earlier drafts does not survive it. Pre-repair, the 17 clean events showed Llama's lift as positive and larger than on the full set ($+0.092$), with $25.0\%$ of its Condition~C reads returning an exactly uniform vector---roughly an order of magnitude above every other model---read as a model-specific sharpness-loss pathology once contamination stopped masking it. Post-repair, both figures collapse: lift on the clean 17 is now $-0.016$ (indistinguishable from Mistral's $-0.020$, and like it spanning zero), and the Condition~C flat rate is $7.4\%$, in the same range as Qwen2.5-7B and no longer an outlier. The $25.0\%$ figure was concentrated in exactly the events later found to carry the unreachable-endpoint corruption (\S\ref{sec:outage}); once those events carry real evaluator input, the elevated flat rate disappears along with the lift advantage. We retract the sharpness-loss-pathology reading: it was evidence about the outage's evaluator input, not about Llama-3.1-8B as an architecture.

\textbf{Qwen2.5-14B leads on lift and both directional-accuracy columns, but is no longer the sharpest on flat rate once the comparison covers its previously-missing events.} It has the largest lift ($+0.069$) and the highest directional accuracy in both With-Sim ($0.900$) and Null ($0.800$) conditions---$5$ and $15$ points, respectively, above the next-best model on each. Its Condition~C flat rate, however, is $5.0\%$---the \emph{highest} of the four models, not the lowest as it appeared under the smaller $17$-event set (where it happened to have $0.0\%$ on the events it had). T6, C6 and T8, the three events that were missing before and are now included, contribute flat reads that were simply invisible to the earlier, narrower comparison; nothing about the recovered events' evidence or scoring is unusual (\S\ref{sec:fourth_model_status}), and there is nothing to indicate they are treated differently by the evaluator. Susceptibility keeps its ordering from the $17$-event set: Qwen2.5-7B still has the largest margin ($+0.173$, was $+0.194$), ahead of Qwen2.5-14B ($+0.153$, was $+0.157$). Qwen2.5-7B remains the most consistent low-risk choice on flat rate specifically: at most $2.5\%$ in every condition across both evaluators, at the cost of the lowest directional accuracy of the four models in Condition~C.

\textbf{Cross-evaluator agreement remains moderate at this scale.} Pooled over all conditions and all available units, Evaluator~A's per-unit mean Brier correlates with Evaluator~C's single read at $r = 0.696$ (Llama, resynced 2026-08-21 alongside the Qwen2.5-14B recovery below), $0.765$ (Mistral), $0.701$ (Qwen2.5-7B) and $0.706$ (Qwen2.5-14B, now over its full $90$ units rather than $72$)---all four in the ``moderate'' band, none near the agreement a reliable instrument should show. Evaluator~C returns a flat vector on \emph{zero} of the $120$ Condition-C reads pooled across all four models ($30$ each, Qwen2.5-14B now included at its full count rather than $24$), the strongest single piece of evidence in this paper that the flat-forecast phenomenon is an artifact of the specific evaluator rather than of deliberation: it is not merely reduced by a second read, as Section~\ref{sec:replication_central} showed, but absent under a differently-implemented one. This comes with a caveat symmetric to Section~\ref{sec:epistemic_mapping}'s finding for Evaluator~J: \texttt{openai/gpt-5.6-luna-pro} rejects this pipeline's strict JSON-schema request outright (\texttt{invalid\_json\_schema}, a malformed \texttt{required} array on \texttt{validated\_scales}) and falls back to unconstrained \texttt{json\_object} mode for every call, so Evaluator~C's zero flat-rate is achieved under the same weaker enforcement level as Evaluator~A, not a stronger one---an instrument-configuration bug independent of the one in Section~\ref{sec:epistemic_mapping}, caught by the same discipline of checking rather than assuming enforcement.

The underlying instrument-noise finding of Section~\ref{sec:replication_noise} also generalises beyond Condition~A: pooling all three conditions, the mean spread across the available Evaluator~A reads per unit---four for Llama, Mistral and Qwen2.5-7B, three for Qwen2.5-14B, which has no original single-pass Evaluator~A record to add as a fourth---is $0.330$ (Llama, resynced 2026-08-21), $0.375$ (Mistral), $0.280$ (Qwen2.5-7B) and $0.208$ (Qwen2.5-14B, now over $90$ units) Brier points, with $7.8$--$13.3\%$ of units reproducing to within $0.01$ across reads. This is larger than the $0.128$-mean/$0.715$-max figure reported there for Condition~A alone, because that figure benefited from Condition~A's dossier being short and largely unambiguous relative to full deliberation transcripts; the instrument is, if anything, less reliable on the transcripts that matter most for the paper's central claim.

\section{The Evaluator Decides; Evidence Volume and Schema Enforcement Do Not Survive a Controlled Test}
\label{sec:volume_mechanism}

Section~\ref{sec:degenerate_flat} raised evidence volume as the one swarm-model-specific variable that could account for the cross-model flat-forecast pattern, and left it unestablished. The per-unit \texttt{evidence\_text} lengths measured above, joined to the flat-unit flags of each scored pass, settle it. Two questions are separable and we take them in turn: whether volume predicts degeneracy \emph{within} a campaign, and whether volume is sufficient to produce it.

\textbf{Within a campaign, longer evidence predicts a flat forecast.} Pooling the $60$ simulated units of each affected campaign, flat units carry markedly more evidence than non-flat ones: median $19{,}016$ against $3{,}009$ characters for Llama-3.1-8B (Mann--Whitney $U=626$, $p = 6.5\times10^{-5}$) and $13{,}906$ against $12{,}074$ for Mistral-7B ($U=683$, $p = 3.1\times10^{-6}$). Qwen2.5-7B, whose units are the shortest of the three campaigns, has no flat units at all and so contributes no within-model contrast.

\textbf{Llama-3.1-8B cannot carry this argument, and we do not rest it on that campaign.} Its high-volume units are exactly its six pre-outage events: all $12$ of its healthy simulated units are flat, against $6$ of the $48$ units affected by the serving failure of Appendix~\ref{sec:outage}. Evidence volume is therefore perfectly confounded there with whether any deliberation occurred at all, and the $19{,}016$-versus-$3{,}009$ contrast is equally consistent with a deliberation account and with an extraction account.

\textbf{Mistral-7B is the clean test, and it shows a monotone dose--response.} That campaign suffered no outage: all $60$ of its simulated units carry a genuine multi-speaker transcript (Table~\ref{tab:evidence_depth}, $0/60$ units with $\leq 1$ speaker), so evidence length varies across units without varying whether deliberation happened. Splitting those units into terciles by \texttt{evidence\_text} length gives flat rates of $1/20$ ($5\%$), $5/20$ ($25\%$) and $13/20$ ($65\%$) from shortest to longest, over a range of only $10{,}192$ to $15{,}557$ characters (Figure~\ref{fig:volume_mechanism}A). A $50\%$ increase in input length moves the degeneracy rate by a factor of thirteen.

\textbf{Volume is nevertheless not sufficient: the same transcripts are not flat under a different instrument.} The decisive comparison holds the input fixed and varies only the evaluator. Llama-3.1-8B's $12$ high-volume units are $12/12$ flat under Evaluator J; under the four Evaluator~A reads of Section~\ref{sec:four_model_extension} they are $0/12$, $1/12$, $2/12$ and $0/12$, and under Evaluator~C, $0/12$. Mistral-7B's $20$ longest units are $13/20$ flat under Evaluator J, then $1/20$, $1/20$, $0/20$, $0/20$ and $0/20$. Figure~\ref{fig:volume_mechanism}B shows both. Nothing about the transcripts changes across those bars; only the model asked to compress them does.

\begin{figure*}[!htbp]
\centering
\includegraphics[width=\textwidth]{figI_volume_mechanism.png}
\caption{Evidence volume as the trigger for degenerate forecasts, and the evaluator as their locus. \textbf{Panel A:} flat-forecast rate by \texttt{evidence\_text} length tercile for Mistral-7B under Evaluator J. Mistral is the clean test because its campaign suffered no serving outage, so all $60$ simulated units carry a real multi-speaker transcript and length varies without a confound. \textbf{Panel B:} the units that drive Panel A---Llama-3.1-8B's $12$ high-volume units and Mistral-7B's $20$ longest---re-scored by every available pass. The transcripts are byte-identical across all six bars in each series; only the instrument differs.}
\label{fig:volume_mechanism}
\end{figure*}

\subsection{Does the failure transfer to other unconstrained evaluators?}
\label{sec:transfer_expt}

\textbf{And enforcement is not the whole of the instrument: the failure is specific to small evaluators, and reproduces on one of them.} The comparison above is suggestive but incomplete, because every pass it compares against---the four Evaluator~A reads and Evaluator~C---has structured-output enforcement \emph{on}. It therefore cannot separate two explanations: that enforcement is what prevents the collapse, or that this particular model is what causes it. We ran the missing arm. All $60$ of Mistral-7B's simulated units were re-scored by five further evaluators, each called with no \texttt{response\_format} field at all---the same unconstrained condition the July instrument ran under---on byte-identical transcripts. If unconstrained decoding over a long transcript were sufficient, all five should degenerate; if the July instrument were unique, none should.

Neither is what happens. Four of the five return \emph{no} flat units: \texttt{deepseek/deepseek-r1} $0/56$ (mean modal $0.667$), \texttt{deepseek/deepseek-r1-distill-llama-70b} $0/54$ ($0.675$), \texttt{qwen/qwen3-32b} $0/57$ ($0.577$) and \texttt{mistralai/ministral-14b-2512} $0/60$ ($0.557$). The fifth does: \texttt{qwen/qwen3-14b} returns an exactly uniform vector for $8/58$ units ($13.8\%$), against Evaluator~J's $19/60$ ($31.7\%$) on the same material. Pooled, the comparison evaluators give $8$ flat units in $285$ scored observations. Two controls matter here. First, prompt simplicity is not the explanation: running \texttt{deepseek/deepseek-r1} through the production evaluator class with the full rubric prompt verbatim and enforcement still disabled gives $0/60$, including $0/19$ on the units Evaluator~J flattened. Second, parameter count alone is not the explanation either, since \texttt{ministral-14b-2512} is the same size as the model that fails and does not fail. The two instruments that degenerate---DeepSeek-R1-14B and Qwen3-14B---are both \emph{small reasoning} models, and the within-family contrast is clean: Qwen3-14B gives $8/58$ where Qwen3-32B, same family and same prompt, gives $0/57$. With one model per cell we offer that as the pattern the data shows, not as an established mechanism.

\textbf{The two failing instruments fail on the same transcripts.} This is the part that turns a curiosity into evidence about the input. Of Qwen3-14B's $8$ flat units, $6$ fall inside the $19$-unit set Evaluator~J flattened, where chance would put $2.6$ (one-sided Fisher exact $p = 0.012$). Its flat units also carry more evidence than its sharp ones---median $13{,}817$ against $12{,}760$ characters---the same direction and nearly the same magnitude as Evaluator~J's own split ($13{,}906$ against $12{,}074$). Two independently trained models, from different families, agree above chance about which transcripts cannot be reduced to an informative distribution, and agree that the harder ones are the longer ones. That is difficult to explain as an idiosyncrasy of one deployment and easy to explain as a property of the input interacting with limited extraction capacity.

\textbf{What this narrows.} An earlier reading of this paper held that ``evidence volume is the trigger and the schema-disabled evaluator is the locus,'' which reads as though any schema-disabled evaluator would do. It would not: three of the five comparison evaluators are at least as capable as Evaluator~J, run in the same condition, and return nothing uniform at all. What separates the instruments that fail from those that do not is the model, and \S\ref{sec:truncation_grid} then shows by direct assignment that neither of the two factors we had credited---input length and enforcement---accounts for the difference. We could not re-run Evaluator~J itself, because the local endpoint no longer exists and no provider we have access to serves DeepSeek-R1-14B; that limits every causal statement in this section to a proxy. Figure~\ref{fig:transfer} shows the full comparison.

\begin{figure*}[!htbp]
\centering
\includegraphics[width=\textwidth]{figK_degeneracy_transfer.png}
\caption{Degenerate consensus across unconstrained evaluators. Every row scores all $60$ of Mistral-7B's simulated units---byte-identical transcripts---with no \texttt{response\_format} field, the identical condition the July instrument ran under. \textbf{Panel A:} flat-forecast rate with $95\%$ Wilson intervals. Four of the five comparison evaluators return nothing uniform; \texttt{qwen/qwen3-14b}, the smallest reasoning model tested, returns $8/58$. The \emph{pipeline rubric prompt} row uses the production evaluator class and its full prompt verbatim, ruling out prompt simplicity as the explanation. \textbf{Panel B:} the same rates restricted to the $19$ units Evaluator~J itself flattened. Qwen3-14B's failures concentrate there---$6$ of its $8$, against $2.6$ expected by chance---so the two failing instruments are failing on the same inputs rather than independently.}
\label{fig:transfer}
\end{figure*}

\subsection{The truncation grid, and what it does to this section's argument}
\label{sec:truncation_grid}

Everything above is observational in the variable that matters. The tercile result compares units that \emph{happen} to carry more evidence with units that happen to carry less, and the instrument contrast compares evaluators that differ in more than enforcement. We therefore ran the manipulation directly: a $2 \times 4$ grid crossing structured-output enforcement against four assigned evidence lengths, over all $60$ of Mistral-7B's simulated units, with the evaluator model, the prompt, the temperature and the transcript identical within each arm. Truncation retains both sections of the record at every level, keeping the log's final lines and uniformly subsampling the action lines, so the levels differ in how much of each section is shown rather than in which sections exist.

The evaluator is \texttt{qwen/qwen3-14b}, chosen because it is the only accessible model that reproduces the degeneracy at all (\S\ref{sec:transfer_expt}); \texttt{deepseek/deepseek-r1} is included as a negative control at three of the four levels. Evaluator J itself cannot be used: its local endpoint no longer exists and no provider we can reach serves DeepSeek-R1-14B. That is the central limitation of what follows, and it means the grid tests this section's mechanism on a proxy rather than on the instrument that produced the finding.

\textbf{Assigned evidence length does not increase degeneracy, and the sign is reversed.} Under \texttt{qwen3-14b} with enforcement off, the flat rate across the four levels is $7/59$ at $\approx\!1.8$k characters, $4/60$ at $\approx\!4.8$k, $0/60$ at $\approx\!7.8$k and $1/60$ untruncated. The \emph{shortest} condition is the worst one, and the contrast between the extremes is significant in the direction opposite to the hypothesis (two-sided Fisher exact $p = 0.032$). Excess sharpness, measured as the modal probability minus $1/K$ so that it is comparable across the mixed-$K$ catalogue, is likewise flat in length: $+0.171$, $+0.190$, $+0.203$, $+0.185$, with overlapping intervals throughout. The negative control behaves the same way and at a different level entirely: \texttt{deepseek-r1} gives $1/60$, $1/49$ and $0/60$ flat with excess sharpness near $+0.32$ at every length.

\textbf{Enforcement, isolated within one model, does nothing.} Pooled over all four lengths, \texttt{qwen3-14b} returns $12/239$ flat units with enforcement off and $11/239$ with it on (Fisher $p = 1.000$), at excess sharpness $+0.187$ and $+0.208$. This is the first test in this paper that varies enforcement \emph{alone}. Every earlier comparison---including the re-scoring of Section~\ref{sec:replication_central}---changed the evaluator model and the enforcement setting together, so none of them could attribute the effect to either.

\begin{figure*}[!htbp]
\centering
\includegraphics[width=\textwidth]{figL_truncation.png}
\caption{The $2 \times 4$ truncation grid over Mistral-7B's $60$ simulated units. Within an arm the evaluator, prompt, temperature and transcript are fixed; only how much of the transcript is shown varies. \textbf{Panel A:} flat-forecast rate with $95\%$ Wilson intervals. Assigning a shorter transcript does not reduce degeneracy---the shortest condition is the worst, and the $\approx\!2$k-versus-untruncated contrast is significant in the reverse direction ($p = 0.032$). \textbf{Panel B:} excess sharpness, the modal probability minus $1/K$, which is zero for a uniform vector at any $K$ and therefore comparable across this mixed-$K$ catalogue. It is flat in length for both evaluators, and the gap between them ($\approx\!+0.19$ for \texttt{qwen3-14b} against $\approx\!+0.32$ for \texttt{deepseek-r1}) is larger at every level than anything length or enforcement does.}
\label{fig:truncation}
\end{figure*}

\textbf{This is a negative result for two of this section's three claims, and we report it as one.} The dose--response of Figure~\ref{fig:volume_mechanism}A is real as an association within Evaluator J and does not reproduce as an effect when length is assigned. Two readings remain open and the data do not separate them: either the volume relationship is specific to Evaluator J, which we cannot re-run, or the tercile pattern was always driven by something that co-varies with length across units rather than by length itself---which is precisely the alternative this experiment was built to exclude, and it does not exclude it. The schema-enforcement claim fares no better under isolation. What survives, and is strengthened rather than weakened, is the claim this paper's title makes: \emph{which model reads the transcript} dominates. The gap between \texttt{deepseek-r1} and \texttt{qwen3-14b} in excess sharpness is roughly $0.13$ at every length, larger than anything either manipulation produces, and it is the same axis along which Evaluator~J differs from every evaluator that does not degenerate.

We leave the recommendation of Section~\ref{sec:replication_implications} in place but re-derive it. Verifying enforcement at runtime remains sound engineering---a pipeline should know what it sent---but on this evidence we cannot promise that enabling it will prevent degenerate extraction, and a benchmark that switches it on and declares the problem solved would be relying on something we tested and could not confirm. The check that does discriminate is the one we already recommend for a different reason: score with a second, stronger evaluator and compare.

\textbf{What this establishes.} Taking the observational and the experimental evidence together: degeneracy is a property of the extracting model, associated with input volume within the one instrument we could not re-run and not caused by it in the two we could. Long transcripts do not intrinsically elicit uniform forecasts---four independent reads and a second architecture recover sharp distributions from the very same text---and the instrument does not degenerate indiscriminately either, since its own flat rate on the same campaigns' short inputs is low. We had attributed the vulnerability to the schema-disabling defect of Section~\ref{sec:epistemic_mapping} and its exercise to input length. Section~\ref{sec:truncation_grid} tests both in isolation and supports neither, so we withdraw the mechanism while keeping the localisation: the defect is real and worth fixing, but it is not what we can show makes the difference. This resolves the question Section~\ref{sec:degenerate_flat} left open in favour of the extraction account, and it does so without appeal to the cross-model comparison whose confounds Section~\ref{sec:serving_note} documents: the dose--response is measured within a single campaign, and the instrument contrast is measured on fixed inputs.

We previously noted an asymmetry here that the 2026-08-16 repair has since resolved rather than merely updated the numbers on. Before the repair, the residual flat units under Evaluator~A relocated to Llama's low-volume arm ($10$, $8$, $12$ and $12$ of $48$ across the four reads, against $0$--$2$ of $12$ high-volume): those $48$ were exactly the outage units, whose \texttt{evidence\_text} was a connection-error dump containing no information about the question (Appendix~\ref{sec:evidence_depth}), for which a uniform posterior is the defensible response. Post-repair, with real transcripts in place of the connection-error dump, the same $48$-unit arm's flat count falls to $3$, $2$, $2$ and $2$ of $48$---in the same low range as the high-volume arm, not elevated. This is consistent with, rather than contradicting, the extraction-channel account: the low-volume arm was never flat \emph{because} it was low-volume; it was flat because it was empty. Once it is no longer empty, it behaves like any other short-but-substantive input. The asymmetry was therefore itself a symptom of the outage, layered on top of the genuine volume/instrument interaction Mistral's clean test establishes; it is retracted as a separate finding.

\textbf{The remaining ambiguity is narrow, and we state it.} A dose--response within Mistral establishes that the amount of transcript shown drives Evaluator J's degeneracy; it does not establish \emph{why}, and it cannot separate raw character count from the properties that co-vary with it in this corpus---more speakers, more distinct claims, more conflicting positions to reconcile. The truncation grid we ran does not resolve this either, and we are explicit about why: showing less of a transcript necessarily shows fewer of its voices. Across our four levels the median number of distinct speakers visible in the representative-actions block falls from $44$ untruncated to $29$, $20$ and $9$, tracking the character count almost exactly. The manipulation therefore varies \emph{amount of deliberation shown} as a bundle, which is the same bundle the natural across-unit variation varies; what it buys is that the variation is now assigned rather than observed, so it cannot be explained by long transcripts differing in kind from short ones. Isolating character count alone would require padding with semantically null filler at a fixed speaker count---a different experiment, and one we have not run.

\parttitle{II}{What Survives the Audit}

\section{What the Benchmark Does Establish: Content-Sensitive Predictive Signal}
\label{sec:predictive_signal}

The preceding sections withdraw a great deal. It is worth stating precisely what survives, because one substantive claim does, and it is the claim the benchmark was built to test in its most defensible form. We evaluate it on the same footing as Table~\ref{tab:four_model_clean}: the $20$ contamination-clean events common to all four campaigns, scored from the per-unit mean of the available Evaluator~A reads rather than a single pass. Table~\ref{tab:predictive_signal} and Figure~\ref{fig:predictive_signal} report it.

\begin{table*}[!htbp]
\centering
\caption{Predictive signal on the $20$ contamination-clean events common to all four campaigns, from the per-unit mean of the available Evaluator~A reads. Chance directional accuracy is $1/K$ per event, averaged; $K$ ranges from $2$ to $9$ across these events, so the null distribution is a Poisson binomial and the reported $p$ is exact and one-sided. $\Delta$ columns are Condition B minus Condition C, with $95\%$ percentile bootstrap intervals ($B=20{,}000$); positive means relevant evidence beats topically unrelated evidence. Recomputed 2026-08-21 on the full $20$-event common set (was $17$; \S\ref{sec:fourth_model_status}).}
\label{tab:predictive_signal}
\footnotesize
\setlength{\tabcolsep}{4pt}
\resizebox{\textwidth}{!}{%
\begin{tabular}{lcccccc}
\toprule
\textbf{Model} & \textbf{Dir.\ acc., B} & \textbf{Dir.\ acc., C} & \textbf{Chance} & \textbf{Exact $p$ (B vs.\ chance)} & \textbf{$\Delta$ dir.\ acc.\ [95\% CI]} & \textbf{$\Delta$ Brier [95\% CI]} \\
\midrule
Llama-3.1-8B\textsuperscript{$\dagger$}  & $0.675$ & $0.500$ & $0.326$ & $3.5\times10^{-4}$  & $+0.175$ $[-0.025, +0.400]$ & $+0.069$ $[-0.058, +0.208]$ \\
Mistral-7B    & $0.850$ & $0.650$ & $0.326$ & $5.3\times10^{-7}$  & $+0.200$ $[-0.050, +0.450]$ & $+0.064$ $[-0.142, +0.261]$ \\
Qwen2.5-7B    & $0.850$ & $0.550$ & $0.326$ & $5.3\times10^{-7}$  & $+0.300$ $[+0.050, +0.550]$ & $+0.173$ $[+0.073, +0.271]$ \\
Qwen2.5-14B   & $\mathbf{0.900}$ & $0.800$ & $0.326$ & $3.1\times10^{-8}$ & $+0.100$ $[-0.100, +0.300]$ & $+0.143$ $[+0.060, +0.234]$ \\
\midrule
\textbf{Pooled} ($80$ units) & $\mathbf{0.819}$ & $0.625$ & $0.326$ & $\mathbf{3.7\times10^{-22}}$ & $\mathbf{+0.194}$ $[+0.075, +0.313]$ & $\mathbf{+0.112}$ $[+0.042, +0.182]$ \\
\bottomrule
\end{tabular}%
}
\smallskip\\
\footnotesize\textsuperscript{$\dagger$}Llama-3.1-8B's row is computed from the re-simulated transcripts of \S\ref{sec:outage}. Its per-model contribution to the surviving finding is the weakest of the four: both of its $\Delta$ intervals now span zero, and its Condition~B directional accuracy ($0.675$) is the lowest here, though still far above the $0.326$ chance level. The pooled result---the one this section's claim actually rests on---excludes zero on both $\Delta$ columns and is driven by all four campaigns rather than by this one. Directional accuracy credits a $k$-way tie at $1/k$; across the $80$ Condition~B units there are $65$ outright hits and a single two-way tie, giving $65.5/80 = 0.819$.
\setlength{\tabcolsep}{6pt}
\end{table*}

\begin{figure*}[!htbp]
\centering
\includegraphics[width=\textwidth]{figJ_predictive_signal.png}
\caption{The benchmark's one surviving positive result, on the $20$ contamination-clean events common to all four campaigns. \textbf{Panel A:} Condition B directional accuracy against the per-event chance level $1/K$ (mean $32.6\%$; $K$ ranges $2$--$9$), with Condition C shown for contrast. \textbf{Panel B:} the Brier gained from relevant rather than topically unrelated evidence, $BS_C - BS_B$, with $95\%$ percentile bootstrap intervals. The pooled interval and two of four per-model intervals exclude zero.}
\label{fig:predictive_signal}
\end{figure*}

\textbf{The simulated forecasts identify realised outcomes far above chance.} Pooling the four campaigns, the swarm's modal option is the outcome that actually occurred in $65$ of $80$ (model, event) pairs outright, with one further unit a two-way tie that includes the outcome; credited at $1/k$ that is $65.5/80 = 81.9\%$ (recomputed 2026-08-21 after the Qwen2.5-14B recovery of \S\ref{sec:fourth_model_status} grew the common clean-event set from $17$ to $20$; was $56.5/68$, $83.1\%$), against a chance level of $32.6\%$. Because the option count varies across these events, the null is a Poisson binomial rather than a coin flip; evaluated exactly, $p = 3.7\times10^{-22}$ pooled and below $4\times10^{-4}$ for every individual model (Llama-3.1-8B is still the weakest individual contributor at $p = 3.5\times10^{-4}$, still far past any conventional threshold). Ten of the twenty events are called correctly by all four models and \emph{no} event is called wrongly by all four. These are events that resolved after every model's public release date, so no campaign could have memorised the outcome, and they are scored from four independent evaluator reads rather than the single pass whose instability Section~\ref{sec:replication_noise} documents.

\textbf{And the swarm is responding to the content of its evidence, not merely to its presence.} Directional accuracy alone cannot distinguish a swarm that reasons from one that inherits a good prior from its base model. Condition~C is the control that can: it runs the identical $300$-agent, $60$-round machinery, injects a structurally matched payload at the same checkpoint, and differs only in that the payload is topically unrelated to the question. Replacing relevant evidence with irrelevant evidence costs $0.112$ Brier $[+0.042, +0.182]$ and $19.4$ points of directional accuracy $[+7.5, +31.3]$ pooled (recomputed 2026-08-21 once the Qwen2.5-14B recovery of \S\ref{sec:fourth_model_status} grew the common clean-event set from $17$ to $20$; was $0.125$ $[+0.053,+0.195]$ Brier and $18.4$ $[+6.6,+30.9]$ points on the 17-event set)---both intervals excluding zero, though only two of the four per-model Brier intervals exclude zero individually, Llama's being $[-0.058,+0.208]$ and Mistral's $[-0.142,+0.261]$. This contrast is a within-simulation comparison: it never touches the shared Condition~A record of Section~\ref{sec:condition_a}, so it is immune to the confound that disqualifies the aggregate lift, and Appendix~\ref{sec:cutoff} shows it survives knowledge-cutoff stratification unchanged.

\textbf{What this does not establish, and we are explicit about it.} Three limits bound the claim.

First, \emph{the simulation is not shown to improve on direct elicitation, and on the recovered 20-event set it now looks worse for one model}. Condition~B's directional accuracy (B vs.\ A) is $0.675$ vs.\ $0.650$ for Llama-3.1-8B (B ahead), $0.850$ vs.\ $0.900$ for Mistral-7B (B \emph{behind}, the first case in this comparison where simulation trails direct elicitation), $0.850$ vs.\ $0.850$ for Qwen2.5-7B (tied), and $0.900$ vs.\ $0.850$ for Qwen2.5-14B (B ahead); recomputed 2026-08-21 on the 20-event set (was three ties and one B-ahead case on the 17-event set). No per-model difference approaches significance on $20$ events either way. Much of the signal in Panel A is therefore attributable to what the base model already knows, not to deliberation. The B-versus-C contrast isolates what deliberation demonstrably contributes---sensitivity to which evidence it is given---and that is a weaker claim than prediction lift. The original Prediction Lift hypothesis (Claim~1) remains unanswerable through Condition~A for the reason given in Section~\ref{sec:condition_a}: it was run once and shared, so no within-model contrast exists there. Section~\ref{sec:cot_baseline} supplies the missing per-model comparison from the other direction, by running a single-agent baseline against the same evidence; it does not overturn this limit so much as sharpen it, since a single call matches the swarm's directional accuracy outright.

Second, \emph{directional accuracy is a coarse criterion}. It scores only which outcome carries the most mass. Condition~B's mean Brier on these events is $0.22$--$0.34$ depending on model---informative, but far from the sharp, well-calibrated distribution a forecasting system would need to be useful operationally. A system can be $88\%$ directionally right and still poorly calibrated, and Section~\ref{sec:degenerate_sharpness} shows this one hedges systematically.

Third, \emph{release-date stratification is a necessary control, not a sufficient one}. It rules out memorisation of the outcome itself. It does not rule out the seed dossier containing strong contemporaneous signal about the likely outcome, which the swarm would then be relaying rather than inferring. Distinguishing those requires holding dossier content fixed while varying deliberation, which this design does not do.

\subsection{A Single-Agent Baseline: One Call Is Not Beaten by Three Hundred Agents}
\label{sec:cot_baseline}

The first of those three limits is the one a reader will press hardest on, and until now we could only concede it, because Condition~A was shared rather than run per model (Section~\ref{sec:condition_a}). This subsection removes that excuse from the other side: instead of re-running the No-Sim condition, we run a \emph{single-agent} control that receives exactly the evidence the swarm's agents received.

\textbf{Design.} For each of the $20$ contamination-clean events and each of four models we issue one retrieval-augmented chain-of-thought call: system prompt casting the model as a calibrated forecaster, the event question and its $K$ options, the same seed dossier the swarm's agents were initialised from, and the same round-$30$ relevant-update payload Condition~B injected---then a demand for step-by-step reasoning and a JSON probability vector over the options, with structured-output enforcement on. Three replicates per cell (replicate~1 at temperature $0$, seed $42$; replicates~2 and~3 at temperature $0.7$, seeds $100$ and $2026$) give $20 \times 4 \times 3 = 240$ calls (recomputed 2026-08-21 on $3$ additional events once Qwen2.5-14B's recovery grew the common clean-event set from $17$ to $20$; \S\ref{sec:fourth_model_status}), all $240$ of which returned parseable output. The comparison is to Condition~B scored exactly as everywhere else in this section, from the per-unit mean of the available Evaluator~A reads.

\textbf{Three ways this comparison can be got wrong, and how we avoid each.} We state them because our own first pass at this analysis made all three mistakes simultaneously and reported a spurious $p = 0.045$ against the swarm.

\emph{(i) Arm matching.} Only two of the four single-agent runs use the same checkpoint as the campaign they are set against. \texttt{mistralai/ministral-8b-2512} is neither the model nor the generation of Mistral-7B-Instruct-v0.1, and \texttt{qwen/qwen3-14b} is a generation ahead of Qwen2.5-14B; those two arms compare different models and cannot be read as a within-model contrast. The primary estimate below is therefore the two identity-matched arms only, and the four-arm pool is reported separately and marked. Even the matched arms differ in precision and serving stack---FP16 through a hosted API against local AWQ-INT4 under vLLM and GGUF under Ollama respectively---which is the same confound Section~\ref{sec:serving_note} documents, here working in the single agent's favour.

\emph{(ii) Aggregation symmetry.} Collapsing the three replicates into one averaged probability vector gives the single-agent arm an ensembling advantage that the swarm arm does not receive: the swarm's four reads re-read \emph{one} simulation, they do not re-run it. Because the headline runs against the swarm, that asymmetry pushes in the same direction as the conclusion, so we report the mean of the per-replicate scores---the expected cost of one call---as primary, and the three-call ensemble only as an upper bound on what the single-agent arm can do at triple cost.

\emph{(iii) Clustering.} The same $20$ events recur across models, so pooled unit-pairs are not independent and a Diebold--Mariano test at $\mathrm{df}=n-1$ over them is anticonservative. Pooled intervals here come from a cluster bootstrap that resamples \emph{events}, and pooled $p$-values from a cluster-robust $t$ on the per-event means ($\mathrm{df}=19$). The difference is not cosmetic: on the four-arm pool the naive test returns $p = 0.016$ where the clustered one returns $p = 0.083$.

\begin{table*}[!htbp]
\centering
\caption{Single-agent RAG+CoT baseline against Condition~B, on the $20$ contamination-clean events (recomputed 2026-08-21; was $17$, \S\ref{sec:fourth_model_status}). $\Delta$ Brier is $BS_{\mathrm{CoT}} - BS_{\mathrm{swarm}}$, so \emph{positive favours the swarm}; intervals are $95\%$ percentile bootstrap ($B=20{,}000$), per-arm over events and pooled over event clusters. Per-arm $p$ is a Diebold--Mariano test with the Harvey--Leybourne--Newbold small-sample correction ($\mathrm{df}=19$); pooled $p$ is the cluster-robust equivalent. The CoT column is the mean of three per-replicate scores, i.e.\ the expected performance of \emph{one} call.}
\label{tab:cot_baseline}
\footnotesize
\setlength{\tabcolsep}{4pt}
\resizebox{\textwidth}{!}{%
\begin{tabular}{lcccccc}
\toprule
 & \multicolumn{2}{c}{\textbf{Dir.\ accuracy}} & \multicolumn{2}{c}{\textbf{Mean Brier}} & & \\
\cmidrule(lr){2-3}\cmidrule(lr){4-5}
\textbf{Arm (CoT model vs.\ campaign)} & \textbf{swarm} & \textbf{1 CoT call} & \textbf{swarm} & \textbf{1 CoT call} & \textbf{$\Delta$ Brier [95\% CI]} & \textbf{$p$} \\
\midrule
\multicolumn{7}{l}{\emph{Identity-matched arms}} \\
Llama-3.1-8B vs.\ Llama-3.1-8B      & $0.675$ & $\mathbf{0.778}$ & $0.545$ & $\mathbf{0.364}$ & $-0.181$ $[-0.366, -0.020]$ & $0.065$ \\
Qwen2.5-7B vs.\ Qwen2.5-7B          & $\mathbf{0.850}$ & $0.790$ & $0.348$ & $\mathbf{0.318}$ & $-0.030$ $[-0.134, +0.081]$ & $0.613$ \\
\textbf{Pooled, matched} ($40$ pairs) & $0.763$ & $\mathbf{0.784}$ & $0.446$ & $\mathbf{0.341}$ & $\mathbf{-0.106}$ $[-0.226, +0.009]$ & $0.109$ \\
\midrule
\multicolumn{7}{l}{\emph{Cross-model arms\textsuperscript{$\ast$} --- not a within-model contrast}} \\
Ministral-8B vs.\ Mistral-7B\textsuperscript{$\ast$}   & $0.850$ & $\mathbf{0.912}$ & $0.378$ & $\mathbf{0.208}$ & $-0.170$ $[-0.369, +0.012]$ & $0.115$ \\
Qwen3-14B vs.\ Qwen2.5-14B\textsuperscript{$\ast$}     & $\mathbf{0.900}$ & $0.858$ & $\mathbf{0.289}$ & $0.280$ & $-0.008$ $[-0.116, +0.093]$ & $0.881$ \\
\textbf{Pooled, all four}\textsuperscript{$\ast$} ($80$ pairs) & $0.819$ & $\mathbf{0.834}$ & $0.390$ & $\mathbf{0.292}$ & $-0.097$ $[-0.199, -0.000]$ & $0.083$ \\
\bottomrule
\end{tabular}%
}
\smallskip\\
\footnotesize\textsuperscript{$\ast$}The CoT model is a different checkpoint from the campaign it is compared against; these rows are reported for completeness and are not evidence about whether deliberation helps a given model. Chance directional accuracy is $32.6\%$ on this event set. Both arms are far above it: pooled over the matched arms the swarm is at $0.763$ (exact Poisson-binomial $p = 1.0\times10^{-9}$) and a single temperature-$0$ call at $0.768$ ($p = 7.6\times10^{-9}$).
\setlength{\tabcolsep}{6pt}
\end{table*}

\begin{figure*}[!htbp]
\centering
\includegraphics[width=\textwidth]{figM_cot_vs_swarm.png}
\caption{The single-agent baseline. \textbf{Panel A:} Brier score for each arm, single call against swarm, on a common axis; the length of the connector is the quantity of interest. The two identity-matched arms are separated from the two cross-model ones. \textbf{Panel B:} the paired difference $BS_{\mathrm{CoT}} - BS_{\mathrm{swarm}}$ with $95\%$ intervals (per-arm bootstrap over events; pooled cluster bootstrap over event clusters). Positive favours the swarm. Open diamonds mark where the point estimate moves if the three CoT replicates are ensembled into one vector instead of scored separately---a $3\times$-cost variant shown to make the aggregation choice visible, not as the primary estimate.}
\label{fig:cot_baseline}
\end{figure*}

\textbf{The result is a non-superiority finding pointing more clearly toward the single call, and we are careful not to overstate it in either direction.} On both identity-matched arms the single call has the lower mean Brier---$0.364$ against $0.545$ for Llama-3.1-8B and $0.318$ against $0.348$ for Qwen2.5-7B---and the higher directional accuracy. Pooled over the matched arms the gap is $0.106$ Brier in the single agent's favour, but the interval is $[-0.226, +0.009]$ and the clustered $p$ is $0.109$: at $20$ events this design still cannot separate that gap from zero, exactly as \S\ref{sec:power_limitation} would predict, though it comes closer than the $17$-event reading did ($p=0.156$). The per-event differences give $d_z = 0.39$, so the achieved power against an effect of the size observed is $37\%$, and $55$ events would be needed for $80\%$ (was $30\%$ and $59$ events on the smaller set). A null here is what an underpowered design returns when the true effect is anywhere from moderately negative to slightly positive; it is not a demonstration of equivalence. What the data do support is the negative: \emph{there is no arm, matched or not, on which the swarm's advantage over a single call is distinguishable from zero}, and on the recovered 20-event set the point estimate now favours the single call on \emph{all four} arms, not three of four as on the 17-event set---Qwen3-14B, previously the one arm where the swarm led on points, now shows a near-zero gap in the single call's favour ($-0.008$ $[-0.116,+0.093]$) rather than a swarm advantage, though it remains a cross-model arm in which neither direction is evidence about Qwen2.5-14B specifically. One further nuance the 20-event set surfaces: two of the six intervals plotted in Figure~\ref{fig:cot_baseline}B --- the Llama-3.1-8B matched arm alone ($[-0.365, -0.020]$) and the all-four pool ($[-0.199, -0.000]$) --- now exclude zero in the single call's favour under their own (non-clustered, per-arm) bootstrap, even though neither is the paper's primary inferential test. The primary test is the cluster-robust one reported above, which does not reach significance on either scope; we report the two individually-significant intervals for transparency, not as an additional finding, since a reader comparing this figure against Table~\ref{tab:cot_baseline} would otherwise wonder why two segments do not touch the dashed zero line.

\textbf{The estimate is stable across the analysis choices, and two of six configurations now exclude zero.} Substituting the three-call ensemble for the single call moves the matched pool from $-0.106$ $[-0.226, +0.009]$ to $-0.135$ $[-0.247, -0.028]$, whose interval excludes zero with clustered $p = 0.034$; taking replicate~1 alone gives $-0.114$ $[-0.244, +0.007]$. The all-four-arm pool shows the same pattern: single $-0.097$ $[-0.199, -0.000]$, ensemble $-0.118$ $[-0.215, -0.025]$ (excludes zero, clustered $p = 0.032$), replicate-1-only $-0.095$ $[-0.202, +0.013]$. Across scopes and variants the point estimate stays between $-0.135$ and $-0.095$, and the lowest clustered $p$ is $0.032$ (recomputed 2026-08-21 on the 20-event set; was $-0.150$ to $-0.052$ and $0.069$ on the 17-event set, where no configuration excluded zero). Only the two three-call-ensemble configurations---a $3\times$-cost variant we report as an upper bound, not the primary estimate---exclude zero; the primary single-call estimate still does not, on either scope. Our discarded first pass reached $p = 0.045$ by combining the four-arm pool, the ensembled CoT vector, and the unclustered $\mathrm{df}=67$ test on the original 17-event data---each of which moves the number the same way. We record this because it is the same class of error the rest of this paper documents in the benchmark itself: a defensible-looking pipeline whose defaults all lean one direction.

\textbf{What this does and does not settle.} It does not replace the per-model Condition~A re-run, which earlier drafts called the single most valuable next experiment and which we now describe as one of two. The two ask different questions: Condition~A as specified uses a different elicitation prompt and a different serving stack from the campaign, whereas this baseline holds the evidence fixed and varies only whether $300$ agents deliberate over it for $60$ rounds. What it does settle is the cost question that motivated the benchmark. A $300$-agent, $60$-round simulation, its transcript, and an LLM evaluator to compress that transcript into a distribution buy no measurable accuracy over one call to the same model with the same evidence---and there is a real possibility, not established at this sample size, that they cost accuracy. It also qualifies the finding this section is named for: the above-chance directional accuracy of Table~\ref{tab:predictive_signal} is \emph{reproduced, and on this reading slightly exceeded, by a single call} ($0.834$ pooled over all four arms against the swarm's $0.819$), so that half of the surviving result is evidence about what these base models know, not about what deliberation adds. The B-versus-C contrast is untouched by this, being a comparison between two simulated conditions with no single-agent analogue in this run, and it remains the part of the finding that is specifically about the simulation.

\textbf{The defensible summary.} On events no model could have memorised, this multi-agent simulation produces forecasts that are strongly informative about how real social events resolved, and that measurably track the content of the evidence deliberated over rather than merely its presence---but the informativeness is matched by a single chain-of-thought call to the same model on the same evidence (\S\ref{sec:cot_baseline}), so only the content-sensitivity contrast is evidence about deliberation as such. Whether deliberation improves on simply asking the model is not established by this campaign; the single-agent baseline says the improvement, if any, is smaller than this design can resolve. We note the ordering this analysis reverses: under the single-pass Evaluator~J reading of Section~\ref{sec:cross_model_stats}, Mistral-7B was the only model whose susceptibility interval excluded zero; under four-read means on contamination-clean events it is the only model whose interval does not. That reversal is a further illustration of Section~\ref{sec:replication_noise}'s point, and a reason to trust the pooled figure over any single-model row.


\section{Discussion and Limitations}
\label{sec:discussion}

\subsection{On the Aggregate Simulation Lift Reversal, and Why We Read It as a Sharpness Story}
\label{sec:lift_reversal}

The result that launched this investigation is that the mean aggregate Brier score for With-Sim (Condition B) \emph{exceeds} the No-Sim baseline (Condition A) for all three models ($BS_B > BS_A$: $0.408 > 0.341$ for Llama-8B, recomputed 2026-08-16 post-repair, was $0.484 > 0.385$; $0.531 > 0.385$ for Mistral-7B; $0.481 > 0.385$ for Qwen-7B). Read as a benchmark result, this looks like a reversal of the ``Wisdom of Crowds'' hypothesis. We now believe that reading is largely wrong, and that the honest ordering of candidate causes is as follows.

\textbf{First, and primarily, measurement of forecast sharpness.} Section~\ref{sec:degenerate}---the phenomenon Part~I opens with, and which Sections~\ref{sec:epistemic_mapping}--\ref{sec:volume_mechanism} subsequently relocate to the extraction instrument---shows that $20.6\%$ of simulated units return an exactly uniform distribution---an extraction non-answer scored as a confident error---against $3.3\%$ under Condition A, and that simulation reduces mean modal probability from $0.658$ to roughly $0.51$ even among non-degenerate units. A strictly proper scoring rule penalises both~\cite{gneiting2007}. Neither requires the simulation to reason worse; both are sufficient by themselves to produce the observed sign, and the fact that directional accuracy is preserved for two of the three models (Section~\ref{sec:degenerate_sharpness}) is evidence that they did.

\textbf{Second, the baseline is not the model's own.} Because Condition A is a shared record (Section~\ref{sec:condition_a}), $BS_A - BS_B$ compares each swarm against a fixed external forecast rather than against its own direct elicitation. The quantity the reversal describes is therefore not the quantity the original Prediction Lift hypothesis asks about.

\textbf{Third, and only then, epistemic explanations.} The bimodality in per-event lift is real: a subset of events (T2, C11, S9, S3, C3) produces positive lift across all three models, while another (T1, C10, C15, S6) is uniformly negative. But we can no longer claim to know why, because the correlation previously used to characterise the split is an artifact (Appendix~\ref{sec:tautology}) and the split itself is defined by the variable it is meant to explain. We retain the SBE/SDE distinction as a description of the data and withdraw the causal reading.

Read together, these three points reframe what this campaign is evidence \emph{for}. It is weak evidence about whether social simulation improves forecasting: the measurement is too confounded and too underpowered to say. It is comparatively strong evidence about the \emph{measurement channel} itself---that asking an unconstrained LLM to compress a long deliberation transcript into a categorical distribution erases forecast sharpness before it erases forecast direction, and that this, not a deficit in collective reasoning, is the defensible explanation for the negative headline number.

Part~I settles which of the two candidate mechanisms is responsible, and the answer is not the one this subsection's earlier drafts proposed. The question posed here---whether the sharpness loss tracks deliberation or the swarm model's own vLLM serving path---turned out to be malformed: the \texttt{probabilities} field is produced by a single fixed evaluator irrespective of which swarm model is under test (Section~\ref{sec:epistemic_mapping}), so the serving stack of the model being benchmarked cannot be the operative variable. Holding the transcripts fixed and varying only the evaluator recovers sharp forecasts (Section~\ref{sec:evaluator_replication}), and a within-campaign dose--response identifies input volume as what exercises the defect (Section~\ref{sec:volume_mechanism}). The remaining agenda is correspondingly narrower: restore enforcement in place and re-measure with the \emph{same} evaluator, re-run Condition~A per model, and characterise event structure using covariates measured independently of either condition---resolution date relative to training cutoff, outcome entropy, option count $K$---before making any further claim about when social simulation helps.

\subsection{Statistical Power Limitation}
\label{sec:power_limitation}

All three campaigns are critically underpowered for detecting the simulation lift effect at conventional significance thresholds. From the empirical power analysis (Table~\ref{tab:cross_model_stats}):

\begin{itemize}
    \item \textbf{Achieved power:} critically low for Mistral ($8.8\%$) and Qwen ($8.5\%$). Llama's observed-power figure is $91.9\%$ on the corrected repair, computed by the same normal approximation ($\text{power} \approx \Phi(|\text{DM}| - z_{0.975})$) that gave $8.5\%$ before the outage events were re-simulated and $13.8\%$ on the defective intermediate repair. That $8.5 \to 13.8 \to 91.9$ trajectory is a warning about the statistic, not a finding about the design: observed power is a monotone restatement of the observed $p$-value and rises automatically as an effect estimate grows. The design-power statement is unchanged---at $E = 30$, detecting $d = 0.5$ at $80\%$ still requires $E \geq 92$
    \item \textbf{Target:} $80\%$ power (conventional threshold)
    \item \textbf{Required sample:} $E \geq 92$ events to detect a medium effect at $\alpha=0.05$, $1-\beta=0.80$
    \item \textbf{Effect size:} all three models show a negative paired effect ($d_z = -0.613$ for Llama on the corrected repair, against $-0.216$ before it; $-0.301$ for Mistral, $-0.215$ for Qwen), with 95\% CI spanning zero for Mistral and Qwen
\end{itemize}

A note on the effect-size statistic. The pipeline's \texttt{summary.json} computes Cohen's $d$ with the pooled two-independent-samples formula ($-0.286$, $-0.440$, $-0.278$) and records $n_a = n_b = 30$. The design is within-subjects---Conditions A and B are evaluated on the same 30 events---so the paired standardised difference $d_z = \bar{d}/s_d$ is the appropriate statistic, and it is smaller in magnitude. We report $d_z$ in Table~\ref{tab:cross_model_stats} and note that the corresponding power figures, which were computed from the two-sample effect size, are if anything optimistic. The significance conclusion is unaffected: it is non-significant either way.

Consequently the Diebold-Mariano test returns non-significant $p$-values for Mistral and Qwen ($p = 0.88$--$0.95$ against the one-sided alternative that simulation helps). Llama is the exception on the corrected repair: $\text{DM} = -3.36$, two-sided $p = 0.0008$---significant, but in the direction of simulation \emph{harming} calibration, and carrying the serving-stack and evaluator substitutions of \S\ref{sec:outage}, so we do not read it as a clean estimate of a simulation effect. We stress what this does and does not license. It does not license the conclusion that the simulation effect is zero. Neither does it license the converse: earlier drafts argued that the cross-model consistency of the contamination-corrected lift ($\approx +0.20$ on clean events) is ``directionally meaningful,'' but Appendix~\ref{sec:contamination} shows those intervals also span zero and that the subset was selected on the outcome. At $E=30$ events with the observed dispersion, this campaign cannot distinguish a real effect of the magnitude we care about from noise, in either direction. That is the honest summary.

This is an inherent limitation of the benchmark's compute economics: a single three-condition event requires approximately 75--115 minutes of A100 GPU time. Expanding to $N=92$ events per model would require $\sim$115--175 GPU-hours per model family, or approximately $350$--$525$ total GPU-hours for three models---feasible for future campaigns but beyond the scope of the current exploratory study. Three remediation strategies are pursued in future work: (i) expanding the event catalog to $N \geq 92$ resolved post-training-cutoff events; (ii) applying one-sided DM tests with a pre-registered directional hypothesis; and (iii) pooling results across model families in a meta-analytic random-effects framework to increase effective sample size without expanding per-model compute.

\subsection{Evaluator Reliability and the 7-Dimension Rubric}
\label{sec:evaluator_reliability}

All three campaigns report \texttt{evaluator\_unstable: true}, with inter-rater Cohen's $\kappa$ values uniformly below the $\kappa \geq 0.8$ reliability threshold across all seven rubric dimensions. Table~\ref{tab:kappa_summary} reports the per-model kappa landscape.

\begin{table}[!htbp]
\centering
\caption{Evaluator Inter-Rater Reliability: Peak and Representative Cohen's $\kappa$ Values per Model ($\kappa < 0.8$ = unreliable for composite scoring).}
\label{tab:kappa_summary}
\small
\begin{tabular}{lccc}
\toprule
\textbf{Dimension} & \textbf{Llama-3.1-8B} & \textbf{Mistral-7B} & \textbf{Qwen2.5-7B} \\
\midrule
prediction\_accuracy    & 0.321 & 0.182 & 0.088 \\
polarization           & 0.241 & 0.167 & 0.122 \\
deliberation\_quality  & 0.209 & 0.228 & 0.062 \\
herd\_effect           & 0.206 & 0.113 & 0.053 \\
convergence            & 0.245 & 0.139 & 0.212 \\
susceptibility         & \textbf{0.322} & \textbf{0.245} & 0.033 \\
information\_diversity & 0.228 & 0.130 & 0.103 \\
\midrule
\textbf{Max $\kappa$}  & \textbf{0.322} & \textbf{0.245} & \textbf{0.212} \\
\textbf{Dropped dims}  & 7 & 7 & 7 \\
\bottomrule
\end{tabular}
\setlength{\tabcolsep}{6pt}
\end{table}

Several observations are notable. First, \textbf{evaluator reliability degrades systematically with decreased model capacity}: Llama-8B achieves the highest kappa values across all dimensions ($\max \kappa = 0.322$, on \texttt{susceptibility}), while Qwen-7B produces the lowest ($\max \kappa = 0.212$). Llama's peak dimension changed with the 2026-08-16 repair (previously \texttt{prediction\_accuracy} at $0.412$, computed on a mix of real and outage-corrupted units); the ordering across models is unchanged, but the specific value should no longer be read as an especially strong per-dimension result---it dropped along with the rest once computed from genuine data throughout. This degradation may reflect that smaller or less socially active models produce more homogeneous, low-information simulation outputs that are harder for the evaluator to discriminate. Second, \textbf{all 7 dimensions are dropped for all 3 models} under the $\kappa \geq 0.8$ reliability cutoff; consequently, the 7-dimension EPI composite score is not computable for any campaign reported here. Third, \textbf{no dimension approaches acceptable reliability}: even the best-performing dimension (\texttt{susceptibility} for Llama-8B at $\kappa = 0.322$) falls into the ``fair to moderate agreement'' category at best, far below the threshold required for scientific validity.

This instability likely reflects two compounding factors: (1) the rubric MCQ labels (\texttt{very\_low}, \texttt{low}, \texttt{high}, \texttt{very\_high}) are too coarse to reliably discriminate across the simulation output space sampled by different events; and (2) the current evaluator architecture uses a fixed small model (DeepSeek-R1-14B) without per-event calibration or few-shot examples. A third factor, confirmed rather than merely suspected, is added in Section~\ref{sec:epistemic_mapping}: the two evaluator calls this reliability measure compares run with all structured-output enforcement disabled, because the evaluator's model identifier matches a keyword filter in the serving client that is unrelated to the rubric's design. What this section's $\kappa$ values measure is accordingly confounded with that defect and should be re-measured once it is fixed before any rubric redesign is evaluated. Future ECN-BENCH versions will redesign the rubric with finer-grained ordinal scales (e.g., 7-point Likert) and event-specific few-shot calibration prompts to achieve $\kappa \geq 0.6$ (substantial agreement) as a minimum reliability threshold for EPI computation.

\subsection{Convergence Dynamics: JSD Analysis and Monotonic Convergence}
\label{sec:jsd_convergence}

The mean inter-round Jensen-Shannon Divergence (JSD) values and monotonic convergence counts from the convergence telemetry reveal qualitatively distinct opinion dynamics across architectures:

\begin{table}[!htbp]
\centering
\caption{Convergence Telemetry Summary and Integrity Audit (all 30 events, $R=60$, Conditions B and C pooled; 60 units and 600 checkpoint values per model).}
\label{tab:convergence_jsd}
\small
\begin{tabular}{lccccc}
\toprule
\textbf{Model} & \textbf{Mean} & \textbf{Pipeline} & \textbf{Placeholder} & \textbf{Constant} & \textbf{Strictly} \\
 & \textbf{JSD} & \textbf{``monotonic''} & \textbf{values} & \textbf{series} & \textbf{decreasing} \\
\midrule
\textbf{Llama-3.1-8B-AWQ} & 0.734 & 48 & 526/600 (87.7\%) & 48/60 & 29 \\
\textbf{Mistral-7B-Instruct-v0.1}   & 0.637 & 0  & 244/600 (40.7\%) & 0/60  & \textbf{133} \\
\textbf{Qwen2.5-7B-GGUF}  & 0.628 & 0  & 166/600 (27.7\%) & 0/60  & \textbf{113} \\
\bottomrule
\end{tabular}
\smallskip\\
\footnotesize ``Placeholder values'' counts checkpoint readings exactly equal to $0.758277$, the constant the telemetry emits when no probed distribution is available. ``Constant series'' counts units whose ten checkpoint values are all identical. ``Strictly decreasing'' counts round-to-round transitions with a JSD decrease exceeding the pipeline tolerance $\epsilon = 0.002$, out of 540 transitions per model.
\end{table}

Earlier drafts of this paper reported a striking finding from the first two columns of Table~\ref{tab:convergence_jsd}: that \emph{only} Llama-3.1-8B exhibits monotonic JSD convergence (48 episodes) while Mistral-7B and Qwen2.5-7B register zero, and interpreted this as a qualitative difference in opinion dynamics across architectures. A direct audit of the underlying \texttt{round\_jsd} arrays shows that this interpretation is exactly inverted, and we retract it.

The telemetry emits a fixed constant, $0.758277$, whenever a checkpoint has no probed belief distribution to compare. For Llama-3.1-8B this constant accounts for $87.7\%$ of all checkpoint readings, and $48$ of its $60$ units have a completely flat ten-point series. The pipeline's monotonicity test flags a series when JSD never increases; a constant series satisfies that condition vacuously. The reported ``monotonic count'' of $48$ is exactly the number of flat, no-data series. Llama-3.1-8B did not converge---it produced no usable convergence trace at all.

Appendix~\ref{sec:outage} closes this finding by supplying its cause. The reason Llama-3.1-8B in particular has no probed belief distribution to compare is that its serving backend became unreachable during the seventh event of the campaign: from that point every one of the $1{,}000$ telemetry probes per unit failed with \texttt{Connection error} and the \texttt{trace} tables contain zero \texttt{interview} rows, so there was nothing for the JSD checkpoint to read. The $46$ Condition B and C units of the 23 affected events, plus the two units of the transitional event S7, account for the $48$ flat series exactly. The placeholder rate and the outage are the same phenomenon measured through two different fields, which is why the symptom is confined to one model and is total rather than partial within it.

Conversely, Mistral-7B and Qwen2.5-7B have no constant series and register $133$ and $113$ strictly-decreasing round transitions respectively, against Llama-3.1-8B's $29$. They score zero on the pipeline's monotonicity flag because their traces are \emph{sawtooth}: checkpoints at rounds 6, 18, 30 and 42 return the placeholder while the intervening checkpoints return real measurements, so any real decline is interrupted by an artificial rise. Figure~\ref{fig:jsd_integrity} shows both effects. Panel A gives the placeholder share per checkpoint; panel B shows the mean trajectories, in which Mistral-7B declines from $0.758$ to $0.549$ and Qwen2.5-7B to $0.575$ while Llama-3.1-8B remains pinned near the placeholder throughout.

The correct statement of the evidence is therefore: \textbf{Claim 2 (Convergence) is untested by this campaign.} The only two models that produced usable traces both show declining inter-round divergence on the measured checkpoints, which is directionally consistent with DeGroot-style belief updating, but the measurement is punctured often enough that we decline to fit an opinion-dynamics model to it. Repairing the telemetry---ensuring a probed distribution exists at every checkpoint, and replacing the vacuous monotonicity flag with a test that requires a minimum number of valid readings---is a precondition for any convergence claim in a future campaign. We flag this as the second major defect of the July 2026 campaign, alongside the shared Condition A of Section~\ref{sec:condition_a}.

\begin{figure*}[!htbp]
\centering
\includegraphics[width=\textwidth]{figA_jsd_integrity.png}
\caption{Convergence telemetry integrity audit, Conditions B and C pooled. \textbf{Panel A:} share of units emitting the placeholder constant $0.758277$ at each checkpoint. Every model is placeholder-only at rounds 6 and 30; Llama-3.1-8B remains so at every checkpoint for $80$--$100\%$ of units. \textbf{Panel B:} mean inter-round JSD trajectory under Condition B. The sawtooth in the Mistral-7B and Qwen2.5-7B traces is the alternation between placeholder checkpoints and real measurements; Llama-3.1-8B's near-flat line at the placeholder value is the absence of measurement, not the absence of belief change.}
\label{fig:jsd_integrity}
\end{figure*}

\subsection{Consolidated Corrections}
\label{sec:errata_summary}

An audit of this kind is only as credible as its own bookkeeping, so we state plainly what this version changes relative to earlier drafts. Twenty-five claims were withdrawn, corrected, restricted or relocated. Table~\ref{tab:errata_summary} groups them by class of change and names the most consequential instance in each; the full row-by-row table, giving the basis for every individual change, is Table~\ref{tab:errata} in Appendix~\ref{sec:errata_appendix}.

Two patterns in that list are worth naming, because they are the reason the audit was necessary at all rather than incidental to it. First, \emph{every} one of the withdrawn quantitative claims was invisible in the pipeline's own summary files and surfaced only on re-derivation from per-unit records---a completion flag, a monotonic-convergence count and a correlation coefficient were each reported as normal values while measuring nothing. Second, the three most consequential corrections are not about the models under test but about the apparatus: what the evaluator was, whether its output was schema-enforced, and whether the serving backend was alive. A benchmark's least-audited component is its own instrumentation, and that is where this paper's errors concentrated.

\begin{table*}[!htbp]
\centering
\caption{Consolidated errata, condensed by class of change. The full row-by-row table, with the basis for each individual change, is Table~\ref{tab:errata} in Appendix~\ref{sec:errata_appendix}.}
\label{tab:errata_summary}
\small
\begin{tabular}{>{\raggedright\arraybackslash}p{3.5cm}c>{\raggedright\arraybackslash}p{10.4cm}}
\toprule
\textbf{Class of change} & \textbf{$n$} & \textbf{Most consequential instance} \\
\midrule
Withdrawn outright & $9$ &
``Event difficulty explains $71.7\%$ of lift variance ($r = 0.847$)'' is arithmetic, not empirical: lift is $BS_A - BS_B$, so $r = +0.862$ is expected under independent $BS_B$ and the observed value on current data is $+0.837$. The non-tautological association is $r = +0.206$ (Appendix~\ref{sec:tautology}) \\
\addlinespace
Corrected on fact & $8$ &
``Mistral-7B-AWQ'' is \texttt{Mistral-7B-Instruct-v0.1} served at FP16, and the evaluator is DeepSeek-R1-14B rather than the GPT-4o-mini named in the specification (Section~\ref{sec:serving_note}) \\
\addlinespace
Downgraded or restricted in scope & $3$ &
The paper's one surviving positive result is halved: a single chain-of-thought call on the same evidence reproduces, and slightly exceeds, the above-chance directional accuracy ($0.834$ against the swarm's $0.819$), so only the B-versus-C content-sensitivity contrast is evidence about deliberation (Section~\ref{sec:cot_baseline}). Also: clean-subset lift no longer ``confirms Claim~1 across all architectures''---all confidence intervals span zero, and the subset was selected on $BS_A$, the lift's own minuend \\
\addlinespace
Defects newly identified & $3$ &
The evaluator's structured-output enforcement is silently disabled whenever its model string matches an \texttt{"r1"} keyword filter---which \texttt{"deepseek-r1-14b"} does, in all three campaigns (Section~\ref{sec:epistemic_mapping}) \\
\addlinespace
Superseded and relocated & $2$ &
Partly. The pre-release stratum carries the negative lift for Qwen2.5-7B ($-0.25$, excluding zero) and Mistral, consistent with memorisation in the shared reference; for Llama-3.1-8B both strata are negative and the post-release one is the larger (Appendix~\ref{sec:cutoff}) \\
\midrule
\textbf{Total} & $\mathbf{25}$ & \\
\bottomrule
\end{tabular}
\setlength{\tabcolsep}{6pt}
\end{table*}

\section{Conclusion and Future Work}
\label{sec:conclusion}

This paper set out to characterize \textbf{loss of sharpness and degenerate consensus}: the tendency of iterative multi-agent LLM deliberation to erase a forecast's informativeness while its direction survives. Using \textbf{ECN-BENCH}, a multi-agent social simulation testbed (the MiroFish-Offline pipeline), we found that $20.6\%$ of simulated forecasts collapsed to an exactly uniform, non-informative distribution and that the remaining forecasts were measurably flattened relative to direct elicitation---two mechanical effects sufficient on their own to produce an aggregate result that, taken at face value (Section~\ref{sec:full_campaign}), would have read as a straightforward failure of synthetic crowd intelligence.

\textbf{The audit then moved the finding one level down, and that relocation is the paper's contribution.} An earlier draft attributed the pattern to the swarm models' own decoding stacks, because it was concentrated in the two vLLM-served models and absent from the Ollama-served one. Reading the scoring code withdrew that account: the \texttt{probabilities} field is produced by a single fixed evaluator regardless of which swarm model is under test, and that evaluator runs with structured-output enforcement silently disabled by a keyword match on its own model name (Section~\ref{sec:epistemic_mapping}). Re-scoring all $270$ transcripts with a schema-enforced instrument cut the flat rate to $6.1\%$ and dissolved the cross-model pattern entirely (Section~\ref{sec:evaluator_replication}); a dose--response measured within the one campaign free of that confound then showed degeneracy rising from $5\%$ to $65\%$ with input length, while the same transcripts came back sharp under five independent reads (Section~\ref{sec:volume_mechanism}). The phenomenon is real and reproducible from the released artifacts, but it is a property of an unconstrained LLM asked to compress a long transcript into a distribution---not of the deliberation being compressed. What generalises from this paper is therefore not a claim about multi-agent deliberation at all, but a requirement: \textbf{a benchmark that extracts probabilities with an LLM must calibrate that extractor as an instrument}, and flat-unit rate is the cheapest diagnostic that separates an extraction artifact from a result.

We have also documented the systems engineering, containerization, and NFS database lock patches necessary to deploy these simulations reliably on air-gapped HPC resources (Appendix~\ref{sec:architecture}), and the multi-criteria (not purely performance-optimal) selection of the $N=300$, $R=60$ production configuration (Appendix~\ref{sec:optimal_config}) at which the finding was measured, since it required a working large-scale pipeline to surface at all.

The supporting benchmark campaign (Appendices~\ref{sec:full_campaign_supporting} and~\ref{sec:results}) across three model families (Llama-3.1-8B-AWQ, Mistral-7B-Instruct-v0.1, Qwen2.5-7B-GGUF, executed July 2026) yields nine further findings, reported as context for the central result rather than as independently validated claims:

\begin{enumerate}
    \item \textbf{Bimodal event structure is real; its cause is not established.} A stable subset (T2, C11, S9, S3, C3) gains from simulation across all models and another (T1, C10, C15, S6) loses, but the $r=0.847$ correlation previously advanced to explain this is an arithmetic artifact of regressing a difference on its own minuend (Appendix~\ref{sec:tautology}); the non-tautological association $r(BS_A, BS_B)$ is only $+0.206$ across events ($+0.130$ over the $90$ individual units). We withdraw the Difficulty--Lift Predictive Rule.
    \item \textbf{The aggregate lift is partly a contamination effect in the shared reference rather than a simulation effect.} Stratifying on each event's resolution date against each model's release date (Appendix~\ref{sec:cutoff}) localises the negative lift to events predating the model's release for Qwen2.5-7B ($-0.25$, interval excluding zero, against $-0.02$ post-release) and for Mistral. Because Condition A is the shared record, that pattern measures memorisation in the reference, not damage from simulation. It does not hold for Llama-3.1-8B, whose lift is negative in both strata and larger on the post-release events ($-0.24$, excluding zero) that it could not have memorised; that campaign is also the one carrying a serving-stack and evaluator substitution (\S\ref{sec:outage}), so we do not attribute the difference to the model. The date criterion agrees with the earlier $BS_A$-based proxy on only $18$--$21$ of $30$ events.
    \item \textbf{Susceptibility ordering holds in sign for two of three models on the single pass, and strengthens into the paper's one positive result when measured properly.} On the single-pass, full-30-event reading (recomputed 2026-08-16 for Llama post-repair, \S\ref{sec:outage}), $BS_C > BS_B$ for Mistral-7B ($+0.115$) and Qwen2.5-7B ($+0.053$); Llama-3.1-8B's sign depends on which of the three August evaluators reads its transcripts ($-0.122$, $+0.093$, $+0.027$), so we treat it as undetermined (\S\ref{sec:cutoff}). Only Mistral-7B's interval excludes zero. Re-measured on four-read means over the $20$ contamination-clean events common to all four campaigns (Section~\ref{sec:predictive_signal}), the pooled margin is $+0.112$ Brier $[+0.042, +0.182]$ and $+19.4$ points of directional accuracy $[+7.5, +31.3]$ (recomputed 2026-08-21 once Qwen2.5-14B's recovery grew the common clean-event set from $17$ to $20$; was $+0.125$ $[+0.053,+0.195]$ and $+18.4$ $[+6.6,+30.9]$), with two of four per-model Brier intervals excluding zero individually---and the model that carried the single-pass result, Mistral-7B, is one of the two that does not. Being a contrast between two simulated conditions, it never touches the shared Condition A record and survives knowledge-cutoff stratification unchanged. This is the benchmark's most trustworthy positive signal, and Section~\ref{sec:predictive_signal} promotes it from a supporting observation to a stated conclusion.
    \item \textbf{The original Convergence hypothesis is untested, and the previously reported result was inverted.} Section~\ref{sec:jsd_convergence} shows the pipeline's ``monotonic convergence'' count of 48 for Llama-3.1-8B in fact counts flat, no-data series ($87.7\%$ placeholder checkpoints, traced to the backend outage of Appendix~\ref{sec:outage}); Mistral-7B and Qwen2.5-7B, reported as showing zero convergence, are the only two models with usable traces and both decline. No opinion-dynamics claim can be sustained until the telemetry is repaired.
    \item \textbf{Interaction volume does not predict calibration.} Llama-3.1-8B (27,087 actions) and Qwen2.5-7B (723) reach nearly identical Brier scores while the older, unquantised Mistral (9,364 actions) scores worst---correlational and confounded with model age, not architecture. Once memorised events are excluded, the three campaigns' Brier scores fall within $0.024$ of one another (Appendix~\ref{sec:cutoff}), leaving no calibration difference to predict.
    \item \textbf{The campaign's run integrity was not uniform, and the artifacts said so.} Appendices~\ref{sec:persona_audit}--\ref{sec:evidence_depth} find the reference model's serving backend died during its seventh event, collapsing its credited $300$-agent persona pool to $11$ event-independent archetypes and extracting forecasts from connection-error logs on 24 of 30 events (23 fully dead plus the transitional seventh)---none of this visible in the result records, every one of which reported success. Completion flags written by the orchestrator are not evidence; the artifacts must be audited independently of them. These 24 events were re-simulated and re-scored 2026-08-16 (\S\ref{sec:outage}); the audit above documents the defect as originally discovered and is retained for that purpose, not as a description of the current data.
    \item \textbf{Evaluator reliability is insufficient for the 7-dimension rubric.} All seven dimensions fall below the reliability cutoff for all three models ($\max \kappa = 0.322$ for Llama post-repair, on \texttt{susceptibility}; was $0.412$ on \texttt{prediction\_accuracy}); the EPI composite is not computable. Section~\ref{sec:epistemic_mapping} traces this to the same schema-disabling defect as item 8 below: both reliability-check calls run with structured-output enforcement off, so $\kappa$ here measures generation-to-generation noise in an unconstrained model, not disagreement between independent raters.
    \item \textbf{A third measurement channel, Epistemic Mapping, is confirmed non-functional locally, with a verified mechanism.} All $2{,}430$ micro-question instances in the production campaigns return ``Missing in LLM response,'' while $41$ OpenRouter-served pilot runs succeed $237/237$ times. Section~\ref{sec:epistemic_mapping} traces this to \texttt{chat\_json()} disabling structured-output enforcement whenever the evaluator's model string matches a ``reasoning model'' keyword filter---which \texttt{"deepseek-r1-14b"} does, uniformly across all three campaigns---combined with a retry loop that never checks \texttt{micro\_epistemic\_mapping} for completeness before silently backfilling a placeholder. This implicates the evaluator's configuration, not its capability, and also invalidates the vLLM-versus-Ollama account offered in an earlier draft for the flat-forecast finding, since the same fixed evaluator produces \texttt{probabilities} regardless of which swarm model is under test.
    \item \textbf{The benchmark comparison itself is confounded on three axes and underpowered.} Condition A was shared rather than run per model, model generation/precision/serving stack vary jointly across campaigns, and achieved statistical power is $\leq 9\%$ at $E=30$ events. None of this weakens the sharpness-loss finding, which is measured at the $n=180$ per-unit level and does not depend on cross-model or cross-condition significance testing.
\end{enumerate}

Taken together, this paper is best understood as a phenomenon report enabled by, and extracted from, an instructive benchmark failure. The infrastructure works: 270 of 270 units completed, with no evaluator fallbacks, across three model families on an air-gapped cluster. The benchmark comparison built on top of that infrastructure does not yet support its original headline claims, for three reasons we consider generalisable to the field: a baseline that was shared rather than per-model, telemetry that failed silently by emitting a plausible constant instead of an error, and a headline correlation that was arithmetic rather than empirical. Each was invisible in the aggregate summary files and surfaced only on re-derivation from per-unit records.

Two things survived that audit, and both sharpened rather than weakened under it. The first is the sharpness-loss finding itself: a documentable tendency for the extraction channel to trade forecast informativeness for consensus, now localised to an interaction between long deliberation transcripts and one schema-disabled instrument (Section~\ref{sec:volume_mechanism}), and worth studying independently of whether the surrounding benchmark's other claims hold up. The second is substantive rather than methodological, and we state it as a positive result rather than as a caveat: \textbf{on events that resolved after every model's release date, this multi-agent simulation produces forecasts that are strongly informative about how real social events turned out}---the correct outcome carries the most probability mass in $82\%$ of (model, event) pairs against a $33\%$ chance baseline, $p = 3.7\times10^{-22}$---\textbf{and those forecasts demonstrably track the content of the evidence deliberated over}, degrading by $0.112$ Brier $[+0.042, +0.182]$ when the injected evidence is replaced with topically unrelated material (Section~\ref{sec:predictive_signal}). Multi-agent LLM simulation is therefore not merely producing plausible-looking text about social events; on this catalogue it carries real, content-dependent signal about their outcomes.

That signal is not, however, shown to come from the deliberation. A single-agent baseline run after the audit---one retrieval-augmented chain-of-thought call per event, given the same seed dossier and the same injected update the swarm's agents received---reproduces, and on the identity-matched pool slightly exceeds, the directional accuracy outright ($0.834$ against the swarm's $0.819$ pooled over all four arms) and carries the lower mean Brier on both arms where the single agent and the campaign use the same checkpoint, though at $20$ events no interval separates the two ($-0.106$ Brier $[-0.226, +0.009]$, clustered $p = 0.109$; Section~\ref{sec:cot_baseline}). The honest statement is a non-superiority one: three hundred agents deliberating for sixty rounds, plus a transcript and an evaluator to compress it, buy no measurable accuracy over a single call, and might cost some. The content-sensitivity result stands apart from this, being a contrast between two simulated conditions for which the single-agent run has no analogue---so what remains specifically attributable to deliberation is that it responds to which evidence it is given, not that it forecasts better than the model does alone. Converting a demonstrated signal into a demonstrated \emph{benefit} now requires two experiments rather than one: the per-model Condition A re-run, and a single-agent arm large enough to resolve a $0.11$ Brier difference, which at the observed variance ($d_z = 0.39$) means $E = 55$ events for $80\%$ power rather than the $20$ available here, at which the achieved power is $37\%$.

Beyond the central finding, this paper establishes an empirically validated \textbf{Sub-15B Exclusion Criterion} (Appendix~\ref{sec:small_model_limits}) for local vLLM HPC execution, rooted in three compounding failure modes: context-window overflow, VRAM saturation, and guided-decoding JSON schema incompatibility. Both Gemma 1.1 2B-IT and CodeGemma 7B-IT were confirmed unable to pass Stage 1 persona extraction, immediately raising \texttt{ValueError: Intelligent persona generation failed} and aborting.

\textbf{Future work} is ordered by what must be fixed before anything else is worth measuring.

\emph{Prerequisites, in order.} (i) \textbf{Partially discharged by Section~\ref{sec:evaluator_replication}.} The experiment this list previously named as highest-priority---restore structured-output enforcement and re-measure the flat-forecast rate---has now been run, by re-scoring all $270$ transcripts with a second evaluator whose enforcement is active rather than by patching the keyword filter in \texttt{chat\_json()}. The flat-forecast rate falls from $20.6\%$ to $6.1\%$ and Condition B degeneracy disappears entirely, which is why this paper now locates the phenomenon in the evaluator. Two parts of the original item remain open and should be done by fixing the defect in place---excluding \texttt{"deepseek-r1-14b"} from the \texttt{"reasoning model"} keyword filter, aliasing it to a served-model-name that does not match \texttt{"r1"}, or adding an explicit code path restoring \texttt{json\_object} enforcement for local reasoning-named evaluators---so that the \emph{same} evaluator can be compared with and without enforcement: whether enforcement raises rubric $\kappa$, and whether it eliminates the micro-question failure, are not answered by swapping to a different model. (ii) \textbf{Calibrate the evaluator as an instrument before drawing any further inference from it.} Section~\ref{sec:replication_noise} finds $4/30$ reproducibility on byte-identical input, with a mean Brier spread of $0.128$ and a maximum of $0.715$, so every future campaign should run the evaluator $k \geq 3$ times per unit and report the between-run spread, and should score with at least two independent evaluators. Any effect smaller than the instrument's own spread is not measurable by this design.

The remaining prerequisite items and the substantive testbed extensions (re-running Condition A per model, repairing the convergence telemetry, decomposing the Brier score, and eight further items) are testbed engineering rather than central-finding follow-ups; we list them in full in Appendix~\ref{sec:extended_future_work}. Recovering the Qwen2.5-14B campaign, listed there in earlier drafts, is no longer outstanding: Appendix~\ref{sec:fourth_model_status} backfills it and Section~\ref{sec:four_model_extension} scores it against the other three models.


% \appendices (plural), not \appendix: this paper has six appendix sections and
% IEEEtran's singular form declares exactly one, which makes every \section after
% it resolve to the same letter ("A") and pushes the real numbering down onto the
% subsections. That is why cross-references used to render as "Appendix M4".
\appendices


\section{Consolidated Errata}
\label{sec:errata_appendix}

Table~\ref{tab:errata} lists every claim withdrawn, corrected or restricted relative to earlier drafts of this paper, so that readers of those drafts can locate the changes without re-reading it. The material subset is summarised in Section~\ref{sec:errata_summary}.

\begin{table*}[!htbp]
\centering
\caption{Consolidated errata: claims withdrawn or revised relative to earlier drafts.}
\label{tab:errata}
\small
\begin{tabular}{>{\raggedright\arraybackslash}p{4.4cm}>{\raggedright\arraybackslash}p{5.6cm}>{\raggedright\arraybackslash}p{5.6cm}}
\toprule
\textbf{Claim in earlier drafts} & \textbf{Status} & \textbf{Basis} \\
\midrule
Event difficulty explains $71.7\%$ of lift variance ($r = 0.847$) &
\textbf{Withdrawn.} The non-tautological association is $r = +0.206$ &
$Lift \equiv BS_A - BS_B$; $r = +0.862$ expected under independent $BS_B$ (\S\ref{sec:tautology}) \\
\addlinespace
Difficulty--Lift Predictive Rule as a budget-screening heuristic &
\textbf{Withdrawn} entirely &
Follows from the above \\
\addlinespace
Only Llama-3.1-8B shows monotonic JSD convergence (48 episodes) &
\textbf{Withdrawn and inverted.} Llama-3.1-8B produced no usable trace; Mistral-7B and Qwen2.5-7B are the only models with real telemetry &
$87.7\%$ of Llama-3.1-8B checkpoints are the constant $0.758277$; 48 = its count of flat series (\S\ref{sec:jsd_convergence}) \\
\addlinespace
Qwen2.5-7B shows a susceptibility sign reversal ($-0.044$) &
\textbf{Corrected} to $+0.053$, same direction as the other two models &
$-0.044$ was a yes-probability delta, not the Brier margin (\S\ref{sec:cross_model_stats}) \\
\addlinespace
The same 11 events are contaminated ``across all three independently trained model families'' &
\textbf{Withdrawn as evidence.} Agreement is guaranteed by construction &
Condition A is one shared record (\S\ref{sec:condition_a}) \\
\addlinespace
Clean-subset lift confirms Claim 1 across all architectures &
\textbf{Downgraded} to directional, non-significant, selection-inflated &
All CIs span zero; subset selected on $BS_A$, the lift's own minuend (\S\ref{sec:contamination}) \\
\addlinespace
Echo-chamber collapse boundary at $R \approx 70$--$80$ &
\textbf{Withdrawn} as a quantitative claim &
Interpolated from two events, one run per cell; C9 delta is $+0.022$ (\S\ref{sec:r_scaling}) \\
\addlinespace
OASIS topology is power-law with $\alpha \approx 2.3$ &
\textbf{Withdrawn.} Not measured &
Manifest records \texttt{oasis-default-unspecified} (\S\ref{sec:interaction_realism}) \\
\addlinespace
Evaluator is GPT-4o-mini &
\textbf{Corrected} to DeepSeek-R1-14B for these campaigns &
Air-gapped compute nodes preclude external API calls (\S\ref{sec:confound}) \\
\addlinespace
Epistemic Alignment Score (micro-question / macro-prediction correlation) not previously reported &
\textbf{Newly audited: $0/2{,}430$ non-functional in the production campaigns, but $237/237$ functional across 41 OpenRouter-served pilot runs.} Root cause confirmed from source: evaluator model name \texttt{"deepseek-r1-14b"} matches a \texttt{"reasoning model"} keyword filter that disables all structured-output enforcement, and the retry loop never checks this one field &
\texttt{chat\_json()} sets \texttt{response\_format=None} for any \texttt{self.model} containing \texttt{"r1"}; \texttt{\_validate\_micro\_mapping()} silently backfills ``Missing in LLM response'' (\S\ref{sec:epistemic_mapping}) \\
\addlinespace
Flat-forecast rate attributed to swarm model's vLLM-vs-Ollama serving backend &
\textbf{Withdrawn.} \texttt{probabilities} is generated by the fixed evaluator, not the swarm model; its schema-disabled condition is constant across all three campaigns and cannot explain a cross-model difference &
Same source-code audit as above; open hypothesis is \texttt{evidence\_text} volume, not serving backend (\S\ref{sec:degenerate_flat}) \\
\addlinespace
$BS = \frac{1}{2}\sum_k (f_k - o_k)^2$, bounded in $[0,1]$ &
\textbf{Corrected} to the unnormalised form, bounded in $[0,2]$ &
21 reported unit scores exceed $1.0$; maximum is exactly $2.0$ (\S\ref{sec:grading}) \\
\addlinespace
Per-event tables verified by spot-check; max discrepancy $\approx 0.18$ &
\textbf{Corrected.} 18 cells wrong, max discrepancy $0.93$ &
Full cell-by-cell recomputation (item 1 above) \\
\addlinespace
SBE $\cup$ SDE partitions the catalogue; SDE $n=19$ &
\textbf{Corrected} to $n=21$; C13 and T3 were unassigned &
$9 + 19 = 28 \neq 30$ (\S\ref{sec:sbe_sde}) \\
\addlinespace
``Mistral-7B-AWQ'', served vLLM+AWQ INT4 &
\textbf{Corrected} to \texttt{Mistral-7B-Instruct-v0.1} at FP16 &
Slurm \texttt{BENCH\_MODEL} resolves to the unquantised repo; no \texttt{--quantization} flag (\S\ref{sec:serving_note}) \\
\addlinespace
``Qwen2.5-14B-AWQ'', served via vLLM &
\textbf{Corrected} to \texttt{qwen2.5:14b} GGUF via Ollama &
vLLM block commented out; \texttt{ollama pull} is the live path (\S\ref{sec:serving_note}) \\
\addlinespace
Mistral's deficit is architecture-dependent over-deliberation &
\textbf{Withdrawn.} Comparison confounded on three axes &
Mistral is 10 months older, unquantised, and the only FP16 model (\S\ref{sec:serving_note}) \\
\addlinespace
At $N=300$ each persona axis is represented by $\geq 75$ distinct agents &
\textbf{Withdrawn as a measured property.} Realised distinct personas per 300-agent pool are $296$ (Mistral), $40$ (Qwen-7B) and $11$ on 23 of 30 events (Llama); $N$ bounds diversity from above only &
Direct count over all $270$ persona files, $81{,}000$ agent records (\S\ref{sec:persona_audit}) \\
\addlinespace
All 270 units completed with no errors and no evaluator fallbacks &
\textbf{Corrected.} True of the recorded fields, false of the runs: 69 Llama units executed against an unreachable endpoint while recording \texttt{error: null} and \texttt{evaluator\_fallback\_used: false} &
$1{,}000$ failed telemetry probes and $0$ \texttt{interview} rows per affected unit (\S\ref{sec:outage}) \\
\addlinespace
Llama-8B shows the richest social interaction ($27{,}087$ actions), best supporting emergent deliberation &
\textbf{Restricted.} Holds over the sweep databases; within the scored 30-event campaign Mistral posts $2\times$ more, and $96.7\%$ of Llama's actions come from 6 of its 30 events &
120 databases per model for the scored campaign vs.\ 672/264/264 in the original audit (Tab.~\ref{tab:action_30event_only}) \\
\addlinespace
Cause of Llama's $87.7\%$ JSD placeholder rate not established &
\textbf{Resolved.} The placeholder and the backend outage are one phenomenon: all $48$ flat series are exactly the $48$ units of the $24$ affected events, with no exceptions in either direction &
Zero \texttt{interview} rows leaves the JSD checkpoint nothing to read (\S\ref{sec:outage}) \\
\addlinespace
Rubric unreliability is a $\kappa$ problem &
\textbf{Extended.} The rubric is also insensitive to its object: single-speaker error-log evidence scores $0.495$ against $0.508$ for a 40-speaker transcript, within model &
Parse of \texttt{evidence\_text} for all $180$ simulated units (\S\ref{sec:evidence_depth}) \\
\addlinespace
Contamination split by $BS_A < 0.15$; date-based test deferred to future work &
\textbf{Superseded and carried out.} Resolution-date stratification agrees with the proxy on only $18$--$21$ of $30$ events; the reversal survives but relocates to the shared reference &
Release dates vs.\ \texttt{events\_raw.json} resolution dates (\S\ref{sec:cutoff}) \\
\addlinespace
Aggregate lift $\Delta = -0.097$ to $-0.147$ &
\textbf{Partly localised.} $-0.25$ on Qwen's pre-release stratum (excluding zero) against $-0.02$ post-release; for Llama both strata are negative and the post-release one is larger &
Cutoff stratification, $B = 20{,}000$ bootstrap (Tab.~\ref{tab:cutoff_lift}) \\
\addlinespace
The swarm's above-chance directional accuracy on clean events is evidence that the simulation carries predictive signal &
\textbf{Restricted.} Evidence that the \emph{base models} do. A single chain-of-thought call on the same dossier and injection reaches $0.834$ against the swarm's $0.819$, and no arm shows a swarm advantage separable from zero ($-0.106$ Brier $[-0.226, +0.009]$ on the identity-matched pool). The B-versus-C contrast is unaffected and remains the deliberation-specific part of the finding &
$240$ single-agent calls over the same $20$ events, four models, three replicates; cluster bootstrap over events (\S\ref{sec:cot_baseline}) \\
\bottomrule
\end{tabular}
\setlength{\tabcolsep}{6pt}
\end{table*}

\section{Per-Event Brier Score Tables}
\label{sec:per_event_tables}

The three tables below give the per-event Condition A/B/C Brier scores behind the aggregates of Section~\ref{sec:full_campaign}. They are reference tables: every cell was recomputed from the released per-unit records, and the errors this recomputation found are recorded in Appendix~\ref{sec:errata_appendix}.

\begin{table*}[!htbp]
\centering
\caption{Per-Event Brier Scores --- Llama-3.1-8B-AWQ ($N=300$, $R=60$). Bold B-score indicates With-Sim $<$ No-Sim (simulation lifts calibration).}
\label{tab:per_event_llama_8b}
\small
\setlength{\tabcolsep}{4pt}
\begin{tabular}{llccccc}
\toprule
\textbf{Event} & \textbf{Type} & \textbf{Cond A (No-Sim)} & \textbf{Cond B (With-Sim)} & \textbf{Cond C (Null)} & \textbf{Lift$_B$} & \textbf{Dir.Acc$_B$} \\
\midrule
\textbf{C1} & Social/Electoral & 0.5600 & $1.0550$ & 0.5150 & $-0.4950$ & 0 \\
\textbf{C10} & Social/Electoral & 0.0036 & $0.0036$ & 0.2834 & $+0.0000$ & 1 \\
\textbf{C11} & Social/Electoral & 1.3150 & $1.3400$ & 0.7550 & $-0.0250$ & 0 \\
\textbf{C12} & Social/Electoral & 0.0005 & $0.1150$ & 0.1034 & $-0.1145$ & 1 \\
\textbf{C13} & Social/Electoral & 0.2273 & $1.3880$ & 1.2500 & $-1.1607$ & 0 \\
\textbf{C14} & Social/Electoral & 0.0900 & $1.0295$ & 0.6100 & $-0.9395$ & 0 \\
\textbf{C15} & Social/Electoral & 0.0050 & $0.9100$ & 0.5700 & $-0.9050$ & 0 \\
\textbf{C2} & Social/Electoral & 0.0050 & $0.9200$ & 1.3400 & $-0.9150$ & 0 \\
\textbf{C3} & Social/Electoral & 0.8100 & $\mathbf{0.7500}$ & 0.7500 & $\mathbf{+0.0600}$ & $\frac{1}{2}$ \\
\textbf{C4} & Social/Electoral & 0.0166 & $0.5000$ & 0.0342 & $-0.4834$ & 1 \\
\textbf{C5} & Social/Electoral & 0.4498 & $1.1450$ & 1.0800 & $-0.6952$ & 0 \\
\textbf{C6} & Social/Electoral & 1.0450 & $1.4150$ & 1.1950 & $-0.3700$ & 0 \\
\textbf{C7} & Social/Electoral & 0.3200 & $0.7000$ & 0.8333 & $-0.3800$ & $\frac{1}{2}$ \\
\textbf{C8} & Social/Electoral & 0.4850 & $0.8000$ & 0.8200 & $-0.3150$ & $\frac{1}{2}$ \\
\textbf{C9} & Social/Electoral & 0.2150 & $0.6534$ & 0.2900 & $-0.4384$ & 1 \\
\textbf{S2} & Senate/Binary & 0.0450 & $0.5000$ & 0.5000 & $-0.4550$ & $\frac{1}{2}$ \\
\textbf{S3} & Senate/Binary & 0.7200 & $\mathbf{0.5000}$ & 0.5000 & $\mathbf{+0.2200}$ & $\frac{1}{2}$ \\
\textbf{S4} & Senate/Binary & 0.3200 & $0.5000$ & 0.5000 & $-0.1800$ & $\frac{1}{2}$ \\
\textbf{S5} & Senate/Binary & 0.5000 & $0.5000$ & 0.5000 & $+0.0000$ & $\frac{1}{2}$ \\
\textbf{S6} & Senate/Binary & 0.1368 & $0.7500$ & 0.7500 & $-0.6132$ & $\frac{1}{2}$ \\
\textbf{S7} & Senate/Binary & 0.2450 & $0.9800$ & 0.7200 & $-0.7350$ & 0 \\
\textbf{S9} & Senate/Binary & 0.7200 & $0.9800$ & 0.9800 & $-0.2600$ & 0 \\
\textbf{T1} & Tech/Economic & 0.0450 & $0.1800$ & 0.3200 & $-0.1350$ & 1 \\
\textbf{T2} & Tech/Economic & 1.2800 & $\mathbf{0.0450}$ & 0.0800 & $\mathbf{+1.2350}$ & 1 \\
\textbf{T3} & Tech/Economic & 0.5450 & $\mathbf{0.2850}$ & 0.2850 & $\mathbf{+0.2600}$ & 1 \\
\textbf{T4} & Tech/Economic & 0.0538 & $1.2000$ & 0.1150 & $-1.1462$ & 0 \\
\textbf{T5} & Tech/Economic & 0.1800 & $\mathbf{0.0450}$ & 0.0450 & $\mathbf{+0.1350}$ & 1 \\
\textbf{T6} & Tech/Economic & 0.7150 & $0.7700$ & 0.9000 & $-0.0550$ & 0 \\
\textbf{T8} & Tech/Economic & 0.4850 & $0.5100$ & 0.2538 & $-0.0250$ & 1 \\
\textbf{T9} & Tech/Economic & 0.0000 & $0.0800$ & 0.0200 & $-0.0800$ & 1 \\
\midrule
\textbf{Mean} & --- & \textbf{0.3846} & \textbf{0.6850} & \textbf{0.5633} & $-0.3004$ & 0.467 \\
\bottomrule
\end{tabular}
\setlength{\tabcolsep}{6pt}
\end{table*}

\begin{table*}[!htbp]
\centering
\caption{Per-Event Brier Scores --- Mistral-7B-Instruct-v0.1 ($N=300$, $R=60$). This table was regenerated directly from \texttt{event\_results.json}; earlier drafts contained 18 incorrect cells across nine C-prefixed events (see Appendix~\ref{sec:data_audit}).}
\label{tab:per_event_mistral_7b}
\small
\setlength{\tabcolsep}{4pt}
\begin{tabular}{llccccc}
\toprule
\textbf{Event} & \textbf{Type} & \textbf{Cond A (No-Sim)} & \textbf{Cond B (With-Sim)} & \textbf{Cond C (Null)} & \textbf{Lift$_B$} & \textbf{Dir.Acc$_B$} \\
\midrule
\textbf{C1} & Social/Electoral & 0.5600 & $0.6667$ & 0.5600 & $-0.1067$ & $\frac{1}{2}$ \\
\textbf{C10} & Social/Electoral & 0.0036 & $0.8000$ & 0.6950 & $-0.7964$ & $\frac{1}{2}$ \\
\textbf{C11} & Social/Electoral & 1.3150 & $\mathbf{0.6250}$ & 0.8000 & $\mathbf{+0.6900}$ & 1 \\
\textbf{C12} & Social/Electoral & 0.0005 & $0.5926$ & 1.2000 & $-0.5921$ & $\frac{1}{2}$ \\
\textbf{C13} & Social/Electoral & 0.2273 & $\mathbf{0.0480}$ & 0.6000 & $\mathbf{+0.1793}$ & 1 \\
\textbf{C14} & Social/Electoral & 0.0900 & $0.7250$ & 0.7150 & $-0.6350$ & 0 \\
\textbf{C15} & Social/Electoral & 0.0050 & $0.2083$ & 0.3000 & $-0.2033$ & 1 \\
\textbf{C2} & Social/Electoral & 0.0050 & $1.1338$ & 0.6200 & $-1.1288$ & 0 \\
\textbf{C3} & Social/Electoral & 0.8100 & $\mathbf{0.7000}$ & 1.1350 & $\mathbf{+0.1100}$ & 0 \\
\textbf{C4} & Social/Electoral & 0.0166 & $0.1368$ & 0.5000 & $-0.1202$ & 1 \\
\textbf{C5} & Social/Electoral & 0.4498 & $0.8000$ & 0.8000 & $-0.3502$ & $\frac{1}{2}$ \\
\textbf{C6} & Social/Electoral & 1.0450 & $\mathbf{0.6650}$ & 0.8000 & $\mathbf{+0.3800}$ & 0 \\
\textbf{C7} & Social/Electoral & 0.3200 & $0.8000$ & 0.7300 & $-0.4800$ & 0 \\
\textbf{C8} & Social/Electoral & 0.4850 & $0.7500$ & 0.8000 & $-0.2650$ & 0 \\
\textbf{C9} & Social/Electoral & 0.2150 & $0.5950$ & 0.8100 & $-0.3800$ & 1 \\
\textbf{S2} & Senate/Binary & 0.0450 & $0.5202$ & 0.5000 & $-0.4752$ & 0 \\
\textbf{S3} & Senate/Binary & 0.7200 & $\mathbf{0.3200}$ & 0.5000 & $\mathbf{+0.4000}$ & 1 \\
\textbf{S4} & Senate/Binary & 0.3200 & $0.7200$ & 0.5000 & $-0.4000$ & 0 \\
\textbf{S5} & Senate/Binary & 0.5000 & $\mathbf{0.3200}$ & 0.3200 & $\mathbf{+0.1800}$ & 1 \\
\textbf{S6} & Senate/Binary & 0.1368 & $0.7500$ & 0.7500 & $-0.6132$ & $\frac{1}{2}$ \\
\textbf{S7} & Senate/Binary & 0.2450 & $\mathbf{0.1800}$ & 1.1250 & $\mathbf{+0.0650}$ & 1 \\
\textbf{S9} & Senate/Binary & 0.7200 & $\mathbf{0.0800}$ & 0.5000 & $\mathbf{+0.6400}$ & 1 \\
\textbf{T1} & Tech/Economic & 0.0450 & $0.5000$ & 0.5000 & $-0.4550$ & $\frac{1}{2}$ \\
\textbf{T2} & Tech/Economic & 1.2800 & $\mathbf{0.0200}$ & 0.5000 & $\mathbf{+1.2600}$ & 1 \\
\textbf{T3} & Tech/Economic & 0.5450 & $0.7800$ & 0.4650 & $-0.2350$ & 0 \\
\textbf{T4} & Tech/Economic & 0.0538 & $0.5300$ & 0.4600 & $-0.4762$ & 1 \\
\textbf{T5} & Tech/Economic & 0.1800 & $0.5000$ & 0.1800 & $-0.3200$ & $\frac{1}{2}$ \\
\textbf{T6} & Tech/Economic & 0.7150 & $0.7500$ & 0.9000 & $-0.0350$ & $\frac{1}{2}$ \\
\textbf{T8} & Tech/Economic & 0.4850 & $0.5400$ & 0.8000 & $-0.0550$ & 1 \\
\textbf{T9} & Tech/Economic & 0.0000 & $0.1800$ & 0.3200 & $-0.1800$ & 1 \\
\midrule
\textbf{Mean} & --- & \textbf{0.3846} & \textbf{0.5312} & \textbf{0.6462} & $-0.1466$ & 0.567 \\
\bottomrule
\end{tabular}
\smallskip \\
\footnotesize A Dir.Acc entry of $\frac{1}{2}$ marks a unit whose modal forecast is tied, so no direction can be scored; the pipeline records these as $0.5$. Tied units are exactly the flat (uniform) forecasts discussed in Section~\ref{sec:degenerate}. The other two per-event tables encode these as $0$, which slightly understates their directional accuracy; the aggregate Dir.Acc$_B$ figures throughout the paper use the pipeline's $0.5$ convention and are unaffected.
\end{table*}

\begin{table*}[!htbp]
\centering
\caption{Per-Event Brier Scores --- Qwen2.5-7B (GGUF Q4\_K\_M via Ollama, $N=300$, $R=60$). Serving stack differs from AWQ/vLLM models; see Section~\ref{sec:serving_note}.}
\label{tab:per_event_qwen_7b}
\small
\setlength{\tabcolsep}{4pt}
\begin{tabular}{llccccc}
\toprule
\textbf{Event} & \textbf{Type} & \textbf{Cond A (No-Sim)} & \textbf{Cond B (With-Sim)} & \textbf{Cond C (Null)} & \textbf{Lift$_B$} & \textbf{Dir.Acc$_B$} \\
\midrule
\textbf{C1} & Social/Electoral & 0.5600 & $\mathbf{0.4650}$ & 0.3350 & $\mathbf{+0.0950}$ & 1 \\
\textbf{C10} & Social/Electoral & 0.0036 & $0.6400$ & 0.4800 & $-0.6364$ & 0 \\
\textbf{C11} & Social/Electoral & 1.3150 & $\mathbf{0.5500}$ & 1.0550 & $\mathbf{+0.7650}$ & 1 \\
\textbf{C12} & Social/Electoral & 0.0005 & $0.1150$ & 0.1338 & $-0.1145$ & 1 \\
\textbf{C13} & Social/Electoral & 0.2273 & $0.6100$ & 0.4438 & $-0.3827$ & 1 \\
\textbf{C14} & Social/Electoral & 0.0900 & $0.2250$ & 0.5000 & $-0.1350$ & 1 \\
\textbf{C15} & Social/Electoral & 0.0050 & $0.6750$ & 0.6350 & $-0.6700$ & 0 \\
\textbf{C2} & Social/Electoral & 0.0050 & $0.1150$ & 0.6400 & $-0.1100$ & 1 \\
\textbf{C3} & Social/Electoral & 0.8100 & $\mathbf{0.5850}$ & 0.9000 & $\mathbf{+0.2250}$ & 1 \\
\textbf{C4} & Social/Electoral & 0.0166 & $0.1250$ & 0.1650 & $-0.1084$ & 1 \\
\textbf{C5} & Social/Electoral & 0.4498 & $\mathbf{0.1400}$ & 0.6531 & $\mathbf{+0.3098}$ & 1 \\
\textbf{C6} & Social/Electoral & 1.0450 & $1.0450$ & 1.0550 & $-0.0000$ & 0 \\
\textbf{C7} & Social/Electoral & 0.3200 & $0.5300$ & 0.7950 & $-0.2100$ & 1 \\
\textbf{C8} & Social/Electoral & 0.4850 & $1.2000$ & 0.6200 & $-0.7150$ & 0 \\
\textbf{C9} & Social/Electoral & 0.2150 & $0.6032$ & 0.5750 & $-0.3882$ & 1 \\
\textbf{S2} & Senate/Binary & 0.0450 & $0.1800$ & 0.3200 & $-0.1350$ & 1 \\
\textbf{S3} & Senate/Binary & 0.7200 & $0.8450$ & 0.9800 & $-0.1250$ & 0 \\
\textbf{S4} & Senate/Binary & 0.3200 & $0.3200$ & 0.1800 & $0.0000$ & 1 \\
\textbf{S5} & Senate/Binary & 0.5000 & $\mathbf{0.3200}$ & 0.3200 & $\mathbf{+0.1800}$ & 1 \\
\textbf{S6} & Senate/Binary & 0.1368 & $0.4150$ & 0.5000 & $-0.2782$ & 1 \\
\textbf{S7} & Senate/Binary & 0.2450 & $\mathbf{0.1250}$ & 0.3200 & $\mathbf{+0.1200}$ & 1 \\
\textbf{S9} & Senate/Binary & 0.7200 & $\mathbf{0.3200}$ & 0.9800 & $\mathbf{+0.4000}$ & 1 \\
\textbf{T1} & Tech/Economic & 0.0450 & $0.9800$ & 0.3200 & $-0.9350$ & 0 \\
\textbf{T2} & Tech/Economic & 1.2800 & $\mathbf{0.0050}$ & 0.7200 & $\mathbf{+1.2750}$ & 1 \\
\textbf{T3} & Tech/Economic & 0.5450 & $\mathbf{0.1950}$ & 0.4550 & $\mathbf{+0.3500}$ & 1 \\
\textbf{T4} & Tech/Economic & 0.0538 & $0.6200$ & 0.2150 & $-0.5662$ & 1 \\
\textbf{T5} & Tech/Economic & 0.1800 & $0.6050$ & 0.3200 & $-0.4250$ & 0 \\
\textbf{T6} & Tech/Economic & 0.7150 & $1.0150$ & 0.7850 & $-0.3000$ & 0 \\
\textbf{T8} & Tech/Economic & 0.4850 & $0.5550$ & 0.5550 & $-0.0700$ & 1 \\
\textbf{T9} & Tech/Economic & 0.0000 & $0.3200$ & 0.0800 & $-0.3200$ & 1 \\
\midrule
\textbf{Mean} & --- & \textbf{0.3846} & \textbf{0.4814} & \textbf{0.5345} & $-0.0968$ & 0.733 \\
\bottomrule
\end{tabular}
\setlength{\tabcolsep}{6pt}
\end{table*}
\section{Supporting Analysis and Run-Integrity Audit}
\label{sec:full_campaign_supporting}

This appendix collects the rest of the July 2026 campaign's findings---event-level structure, interaction logs, contamination analysis, and the run-integrity audit---together with the long-form logs behind them. It is context for the audit of Part~I, not a set of independently validated benchmark results: every quantity here carries the same caveats as the body (shared Condition~A reference, three-way provenance confound, low statistical power) and, where it derives from a score rather than from an artifact on disk, comes from the single Evaluator~J pass whose defect Section~\ref{sec:epistemic_mapping} documents. It is placed in an appendix for that reason, and because several of its subsections are reference logs rather than argument. The claims it withdraws or revises are consolidated in Appendix~\ref{sec:errata_appendix}.

\subsection{Simulation-Benefiting vs.~Simulation-Degrading Events: A Structural Analysis}
\label{sec:sbe_sde}

\begin{figure*}[!htbp]
\centering
\includegraphics[width=\textwidth]{fig6_volatility_vs_lift.png}
\caption{Event Difficulty (BS$_A$) vs.~Simulation Lift (BS$_A$ - BS$_B$). Each point represents one of the 90 (event, model) pairs, with marker shape encoding model architecture and point color encoding taxonomy category. The apparent correlation ($r = 0.847$, the value as originally reported) is an arithmetic artifact: the vertical axis contains the horizontal axis as its minuend, and $r = +0.862$ is expected here even under an uninformative BS$_B$. This panel is the pre-repair rendering; on the corrected data of \S\ref{sec:outage} the same relationship is $r = +0.837$, which changes none of the argument. See Figure~\ref{fig:tautology} for the same data with the shared term removed ($r = +0.206$) and Appendix~\ref{sec:tautology} for the derivation. This figure is retained to document the relationship as originally reported; it should not be read as evidence that baseline difficulty predicts simulation benefit. Annotated events show the most extreme cases.}
\label{fig:volatility_lift}
\end{figure*}

A central empirical finding of this campaign is the existence of two structurally distinct event classes, which we formalize as \textbf{Simulation-Benefiting Events} (SBE) and \textbf{Simulation-Degrading Events} (SDE).

\subsubsection{SBE Structural Profile}

Figure~\ref{fig:volatility_lift} plots all 90 (event, model) pairs on the difficulty--lift plane from which the two clusters are read; as its caption and Appendix~\ref{sec:tautology} set out, the trend visible there is arithmetic, so the clusters below are a descriptive partition rather than evidence of a difficulty--benefit relationship.

The SBE cluster (T2, C11, S9, S3, C3, C6, S5, C1, S7; $n=9$) exhibits the following defining characteristics:

\begin{itemize}
    \item \textbf{High baseline uncertainty:} Mean $BS_A = 0.799$ (vs.~$0.188$ for SDE). These are events where the No-Sim direct elicitation is poorly calibrated, indicating the LLM has limited or contradictory prior knowledge.
    \item \textbf{Strong B$<$C gap:} Mean B-to-C gap $= +0.238$, indicating that social deliberation with relevant information (Cond B) substantially outperforms null noise injection (Cond C). Unlike the lift, this contrast shares no term between its two sides and is therefore not subject to the artifact of Appendix~\ref{sec:tautology}.
    \item \textbf{Cross-model consistency:} SBE events show positive lift across all 3 models in 5/9 cases ($55.6\%$): T2, C11, S9, C3 and S7. The remaining four (S3, S5, C1 with 2/3; C6 with 1/3) are majority- rather than unanimously-positive.
    \item \textbf{Diverse taxonomy:} SBE events span all three taxonomy categories (4 Social/Electoral, 4 Senate/Binary, 1 Technical/Economic), suggesting high-uncertainty events exist across domains.
\end{itemize}

\subsubsection{SDE Structural Profile}

The SDE cluster is the complement of the SBE set: T1, C10, C15, S6, T5, C2, C9, S2, C14, C12, C7, S4, T9, T4, T6, T8, C8, C5, C4, plus C13 and T3 ($n=21$; earlier drafts listed only 19 and left C13 and T3 unassigned, so that the two clusters summed to 28 rather than 30). It exhibits contrasting characteristics:

\begin{itemize}
    \item \textbf{High prior certainty:} Mean $BS_A = 0.207$ over the corrected $n=21$ set. These events already have correct directional forecasts under the shared Condition A reference.
    \item \textbf{Near-null B$<$C gap:} Mean B-to-C gap $= +0.015$, an order of magnitude below the SBE value of $+0.238$, indicating the swarm barely distinguishes real information from null noise on these events---the social simulation is operating near the susceptibility floor.
    \item \textbf{Cross-model coherence in degradation:} 0/21 SDE events show positive lift across all 3 models simultaneously.
    \item \textbf{Mechanism:} For SDE events, the injected information at $R=30$ conflicts with the swarm's already-formed high-certainty collective prior. Rather than updating productively, agents engage in polarization dynamics that \emph{amplify} the conflict, pushing probability mass away from the correct answer. This is analogous to the ``groupthink inversion'' phenomenon in human deliberation research~\cite{surowiecki2005}.
\end{itemize}

\subsubsection{The Difficulty--Lift Correlation Is an Arithmetic Artifact}
\label{sec:tautology}

The Pearson correlation between $BS_A$ (event difficulty at baseline) and mean cross-model lift is $r = +0.837$ ($r^2 = 0.700$) on the corrected data---it was $r = 0.847$ ($r^2 = 0.717$) before the repair of \S\ref{sec:outage}, and the argument below is insensitive to which value is used---and earlier drafts of this paper read it as evidence that \textbf{event uncertainty explains 71.7\% of the variance in simulation lift}. That reading is not defensible, and we retract it.

The quantity on the vertical axis is $Lift = BS_A - BS_B$, which contains the quantity on the horizontal axis as its minuend. Regressing a difference on one of its own terms produces a large positive correlation by construction, with no contribution from the data. If $BS_B$ were statistically independent of $BS_A$, the expected correlation would be
\begin{equation}
    r_{\text{mech}} = \frac{\sigma_A}{\sqrt{\sigma_A^2 + \sigma_B^2}}
      = \frac{0.380}{\sqrt{0.380^2 + 0.223^2}} = 0.862 ,
\end{equation}
using the observed dispersions $\sigma_A = 0.380$ and $\sigma_B = 0.223$. The observed $r = +0.837$ is thus \emph{below} the value expected under pure arithmetic coupling with an uninformative $BS_B$. The correlation carries no evidence for a difficulty--benefit relationship.

The non-tautological test is the association between baseline difficulty and simulated performance directly, $r(BS_A, BS_B)$, which shares no term between axes. That correlation is $\mathbf{+0.206}$ ($r^2 = 0.042$) on the $30$ event means, $+0.130$ ($r^2 = 0.017$) pooled over the $90$ individual units, and $+0.251$, $-0.078$, $+0.172$ for Llama-3.1-8B, Mistral-7B and Qwen2.5-7B individually. Llama's is the only one that moved with the repair of \S\ref{sec:outage} (from $+0.127$), and at $n=30$ on a single evaluator pass we do not read the change as informative; it is reported so that the appendix matches the released data. Figure~\ref{fig:tautology} plots both panels on the same data. \textbf{Baseline difficulty does not predict how well the swarm forecasts an event.} We therefore withdraw the ``Difficulty--Lift Predictive Rule'' proposed in earlier drafts, including its use as a screening heuristic for allocating simulation budget: on this evidence, a preliminary No-Sim elicitation tells a practitioner nothing about whether simulation will help.

We retain the SBE/SDE partition as a \emph{descriptive} split, because the per-event lift values genuinely are bimodal, but we emphasise that ``SBE'' is largely a relabelling of ``events on which the shared Condition A reference happened to score poorly.'' Since $BS_B$ is roughly constant across strata, selecting events by high $BS_A$ mechanically selects events with high $BS_A - BS_B$. Establishing a genuine structural determinant of simulation benefit requires a covariate measured independently of both conditions---for instance resolution date relative to training cutoff, outcome entropy, or option count---which we identify as the priority for the next campaign.

\subsubsection{Two Units Read End to End: What the Partition Looks Like Up Close}
\label{sec:case_studies}

The partition above is a statement about scores. This subsection reads the two extremes as documents. We take both from the Llama-3.1-8B campaign so that the contrast is internal to one model, one serving stack and one evaluator pass, and because that campaign is the only one whose evidence carries no telemetry-probe content (Table~\ref{tab:probe_share}); the same pair read on Qwen2.5-7B would confound the comparison with the extraction defect of \S\ref{sec:evidence_depth}. Probability vectors below are the mean of the four Evaluator~A reads.

\textbf{C11, the extreme SBE: ``Who will win the 2024 Nobel Peace Prize?'' (ground truth: Nihon Hidankyo).} Condition~A is not merely uncertain but wrong: it puts $0.47$ on \emph{Other} and only $0.12$ on the correct option, for $BS_A = 1.055$ on a five-option question. Condition~B moves the mode onto Nihon Hidankyo at $0.30$ and halves the loss to $BS_B = 0.609$. The transcript, however, does not contain a deliberation that arrives at that answer. Its sixty lines are generic advocacy---the United Nations as a model of cooperation, sustainable development, systemic injustice---and no agent names the eventual laureate on the merits. What the transcript does contain is the round-30 injection, which states that observers note ``an unusual thematic focus from Oslo on the legacy of atomic-bomb survivor advocacy'' and ``suggests a logically stronger probability for Hidankyo-aligned organizations.'' It also contains two lines spoken by an agent whose display name \emph{is} the ground-truth option, an instance of the label leak quantified in \S\ref{sec:evidence_depth}. The lift on this event is therefore best read as retrieval of an injected cue, assisted by a leaked label, rather than as an emergent forecast; and a single RAG+CoT call given the same dossier and the same injection scores $0.638$, which is not distinguishable from the swarm's $0.609$.

\textbf{S6, the extreme SDE: ``Which party wins the most seats in the 2024 Austrian legislative election?'' (ground truth: FP\"O).} Here Condition~A is very nearly right---$0.82$ on FP\"O, $BS_A = 0.053$---and the simulation destroys it: Condition~B places $0.71$ on \emph{Other} and $0.05$ on FP\"O, for $BS_B = 1.448$. S6 is one of the six events that preceded the serving outage, so unlike most of this campaign it carries a dense, fully logged deliberation: $251$ content actions from $185$ distinct agents across both platforms, $91{,}582$ characters. Reading it explains the collapse without recourse to polarization dynamics. The discussion is not about Austria. Agents converge, agreeably and at length, on youth environmental leadership, intergenerational sustainability curricula, frameworks for evaluating intersectional policy, and applications of AI to urban planning; the recurring rhetorical move is ``I completely agree.'' The single on-topic prediction in the sampled evidence forecasts a win for the Greens, which is not one of the four options. A sixty-line sample of that material contains no reason to prefer FP\"O, and the evaluator's mass drifts to the residual category. One CoT call on the same dossier keeps the correct answer at $BS = 0.097$.

\textbf{What the pair shows.} The two events are conventionally read as deliberation helping where the prior is weak and hurting where the prior is strong. The documents support a plainer mechanism: in both cases the evaluator is scoring the topical content of a sixty-line sample, and the sample's relationship to the question is what moves the score. C11 improves because the sample contains a directly relevant injected sentence; S6 collapses because a large volume of fluent, on-consensus, off-topic text displaces a correct prior. This also predicts the null we report elsewhere---that transcript volume does not correlate with lift ($r = +0.016$ over the forty matched-arm units)---since the extractor caps every sample at sixty lines regardless of how much was produced, so volume varies between logging regimes rather than between events. We offer this as an interpretation of two units, not as a tested mechanism; the corresponding experiment is the extractor ablation proposed in \S\ref{sec:evidence_depth}.

\begin{figure*}[!htbp]
\centering
\includegraphics[width=\textwidth]{figD_difficulty_tautology.png}
\caption{The difficulty--lift correlation is arithmetic, not empirical. \textbf{Panel A} reproduces the relationship the pipeline reported: mean simulation lift against baseline difficulty, $r = +0.837$. Because $Lift \equiv BS_A - BS_B$, a correlation of $r = +0.862$ is expected here even if $BS_B$ were statistically independent of $BS_A$; the observed value is below that null. \textbf{Panel B} plots the same 30 events with the shared term removed---baseline difficulty against With-Sim score---giving $r = +0.206$ ($r^2 = 0.042$). Marker colour encodes taxonomy category. Each point is the mean over the three model campaigns.}
\label{fig:tautology}
\end{figure*}

\subsection{Model Interaction Realism: A Critical Evaluation}
\label{sec:interaction_realism}

\begin{figure*}[!htbp]
\centering
\includegraphics[width=\textwidth]{fig5_jsd_convergence.png}
\caption{\textbf{Superseded.} Round-by-round JSD profiles for selected high-lift (left) and low-lift (right) events using Llama-3.1-8B, Condition B, as generated for earlier drafts. This figure should be read only alongside Figure~\ref{fig:jsd_integrity} and Section~\ref{sec:jsd_convergence}: $87.7\%$ of Llama-3.1-8B's checkpoint readings are the constant placeholder $0.758277$, so the apparent differences between the two panels do not reflect measured belief dynamics. It is retained for transparency about what the pipeline emitted, not as evidence of convergence.}
\label{fig:jsd_convergence}
\end{figure*}

\begin{figure*}[!htbp]
\centering
\includegraphics[width=\textwidth]{fig7_directional_matrix.png}
\caption{Directional Accuracy Matrix: Condition B. Each cell shows WIN/50-50/LOSS for the given (model, event) pair. Green = correct direction, orange = wrong direction, neutral = 50-50. Cross-model agreement is highest for SBE events (T2, C11, S9 show WIN across all three models), while SDE events show heterogeneous patterns. Qwen-7B achieves the most robust WIN preservation, losing only 7 events to LOSS in Condition B vs.~12 for Mistral-7B.}
\label{fig:directional_matrix}
\end{figure*}

A critical question for the validity of ECN-BENCH as a research benchmark is whether the agent interaction dynamics exhibit \textbf{realistic social behavior} analogous to human deliberation, or whether they represent superficial textual mimicry without genuine epistemic updating. The empirical data, qualitative analysis of simulation logs, and convergence telemetry permit the following structured assessment. Figure~\ref{fig:jsd_convergence} reproduces the round-by-round convergence profiles on which an earlier draft based this assessment; it is retained only as a record of what the pipeline emitted, and Section~\ref{sec:jsd_convergence} explains why its apparent structure is a placeholder artifact rather than measured belief dynamics.

\subsubsection{Evidence for Genuine Epistemic Interaction}

Several indicators are consistent with selective belief updating in the simulated network:

\textbf{(i) Belief trajectory divergence between Conditions B and C:} For SBE events (T2, S9, C11), the With-Sim and Null Control conditions produce substantially different final probability estimates despite starting from identical persona sets. The B-to-C gap of $+0.238$ for SBE events arises from the fact that the round-30 injection enters the telemetry-probe context differently under Conditions B and C, allowing agents in Condition B to condition on relevant signal while the Null Control introduces irrelevant noise.

\textbf{(ii) Declining JSD on the models with usable telemetry:} Mistral-7B and Qwen2.5-7B---the only two models whose convergence telemetry is not dominated by placeholder values (Section~\ref{sec:jsd_convergence})---show mean inter-round divergence falling from $0.758$ to $0.549$ and $0.575$ respectively over the measured checkpoints, which is directionally consistent with DeGroot sequential belief updating. We attach no weight to the corresponding Llama-3.1-8B trace, which is $87.7\%$ placeholder. This indicator is accordingly weaker than earlier drafts claimed, and applies to the opposite pair of models.

\subsubsection{Action-Log Audit: Model-Dependent Social Interaction Density (Authoritative .db Counts)}

A direct audit of the authoritative SQLite \texttt{trace} tables across 1,380 database files (Tables~\ref{tab:action_density_30event}--\ref{tab:action_breakdown}) reveals a four-tier model-dependent social-action density. Two scope conditions on that audit must be read alongside it, and both are established in Appendix~\ref{sec:outage}. First, the 1,380 files span every campaign including the $N$- and $R$-sweeps, whereas the Brier scores of this paper come from the 30-event campaign alone, which contributes exactly 120 databases per model; restricted to those, the ranking below changes and Table~\ref{tab:action_30event_only} supersedes the counts quoted here for any statement about the scored campaign. Second, $96.7\%$ of Llama-3.1-8B's social actions in the 30-event campaign fall in the six events that preceded a serving-backend failure, so the four-tier ordering is a statement about that model's healthy runs and not about its campaign. We note that earlier drafts of this analysis relied on the \texttt{actions.jsonl} summary files, which systematically undercount organic actions due to a cursor-tracking bug during telemetry-probe rounds (organic actions written during probe rounds are read by the cursor but discarded, as only \texttt{INTERVIEW} actions are logged); the \texttt{.db} counts reported here supersede those earlier figures.

\textbf{Llama-3.1-8B-AWQ---Richest social interaction.} Llama-3.1-8B-AWQ produces the richest social interaction: 27,087 total social actions (40.31/DB), including 8,079 posts, 5,421 comments, 5,993 follows, and 3,010 comment-likes. The high \texttt{follow} and \texttt{comment\_like} volumes indicate that Llama-8B agents build dense peer networks and engage deeply with peers' contributions---the strongest evidence of genuine social deliberation in the benchmark.

\textbf{Mistral-7B---Antagonistic/\allowbreak broadcast profile.} Mistral-7B is the second most interactive (9,364 actions, 35.47/DB) but exhibits a qualitatively different profile: high posting (4,476) and liking (1,946) but disproportionately high disliking (67 vs.\ Llama's 7) and low comment-liking (167 vs.\ Llama's 3,010). This antagonistic/\allowbreak broadcast-style pattern---agents express agreement or disagreement with posts but rarely engage in threaded comment-level discussion---is consistent with conformity and polarization dynamics rather than deliberative consensus-building.

\textbf{Qwen2.5-14B---Analysis-focused interaction.} Qwen2.5-14B (full 30-event campaign, 180 DBs) produces moderate, analysis-focused interaction (2,407 actions, 13.38/DB): 1,215 posts and 585 comments but only 11 follows. Agents write substantive contributions and reply to peers but build almost no social network---a ``lone analyst'' pattern distinct from Llama's network-building and Mistral's broadcast style.

\textbf{Qwen2.5-7B---Near-dormant.} Qwen2.5-7B is the only near-dormant model (723 actions, 2.74/DB), with minimal activity across all action types (592 posts, 108 comments, 16 likes, 1 follow). Under the deployed configuration, its social-interaction channel is effectively dormant.

\textbf{Mechanism implication.} The mechanism of measured Brier lift is therefore strongly model-dependent. For Qwen-7B, lift is attributable to diverse-persona ensemble aggregation plus mid-simulation injection---the social channel is dormant. For Qwen-14B, lift includes a modest analysis-focused contribution component (posts and comments without network-building). For Mistral-7B, lift reflects a mixture of persona-ensemble, antagonistic social interaction (high dislike, low comment-like), and injection effects. For Llama-8B, lift includes a genuine social-deliberation component: agents build peer networks (5,993 follows), engage in threaded discussion (5,421 comments, 3,010 comment-likes), and their posts enter peers' later telemetry-probe timelines. Claims of ``emergent crowd intelligence through social deliberation'' are best supported for Llama-8B; partially supported for Mistral and Qwen-14B; and not supported for Qwen-7B. These are statements about the observed action traces, not about the architectures, for the provenance reasons of Section~\ref{sec:serving_note}. We further restrict the Llama-8B half of that sentence: within the scored 30-event campaign it holds on six events and is false on the other twenty-three, where the swarm produced four social actions in total (Appendix~\ref{sec:outage}). The mechanism statement for Llama-8B should therefore be read as conditional on a working serving backend, which the campaign had for one fifth of its events.

\textbf{Correlational observation.} A four-model correlational observation emerges: interaction VOLUME alone does not predict Brier calibration. Llama-8B (most interaction, 40.31/DB) and Qwen-7B (least interaction, 2.74/DB) achieve nearly identical aggregate Brier scores ($BS_B$ = 0.484 vs.\ 0.481), while Mistral-7B (second-most interaction, 35.47/DB) has the worst Brier ($BS_B$ = 0.531). One candidate distinguishing factor is interaction \emph{quality}: Llama-8B's high \texttt{follow} + \texttt{comment\_like} co-occurs with good calibration, while Mistral's high \texttt{dislike} + low \texttt{comment\_like} co-occurs with poor calibration. We stress that this is a correlational observation over four models with no attempt at control, and that a competing explanation is available for the same pattern: the model with the antagonistic profile is also the oldest by roughly ten months and the only one served at full precision (Section~\ref{sec:serving_note}). Four points cannot separate interaction quality from model generation.

\subsubsection{Evidence for Behavioral Homogenization (Critical Limitation)}

The simulation logs also reveal systematic limitations:

\textbf{(i) Prior anchoring in SDE events:} For events with low $BS_A$ (e.g., C10, C12, C15), the baseline No-Sim evaluation already reflects a strong, confident prior probability. When the simulation begins, agents anchor their responses to this pre-existing LLM prior rather than engaging with the seed context analytically. At $R=30$, the injected intelligence is processed as a \emph{perturbation to a settled consensus} rather than new evidence for consideration, triggering social pressure dynamics that paradoxically increase miscalibration.

\textbf{(ii) Cross-model behavioral convergence (Directional Matrix, Figure~\ref{fig:directional_matrix}):} For SDE events, all three models show identical directional errors (LOSS) despite using different architectures, persona DNA, and serving stacks. This cross-architectural agreement on wrong answers indicates that the simulation is amplifying a shared LLM prior bias rather than discovering event-specific signal. True social deliberation would be expected to show more disagreement between architecturally distinct models on borderline cases.

\subsubsection{Topology and Echo-Chamber Risks}

Earlier drafts stated that the OASIS Twitter topology uses a power-law degree distribution with $\alpha \approx 2.3$. We withdraw that figure: it was not measured, and the run manifest for every campaign records \texttt{degree\_distribution\_descriptor: "oasis-default-unspecified"} with \texttt{clustering\_coefficient: null}. The graphs were built by the framework's default \texttt{generate\_twitter\_agent\_graph} and \texttt{generate\_reddit\_agent\_graph} constructors, whose realised degree distribution we did not instrument. The structural concern nevertheless stands in qualitative form: if the generated graph is heavy-tailed, a small number of highly-connected ``hub'' agents could dominate information flow, so that a position posted early by a hub reaches later agents through multiple social paths and produces \emph{correlated} rather than independent belief updates. However, the .db audit shows the social-action channel is substantial for Llama-8B (27,087 actions) and Mistral-7B (9,364) but near-dormant for Qwen-7B (723). For Llama-8B, the high follow count (5,993) creates dense network connectivity where echo-chamber effects could theoretically amplify hub opinions; for Mistral-7B, the high dislike count (67) suggests early polarization dynamics. For Qwen-7B and Qwen-14B (minimal follow: 1 and 11 respectively), topology effects are moot since agents do not build networks. Empirical validation of echo-chamber dynamics requires graph-level analysis of the follow network, which is pending.

\subsubsection{Implications for Benchmark Validity}

The action-log audit establishes an important, model-dependent scope condition: for Qwen-7B, ECN-BENCH primarily measures persona-ensemble probabilistic aggregation under mid-simulation injection (social channel dormant). For Qwen-14B, it additionally captures analysis-focused posting and commenting without network-building. For Mistral-7B, it captures a persona-ensemble + antagonistic-social + injection mixture. For Llama-8B, it captures genuine social deliberation with peer-network building and threaded engagement. This does not invalidate the measured Brier scores---it indicates that different mechanisms are in play across the campaigns. We say ``across campaigns'' rather than ``by architecture'' deliberately: the campaigns differ in model generation, precision and serving stack as well as architecture (Section~\ref{sec:serving_note}). We note that the difficulty--lift correlation invoked in earlier drafts to delimit where the persona ensemble adds value is an artifact (Appendix~\ref{sec:tautology}) and no longer supports that boundary; the SBE/SDE split remains descriptive only.

To strengthen the social-interaction channel and enable claims of deliberative emergence, we propose the following architectural modifications: (1) \textbf{Forcing a minimum posting rate} for Qwen-7B (near-dormant); (2) \textbf{Network-building incentives} (follow/reply requirements) for Qwen-14B (analysis without network); (3) \textbf{Anti-antagonism mechanisms} (reduce dislike, incentivize comment-level engagement) for Mistral (antagonistic/\allowbreak broadcast pattern), after first re-testing that model at matched generation and precision; (4) \textbf{Reply-to-peer requirements} and \textbf{information drip control} to extend Llama-8B's already-rich deliberation to all event types; (5) \textbf{Diversity penalties} on repeated content across all models.

\subsection{Temporal Data Leakage and Training Contamination: Reinterpreting Simulation Lift}
\label{sec:contamination}

\begin{figure*}[!htbp]
\centering
\includegraphics[width=\textwidth]{fig9_contamination_analysis.png}
\caption{Temporal Data Leakage and Contamination Analysis. Panel A: Condition A Brier score per event colored by contamination likelihood (Red = Contaminated $BS_A < 0.15$, Yellow = Borderline $0.15 \le BS_A < 0.35$, Green = Likely Clean $BS_A \ge 0.35$). 11 of 30 events (36.7\%) fall below the contamination threshold; because Condition A is a single shared record (Section~\ref{sec:condition_a}), this classification is identical across models by construction rather than by agreement. Panel B: $BS_A$ vs.\ Lift scatter with contamination coloring---see Figure~\ref{fig:tautology} and Appendix~\ref{sec:tautology} for why the trend in this panel is arithmetic. Panel C: contamination-corrected lift, which moves from $\approx -0.10$ to $\approx +0.19$--$+0.20$; Figure~\ref{fig:lift_stratified} adds the confidence intervals, all of which span zero, and Appendix~\ref{sec:contamination} explains why thresholding on $BS_A$ makes this reversal partly self-fulfilling.}
\label{fig:contamination_analysis}
\end{figure*}

A crucial methodological insight arising from our empirical campaign is the presence of severe \textbf{temporal data leakage and pre-training contamination}. In benchmarking predictive multi-agent systems, a fundamental prerequisite is that the evaluation events represent genuine \emph{forward-looking uncertainty} relative to the model's knowledge cutoff. If an event outcome was already fixed and heavily represented in the pre-training text corpus or instruction tuning datasets of the LLM, the model does not perform probabilistic forecasting; instead, it performs static factual retrieval (memorized recall).

\subsubsection{Quantifying Contamination Severity}

Tables~\ref{tab:per_event_llama_8b}--\ref{tab:per_event_qwen_7b} reveal that for $11$ of the $30$ events ($36.7\%$), the Condition A baseline Brier score is near zero ($BS_A < 0.15$): C10, C12, C14, C15, C2, C4, S2, S6, T1, T4 and T9. For seven of these, the score is essentially zero ($BS_A \le 0.045$; e.g.\ C12: $0.0005$, C10: $0.0036$, T9: $0.0000$).

We must correct a claim made in earlier drafts, which reported that ``the exact same 11 events are classified as contaminated across all three independently trained model families'' and treated this unanimity as corroborating evidence. That agreement is vacuous. As established in Section~\ref{sec:condition_a}, Condition A was executed once and copied to all three campaigns, so the three models share a single $BS_A$ vector by construction and could not possibly have disagreed. The unanimity is an artifact of the campaign design and carries no evidential weight whatsoever. What remains is the weaker but still meaningful observation that a single direct-elicitation pass answers 11 of 30 events almost perfectly.

That pattern is consistent with these events (such as S2: ``US Senate Control 2024'', C10: ``2024 UK Prime Minister'', T9: ``OpenAI AGI 2025'') having resolved before the training cutoff or having had an overwhelming, unambiguous consensus in pre-training data---in which case Condition A is performing memorised recall rather than forecasting. It is also consistent with the events simply being easy. Distinguishing the two requires the external, date-based test described below; the present data cannot separate them, and we do not claim it does.

\subsubsection{The Simulation vs. Memorization Conflict}

This contamination explains the apparent aggregate "simulation degradation" observed in Section~\ref{sec:cross_model_stats}. When $N=300$ agents engage in social deliberation for $R=60$ rounds on an event where the LLM already possesses $99\%+$ memorized certainty (Condition A $BS_A \approx 0.00$):
\begin{enumerate}
    \item Social interaction forces agents to generate diverse roleplay commentary, explore alternative scenarios, and express nuanced probabilities.
    \item This deliberation process inevitably shifts the swarm's collective probability distribution away from the single deterministic point mass ($1.0$ or $0.0$) toward a moderate distribution (e.g., $0.70$ / $0.30$).
    \item When evaluated against the binary ground truth, this spread of probability mass increases the Brier score from $0.00$ to $0.40$--$0.60$, resulting in a large negative simulation lift ($\Delta \approx -0.40$ to $-0.93$; see Event T1, C15, S6).
\end{enumerate}

Thus, on contaminated events, \textbf{the social simulation is competing against static memorized recall}. Any social deliberation that introduces epistemic uncertainty relative to perfect memorization is mathematically penalized by the Brier metric.

\subsubsection{Baseline Score as a Possible Contamination Proxy}

Earlier drafts advanced this as a reinterpretation of the correlation between $BS_A$ and lift. That correlation is an artifact (Appendix~\ref{sec:tautology}) and cannot support any interpretation, so the argument below rests instead on the unanimous cross-model agreement on which events are near-perfectly answered at baseline, which is a fact about $BS_A$ alone and does not involve the lift. On that basis we propose---as a hypothesis rather than a demonstrated result---that \textbf{$BS_A$ may serve as a proxy for the degree of training data contamination}.
\begin{itemize}
    \item \textbf{Low $BS_A$ ($< 0.15$):} High contamination. The model recalls the answer. Social simulation disrupts static recall, producing negative lift.
    \item \textbf{High $BS_A$ ($\ge 0.35$):} Genuine prediction task (uncontaminated). The model faces true uncertainty. Social deliberation synthesizes disparate signals, producing strong positive lift.
\end{itemize}

\subsubsection{Contamination-Corrected Lift Reversal}

To evaluate the true efficacy of multi-agent social simulation for genuine prediction, we isolate the subset of \textbf{13 Likely Clean Events} ($BS_A \ge 0.35$, where models exhibit genuine baseline uncertainty with mean $BS_A = 0.7408$).

As illustrated in Figure~\ref{fig:contamination_analysis} (Panel C) and Figure~\ref{fig:lift_stratified}, when contaminated events are excluded:
\begin{itemize}
    \item \textbf{Llama-3.1-8B-AWQ:} $BS_A = 0.7408 \to BS_B = 0.5397$. Mean Lift $= +0.2010$, 95\% bootstrap CI $[-0.020, +0.460]$ (from $-0.0993$ on the full dataset).
    \item \textbf{Mistral-7B-Instruct-v0.1:} $BS_A = 0.7408 \to BS_B = 0.5397$. Mean Lift $= +0.2010$, 95\% CI $[-0.025, +0.457]$ (from $-0.1466$).
    \item \textbf{Qwen2.5-7B-GGUF:} $BS_A = 0.7408 \to BS_B = 0.5569$. Mean Lift $= +0.1838$, 95\% CI $[-0.063, +0.447]$ (from $-0.0968$).
\end{itemize}
(The Llama-3.1-8B and Mistral-7B clean-subset means are identical to four decimal places. We verified this is a numerical coincidence: their per-event Condition B scores differ on 12 of the 13 clean events and happen to sum to the same total, $7.0167$. It is not a duplicated row.)

Three qualifications must be attached to this reversal, and together they prevent us from claiming that Claim 1 is confirmed.

\textbf{First, the subset is selected on the dependent variable.} The ``clean'' set is defined by $BS_A \geq 0.35$, and the outcome being averaged is $BS_A - BS_B$. Since $BS_B$ is close to constant across strata (Figure~\ref{fig:lift_stratified}), conditioning on a high $BS_A$ mechanically selects events with a high $BS_A - BS_B$. This is the same arithmetic coupling identified in Appendix~\ref{sec:tautology}, now operating through a selection rule rather than a regression. A positive mean lift on this subset is close to guaranteed and is therefore weak evidence for the hypothesis.

\textbf{Second, none of the three intervals excludes zero.} At $n=13$ events the bootstrap CIs all span zero comfortably, so even taken at face value the reversal is not statistically distinguishable from no effect.

\textbf{Third, the contamination criterion is itself derived from the model's own output} rather than from an external fact about the event. Classifying an event as contaminated because the model scored well on it makes contamination unfalsifiable: any event the model gets right is reclassified out of the sample. A defensible contamination test must use an event-level covariate that is independent of performance---resolution date relative to the model's training cutoff being the obvious candidate, since it is documented for all three model families and requires no access to model outputs. Appendix~\ref{sec:cutoff} carries out that test and supersedes the $BS_A$ threshold used in this subsubsection; the reversal survives the exogenous criterion, but its interpretation changes.

We therefore state the finding as follows: \emph{on the subset of events where the shared No-Sim reference is poorly calibrated, simulation reduces Brier loss by $\approx 0.19$--$0.20$ on average across all three models; this is directionally consistent with the Wisdom of Crowds hypothesis, is not statistically significant at $n=13$, and is substantially inflated by selection on the outcome.} The mechanism of whatever lift is real remains model-dependent, as set out in Appendix~\ref{sec:interaction_realism}.

\begin{figure*}[!htbp]
\centering
\includegraphics[width=\textwidth]{figC_lift_stratified.png}
\caption{Simulation lift by event stratum, with $95\%$ percentile bootstrap confidence intervals ($B=20{,}000$). The sign of the aggregate lift flips from negative to positive as the event set is progressively restricted by the contamination threshold---but the threshold is applied to $BS_A$, which is the minuend of the lift being plotted, so the flip is partly a selection artifact (Appendix~\ref{sec:tautology}). Every interval, in every stratum, includes zero.}
\label{fig:lift_stratified}
\end{figure*}

\subsubsection{Knowledge-Cutoff Stratification: the Negative Lift Lives Entirely in Pre-Release Events}
\label{sec:cutoff}

The two preceding subsubsections classify events as contaminated by thresholding the shared baseline at $BS_A < 0.15$, and state plainly that stratifying on a measured score makes the resulting lift reversal partly self-fulfilling. We now replace that proxy with a criterion exogenous to every measured quantity: an event's resolution date against each model's public release date. A model released before an event resolved cannot have memorised its outcome, whatever its documented data cutoff; release date is therefore a hard upper bound on possible leakage and, unlike $BS_A$, it is fixed before any forecast is made. We use resolution date rather than the dossier's information cutoff because what a model may have memorised is the \emph{outcome}, which becomes public at resolution. Resolution dates come from \texttt{data/events\_raw.json}; the bounds are 2023-09-27 (Mistral-7B-Instruct-v0.1), 2024-07-23 (Llama-3.1-8B-Instruct, whose documented pretraining cutoff is earlier still, December 2023) and 2024-09-19 (Qwen2.5-7B and Qwen2.5-14B).

Ten of the thirty events resolved before at least one model's release (C2, C3, C4, C7, C10, C12, T1, T3, T4, T5); the remaining twenty postdate every model's release and form a \emph{globally clean} subset (C1, C5, C6, C8, C9, C11, C13, C14, C15, S2, S3, S4, S5, S6, S7, S9, T2, T6, T8, T9). Table~\ref{tab:cutoff_lift} and Figure~\ref{fig:cutoff} give the stratified result.

\begin{table*}[!htbp]
\centering
\caption{Simulation lift stratified by whether the event resolved before or after the model's release, with $95\%$ percentile bootstrap intervals ($B = 20{,}000$). Mistral is the oldest model and therefore has only one pre-release event, so its stratification is uninformative and is reported for completeness. \textsuperscript{$\dagger$}Llama's rows are computed on an \emph{instrument-consistent} basis: Conditions A, B and C all scored by \texttt{deepseek-v4-flash-0731}, so the contrast is not confounded by which evaluator read which condition. This is necessary because the 24 events re-simulated on 2026-08-16 (\S\ref{sec:outage}) cannot be scored by the July local evaluator that produced the Mistral and Qwen rows. Under the other two August evaluators the same events give all-event lift $-0.300$ [$-0.471$, $-0.124$] (\texttt{deepseek-r1}) and $+0.046$ [$-0.106$, $+0.198$] (\texttt{gpt-5.6-luna-pro}); the post-release stratum is $-0.343$ and $+0.077$ respectively. The direction of Llama's lift is therefore robust in two of three readings and its magnitude is not; see \S\ref{sec:cutoff}(1).}
\label{tab:cutoff_lift}
\small
\begin{tabular}{llccccl}
\toprule
\textbf{Campaign} & \textbf{Event subset} & $n$ & $BS_A$ & $BS_B$ & $BS_C$ & \textbf{Lift $BS_A - BS_B$ [95\% CI]} \\
\midrule
\multirow{3}{*}{Llama-3.1-8B\textsuperscript{$\dagger$}}
 & all events    & $30$ & $0.283$ & $0.477$ & $0.570$ & $\mathbf{-0.195\ [-0.361, -0.030]}$ \\
 & post-release  & $21$ & $0.321$ & $0.563$ & $0.625$ & $\mathbf{-0.242\ [-0.436, -0.062]}$ \\
 & pre-release   & $9$  & $0.193$ & $0.278$ & $0.443$ & $-0.085\ [-0.400, +0.271]$ \\
\midrule
\multirow{3}{*}{Qwen2.5-7B}
 & all events    & $30$ & $0.385$ & $0.481$ & $0.535$ & $-0.097\ [-0.249, +0.063]$ \\
 & post-release  & $20$ & $0.478$ & $0.497$ & $0.581$ & $-0.019\ [-0.208, +0.195]$ \\
 & pre-release   & $10$ & $0.198$ & $0.451$ & $0.442$ & $\mathbf{-0.253\ [-0.486, -0.028]}$ \\
\midrule
\multirow{2}{*}{Mistral-7B}
 & all events    & $30$ & $0.385$ & $0.531$ & $0.646$ & $-0.147\ [-0.313, +0.028]$ \\
 & post-release  & $29$ & $0.392$ & $0.532$ & $0.662$ & $-0.141\ [-0.314, +0.038]$ \\
\bottomrule
\end{tabular}
\setlength{\tabcolsep}{6pt}
\end{table*}

\begin{figure*}[!htbp]
\centering
\includegraphics[width=\textwidth]{figH_cutoff_lift.png}
\caption{Both panels regenerated 2026-08-16 from the corrected campaign files (\S\ref{sec:outage}). Left: simulation lift for each campaign over all events, over events postdating the model's release, and over events predating it. For Qwen2.5-7B the aggregate negative lift remains attributable to the pre-release stratum, whose interval excludes zero ($-0.25$) while the post-release interval is centred on zero; for Llama-3.1-8B the pattern is inverted, its post-release stratum carrying the larger negative lift ($-0.34$, $n=21$) and its pre-release interval spanning zero ($n=9$). Mistral-7B's pre-release stratum contains a single event and its point is plotted without an interval for that reason. Right: mean Brier score by condition over the twenty events no model could have memorised. Condition A is identical across all three campaigns ($0.478$) because it is a single shared No-Sim record (Section~\ref{sec:condition_a}), restored rather than re-elicited by the repair. Condition B is where the campaigns separate: $0.793$ $[0.625, 0.959]$ for Llama-3.1-8B against $0.473$ and $0.497$ for Mistral-7B and Qwen2.5-7B, an interval that does not overlap either. That gap is confounded with the serving-stack and evaluator substitutions Llama's repaired events carry, and we do not read it as a model comparison.}
\label{fig:cutoff}
\end{figure*}

Four conclusions follow, and the first is the load-bearing one.

\textbf{(1) The contamination reading survives for Qwen2.5-7B but not for Llama-3.1-8B, whose negative lift now sits in the post-release stratum.} This subsubsection was originally built on Llama's contrast between $-0.323$ over its nine pre-release events and $-0.003$ over the twenty-one post-release ones. An intermediate, defective repair on 2026-08-16 appeared to inverse that contrast; that inversion was an artifact and has been withdrawn---the defect and its mechanism are documented in \S\ref{sec:outage}, and the corrected data reported here comes from re-simulating and re-scoring those events properly. On the corrected data both strata carry the same sign: post-release lift is $-0.242$ $[-0.436, -0.062]$, excluding zero, and pre-release lift is $-0.085$ $[-0.400, +0.271]$, spanning it. The stratified contrast that motivated the contamination reading is therefore absent for Llama: its negative lift is not localised to events it could have memorised, and is if anything larger on the events it could not. For Qwen2.5-7B the original contrast is unchanged and still holds: $-0.253$ (excludes zero, pre-release) against $-0.019$ (post-release), as does Mistral's single pre-release event. We read this as a difference between models on valid data rather than as an anomaly: the memorisation account explains Qwen's aggregate lift and does not explain Llama's, and with $E=30$ and the evaluator spread quantified in Table~\ref{tab:cutoff_lift} the design cannot adjudicate why. What the corrected data does remove is any need to invoke an unexplained mechanism: the earlier ``inversion'' was a scoring defect, not a property of the campaign.

\textbf{(2) On events every model had to genuinely forecast, the three campaigns are no longer indistinguishable---Llama is the outlier.} Over the globally clean twenty events, $BS_B$ is $0.548$ (Llama, corrected repair, instrument-consistent pairing), $0.473$ (Mistral) and $0.497$ (Qwen2.5-7B), and lift is $-0.211$ $[-0.405, -0.034]$ (Llama, excluding zero), $+0.005$ $[-0.190, +0.223]$ and $-0.019$ $[-0.208, +0.195]$. The cross-model $BS_B$ spread is $0.075$ here, and Llama is the worst-performing campaign on the clean subset rather than statistically tied with the other two. Two confounds sit on top of this comparison and neither can be removed: Llama's twenty-four re-simulated events ran on a different serving stack from the other campaigns (\S\ref{sec:outage}), and they are scored by a different evaluator than the Mistral and Qwen rows because the July instrument cannot be applied to simulations that did not exist in July. This complicates the decision of Section~\ref{sec:serving_note} to downgrade all three-way model comparisons to descriptive: it is still the right call given the provenance confound, but it can no longer be justified by claiming there is nothing left to explain once memorised events are removed---for Llama there now is, and the design cannot say whether the cause is the model, the serving stack, or the judge.

\textbf{(3) Mistral's susceptibility margin survives; Llama's sign is an evaluator effect, not a data effect.} Over the globally clean subset, $BS_C - BS_B$ is $+0.172$ $[+0.068, +0.291]$ for Mistral-7B (unchanged) and $+0.084$ $[-0.055, +0.226]$ for Qwen2.5-7B. For Llama-3.1-8B the same twenty events give $+0.085$ $[-0.096, +0.283]$ under \texttt{deepseek-v4-flash-0731}, $+0.024$ $[-0.147, +0.197]$ under \texttt{gpt-5.6-luna-pro}, and $-0.153$ $[-0.245, -0.069]$ under \texttt{deepseek-r1}---three readings of one set of transcripts, two centred on the hypothesis' side and one excluding zero against it. An earlier version of this paper reported the negative value alone as a clean-data reversal that the paper could not explain. That reading is withdrawn on two grounds. First, the $-0.198$ figure it rested on came from the defective repair of \S\ref{sec:outage}, in which every unit was scored against an empty evidence field, so Conditions B and C differed only by a letter in the prompt; on the corrected simulations that value does not reproduce. Second, what remains is a disagreement among evaluators applied to identical data, and the three judges differ widely in absolute severity as well (mean Brier $0.32$, $0.48$ and $0.63$), with pairwise directional agreement of only $48$--$79\%$. Susceptibility is a contrast between two simulated conditions and never touches Condition A, so it is insulated from the Condition-A defect---but it is fully exposed to instrument variance, which at $E=30$ is large enough to determine the sign. We therefore report the spread rather than a point estimate, and treat Llama's susceptibility as undetermined by this design.

\textbf{(4) The $BS_A < 0.15$ proxy and the date criterion agree on only $18$--$21$ of $30$ events, and the disagreements are diagnosable.} The proxy flags eleven events (C2, C4, C10, C12, C14, C15, S2, S6, T1, T4, T9). Four of its flags cannot be memorisation: S2 and S6 resolved in November and September 2024, after every model's release---though S6, resolving 2024-09-29, clears the Qwen2.5 bound by only ten days and is the one event whose classification a release-date criterion holds weakly---and C12 and T9 are simply easy given their sanitized dossiers. C12's context is frozen at ``Eve of Finals Game 5'' with Boston leading the series $3$--$1$, and T9 asks whether OpenAI announced AGI in 2025, where the base-rate answer \texttt{NO} requires no knowledge of the outcome. Conversely the proxy misses C3, C7, T3 and T5, which resolved before every relevant release. The two criteria fail in opposite directions: a low $BS_A$ conflates memorisation with an easy dossier, while the date criterion establishes only the \emph{possibility} of memorisation. We report the date criterion as primary because only it is safe to stratify on, and we recommend that any forecasting benchmark built on resolved markets publish the resolution date of every event next to each model's cutoff, so that this stratification is available to readers by construction.

\textbf{The five events the outage never touched are all degenerate.} Intersecting this stratification with the run-integrity audit of Appendix~\ref{sec:outage} isolates the units that were simultaneously post-release and served by a healthy backend in the \emph{original} July campaign: five events (S2, S3, S4, S5, S6), ten simulated units. Every one of those ten returns an exactly uniform probability vector---$\{0.5, 0.5\}$ for the four binary events and $\{0.25, 0.25, 0.25, 0.25\}$ for S6---under Condition B and under Condition C alike, giving $BS_B = BS_C = 0.550$ against $BS_A = 0.344$. Within those five, degenerate consensus is not a tendency, it is the entire outcome. Since the corrected repair the reference campaign is no longer confined to them---the other twenty-four events now carry valid re-simulated transcripts---but they remain the only Llama events whose data was produced end-to-end on the original local serving stack, and so the only ones free of the substitution confound.

\textbf{A separate design defect affects Condition B on seven events, in every campaign.} The step-30 injection bank is keyed by event identifier and was authored against the stale catalogue described in \S\ref{sec:outage}. For the seven identifiers where the two catalogues disagree (C6, T1, T3, T4, T5, T6, T8), Condition B's ``relevant update'' is therefore topically unrelated to the question the simulation actually ran: the 2025 Papal Conclave event receives a core-PCE inflation bulletin, the Taiwanese presidential election receives a Bitcoin treasury announcement, the Irish general election receives an OpenAI product note. This is not a consequence of the outage or of either repair---it is present in the original Mistral-7B and Qwen2.5-7B campaigns as well, both of which simulated the correct questions under the incorrect injections. On those seven events Condition B is thus effectively a second null control, and the data behave accordingly: $BS_C - BS_B$ is $-0.023$ $[-0.181, +0.126]$ on the seven versus $+0.157$ $[+0.034, +0.287]$ on the other twenty-three for Mistral-7B, and $-0.187$ versus $+0.126$ for Qwen2.5-7B. The susceptibility effect this paper reports exists only where the injection is on-topic. We leave the injections unchanged rather than re-authoring them for the re-simulated Llama events, which would have given one campaign better stimuli than the others, and record the seven events as a known limitation of the instrument.

\subsubsection{Normalized Brier Skill Score (BSS) Evaluation}

\begin{figure*}[!htbp]
\centering
\includegraphics[width=\textwidth]{fig10_bss_comparison.png}
\caption{Normalized Brier Skill Score (BSS) Comparison across the 19 Filtered Events ($BS_A \ge 0.15$). $BSS = 1.0$ indicates perfect forecasting, $BSS = 0.0$ indicates uniform random guessing. Across all three model architectures, Condition A direct elicitation yields marginal skill over random chance ($BSS_A = +0.0924$), whereas Condition B social simulation elevates forecasting skill to $BSS_B \approx +0.2264 \to +0.2476$ (a net $+0.134 \to +0.155$ skill gain).}
\label{fig:bss_comparison}
\end{figure*}

\begin{figure*}[!htbp]
\centering
\includegraphics[width=\textwidth]{fig11_bss_by_format.png}
\caption{BSS Skill Lift by Question Format: Boolean ($K=2$, $n=7$) vs.\ MCQ ($K > 2$, $n=12$). Boolean questions exhibit dramatic simulation skill gains ($\Delta BSS = +0.158 \to +0.261$). MCQ questions maintain positive forecasting skill over random guessing ($BSS_B \approx +0.17 \to +0.25$), though probability dispersion across multiple options reduces net lift compared to binary choices.}
\label{fig:bss_by_format}
\end{figure*}

To evaluate forecasting capability without option-count distortion, we convert raw Brier scores into the normalized \textbf{Brier Skill Score (BSS)}. On the 19 filtered events:
\begin{itemize}
    \item \textbf{Llama-3.1-8B-AWQ:} Baseline skill $BSS_A = +0.0924 \to$ Condition B $BSS_B = \mathbf{+0.2264}$ (Skill Lift $= \mathbf{+0.1340}$).
    \item \textbf{Mistral-7B-Instruct-v0.1:} Baseline skill $BSS_A = +0.0924 \to$ Condition B $BSS_B = \mathbf{+0.2476}$ (Skill Lift $= \mathbf{+0.1552}$).
    \item \textbf{Qwen2.5-7B-GGUF:} Baseline skill $BSS_A = +0.0924 \to$ Condition B $BSS_B = \mathbf{+0.2309}$ (Skill Lift $= \mathbf{+0.1385}$).
\end{itemize}

Figure~\ref{fig:bss_comparison} illustrates this BSS elevation. Direct single-pass elicitation (Condition A) achieves less than $10\%$ skill above uniform random guessing ($BSS_A = +0.0924$). In contrast, multi-agent social interaction (Condition B) more than doubles forecasting skill, elevating $BSS_B$ to $+0.2264 \to +0.2476$.

\subsubsection{Question Format Analysis: Boolean vs.\ MCQ}

Decomposing the 19 filtered events into \textbf{Boolean} ($K=2$, $n=7$) and \textbf{Multiple Choice (MCQ)} ($K > 2$, $n=12$) reveals significant structural differences in swarm performance (Figure~\ref{fig:bss_by_format}):

\begin{itemize}
    \item \textbf{Boolean Questions ($K=2$):} Social simulation produces dramatic skill gains across all models. BSS increases from $BSS_A = -0.1329$ (sub-random baseline) to $BSS_B = +0.1829 \to +0.3886$. Mistral-7B achieves the highest Boolean skill lift ($\Delta BSS = +0.5214$). Binary consensus formation enables swarms to efficiently concentrate probability mass onto the realized outcome.
    \item \textbf{MCQ Questions ($K > 2$):} For multi-option questions (up to $K=9$), Condition B maintains solid skill above random guessing ($BSS_B = +0.1654 \to +0.2517$). However, agent deliberation causes \emph{probability dispersion} across secondary choices, moderating net skill lift relative to Condition A ($\Delta BSS = -0.058 \to +0.028$). Crucially, Condition B consistently outperforms Null Control Condition C across all MCQ events (Susceptibility Margin $= +0.061 \to +0.193$), confirming that the swarm actively integrates real-world signal over noise even under multi-option dispersion.
\end{itemize}

This finding establishes an essential methodological mandate for future agentic benchmark design: benchmarks evaluating predictive multi-agent systems must implement strict temporal filtering or post-cutoff event generation to prevent memorized recall from corrupting crowd intelligence evaluation.

\subsection{Reproducibility and Data Audit}
\label{sec:data_audit}

All aggregate statistics reported in this paper (per-model mean $BS_A$/$BS_B$/$BS_C$, directional accuracy, Diebold-Mariano statistics, the 11/6/13 contaminated/borderline/clean event split, and the Brier Skill Score values) are reproduced from the programmatic \texttt{summary.json} of the three completed campaigns, in which all 270 (model, event, condition) units are recorded as having completed with no errors and no evaluator fallbacks. Appendix~\ref{sec:outage} shows that those recorded flags are not evidence of a successful run: 69 of the 270 executed against an unreachable model endpoint and are flagged clean regardless. For this revision we additionally recomputed every reported quantity independently from the raw \texttt{event\_results.json} records; the script that does so and that generates Figures~\ref{fig:susceptibility}, \ref{fig:tautology}, \ref{fig:jsd_integrity}, \ref{fig:degenerate} and \ref{fig:lift_stratified} is released with the code, as is the run-integrity script behind Figures~\ref{fig:outage}, \ref{fig:evidence_depth} and \ref{fig:cutoff}, so that each figure can be traced to the per-unit records.

That recomputation surfaced three discrepancies with earlier drafts, which we record here in full.

\textbf{(1) Per-event table errors.} A full cell-by-cell comparison---not the spot-check described in earlier drafts---found \textbf{18 incorrect cells}, all in the Mistral-7B table and all confined to its nine alphabetically-first C-prefixed events (C1, C10, C11, C12, C13, C14, C15, C2, C3). Errors reach $0.93$ Brier (C2: reported $0.2000$, actual $1.1338$), not the ``$\approx 0.18$ maximum'' previously stated. The Llama-3.1-8B and Qwen2.5-7B tables match the raw records exactly. Because the aggregate means were always computed programmatically, no mean, test statistic or conclusion in the paper depended on the erroneous cells; Table~\ref{tab:per_event_mistral_7b} has been regenerated from source.

\textbf{(2) Metric conflation in the susceptibility column.} The $\Delta_{B-C}$ value previously reported for Qwen2.5-7B ($-0.044$) was read from \texttt{content\_susceptibility.delta.B\_minus\_C} (a mean yes-probability difference) while the other two rows were read from the Brier margin, producing a spurious sign reversal. Corrected in Table~\ref{tab:cross_model_stats} and Section~\ref{sec:cross_model_stats}.

\textbf{(3) Condition A provenance.} Condition A is shared across models, and earlier drafts justified this as an intended consequence of deterministic decoding (temperature~0, seed~42). That justification is incorrect: determinism makes a \emph{single} model reproducible, but cannot make three different model families emit byte-identical probability vectors on 30 events, which is what the records show. Condition A was generated once and reused. Section~\ref{sec:condition_a} treats this as a design defect and traces its consequences. The replication of Section~\ref{sec:evaluator_replication} turns this defect to diagnostic use: because the Condition A input is byte-identical across campaigns, re-evaluating each campaign separately constitutes a triplicate re-read of the same text, which is how the instrument's reproducibility ($5/30$) was measured.

\subsubsection{Two Evaluator Passes, Released Separately}
\label{sec:data_audit_two_passes}

The replication of Section~\ref{sec:evaluator_replication} adds a second scored pass over the same simulation artifacts, and the two must not be mixed. Each campaign directory therefore carries both, each internally consistent in $90/90$ units:

\begin{itemize}
    \item \texttt{event\_results.json} --- Dataset J, the July pass (DeepSeek-R1-14B). Reproduces that campaign's \texttt{summary.json} exactly, and is the basis of every number in Sections~\ref{sec:full_campaign}--\ref{sec:degenerate}.
    \item \texttt{event\_results\_evaluatorB\_openrouter.json} --- Dataset A, the August pass (\texttt{deepseek-v4-flash-0731}), with all probability-derived metrics recomputed from that pass's own vectors.
\end{itemize}

We record the hazard that produced this arrangement, since it is a general one. The re-evaluation utility overwrites \texttt{probabilities}, \texttt{mcq\_dimensions}, \texttt{validated\_scales} and \texttt{micro\_epistemic\_mapping}, but does \emph{not} recompute \texttt{brier}, \texttt{rps}, \texttt{directional\_accuracy} or the \texttt{calibration\_*} fields, which are written once by the protocol runner. Running it in place therefore produced files in which the stored Brier described a probability vector the record no longer contained---in $82$, $86$ and $87$ of $90$ units for Llama, Mistral and Qwen respectively---while remaining superficially well-formed and passing the pipeline's own completeness checks. Any analysis reading \texttt{brier} from those files would have silently mixed two instruments. The reconciliation utility (\texttt{separate\_evaluator\_datasets.py}) rebuilds both datasets from the preserved snapshots, reports the consistency count for each, and retains the inconsistent intermediate file for audit rather than deleting it. We recommend that any benchmark storing derived metrics alongside the predictions they derive from either recompute them atomically with the prediction or validate their consistency on load; a stored metric that cannot be regenerated from its own record is not auditable.

A second, repaired hazard is recorded here for completeness, because the shape of it generalises. Dataset A was first generated against a stale copy of the event catalogue shipped under \texttt{ECN-HPC-DEPLOY/data/}, in which the identifiers C6, T1, T3, T4, T5, T6 and T8 carry different events; the \texttt{micro\_epistemic\_mapping} field of the corresponding $21$ units per campaign therefore answered micro-questions belonging to the wrong event, while the probability vectors---built from each record's own question, evidence and option list---remained correct. Nothing about the affected files was malformed: two catalogues with overlapping identifiers and disjoint contents produce records that validate cleanly and are wrong only by cross-reference. It was caught by comparing each unit's stored micro-question ids against the catalogue, which is now a released check (\texttt{find\_contaminated\_micro\_events.py}); the $63$ affected units were re-evaluated against the authoritative \texttt{data/events\_raw.json}, the checker reports zero remaining mismatches, and the catalogue path is corrected in the released scripts. The generalisable lesson is that an event catalogue keyed by short identifiers needs a content hash or version recorded in each run manifest; identifier collision across catalogue versions is otherwise undetectable from the outputs alone.

\subsubsection{Persona Pool Audit: the 4-Axis DNA Is Not Realised as Specified}
\label{sec:persona_audit}

Appendices~\ref{sec:optimal_bounds} and~\ref{sec:cross_model} justify $N=300$ partly on cognitive-diversity grounds, stating that at that population each of the four persona-DNA axes (Worldview, Motivation, Style, Biases) is represented by $\geq 75$ distinct agents. That claim was never measured. We measure it here by counting the persona records the pipeline actually wrote: all $270$ \texttt{reddit\_profiles.json} files, $300$ agents each, $81{,}000$ agent records in total. Table~\ref{tab:persona_audit} reports the result, and it does not support the claim for two of the three campaigns.

\begin{table*}[!htbp]
\centering
\caption{Persona pool as written to disk, per campaign. ``Distinct personas'' counts unique persona strings among the 300 agents of one unit; the boilerplate bio is the event-independent string \texttt{"[profession] based in [country], engaging on emerging developments."}; template compliance is the fraction of agents whose persona string emits all four labelled axes. Persona sets are fixed per event: Conditions A, B and C of the same event share one pool, so the per-condition figures are identical.}
\label{tab:persona_audit}
\footnotesize
\setlength{\tabcolsep}{3pt}
\resizebox{\textwidth}{!}{%
\begin{tabular}{lccccc}
\toprule
\textbf{Campaign} & \textbf{Distinct personas / 300} & \textbf{min / median / max} & \textbf{Boilerplate bio} & \textbf{4-axis compliance} & \textbf{Units $\leq$ 20 distinct} \\
\midrule
Mistral-7B-Instruct-v0.1 & $295.6$ & $263 / 298 / 300$ & $0.7\%$  & $97.7\%$ & $0/90$ \\
Llama-3.1-8B-AWQ         & $71.6$  & $11 / 11 / 271$    & $79.0\%$ & $16.2\%$ & $69/90$ \\
Qwen2.5-7B-GGUF          & $38.0$  & $19 / 40 / 57$     & $92.0\%$ & $1.7\%$  & $9/90$ \\
\bottomrule
\end{tabular}%
}
\setlength{\tabcolsep}{6pt}
\end{table*}

Three observations follow.

\textbf{(1) Only Mistral realises the specified design.} Its pools carry a median of $298$ distinct personas out of $300$ and emit the four labelled axes in $97.7\%$ of agent records. Whatever else is confounded about that campaign (Section~\ref{sec:serving_note}), its Stage~1 persona extraction worked.

\textbf{(2) Qwen2.5-7B generates roughly 38 event-specific stakeholders and fills the remaining $\sim$260 agent slots with a boilerplate template.} Its distinct-persona count never exceeds $57$. It also does not follow the four-axis template: it writes \texttt{- Biased against [X]} in place of \texttt{- Biases: [X]}, so a strict parse of the Biases axis succeeds on $1.7\%$ of its agents. The persona ensemble that Appendix~\ref{sec:interaction_realism} identifies as this model's operative mechanism is therefore an ensemble of about $38$ archetypes, not $300$.

\textbf{(3) In 23 of 30 events, the Llama-3.1-8B pool is a static fallback library of 11 archetypes, identical across events.} Those 69 units decompose as $5 \times 45$ agents on five generic worldview archetypes (\emph{grassroots activist}, \emph{techno-optimist}, \emph{skeptical contrarian}, \emph{traditionalist moralist}, \emph{globalist bureaucrat}), $30$ agents on \texttt{"Worldview: Pragmatic and data-driven"}, and $5 \times 9$ agents on \texttt{"Pragmatic institutionalist; filters truth through [profession] experience"}. Not one of the 11 refers to the event under forecast. The remaining $21$ units carry $270$--$271$ distinct, event-specific personas. The split is not random and not event-driven; Appendix~\ref{sec:outage} identifies its cause.

The generalisable point is that persona diversity is an \emph{output} of the model under test---Stage~1 extraction is an LLM call---and therefore has to be measured per campaign rather than assumed from $N$. A benchmark that scales agent population as its primary diversity knob should report realised distinct-persona counts alongside $N$, because the two came apart here by a factor of $27$.

\subsubsection{A Mid-Campaign Backend Outage Voids 23 of the 30 Llama-3.1-8B Events}
\label{sec:outage}

\begin{figure*}[!htbp]
\centering
\includegraphics[width=\textwidth]{figF_llama_outage.png}
\caption{The Llama-3.1-8B campaign's Condition B units in wall-clock execution order. The serving backend fails during the seventh event (S7) and never recovers: realised persona diversity drops from $271$ to the 11-archetype fallback, social actions in the authoritative \texttt{.db} trace drop from $448$--$539$ per unit to exactly $4$, every one of the $1{,}000$ telemetry probes fails with \texttt{Connection error}, and the \texttt{interview} rows through which the forecast is extracted fall to zero. All 23 post-break events record byte-identical failure counts, which is the signature of a dead endpoint rather than of noise. Every one of these units is nevertheless recorded as \texttt{simulation\_status: "completed"} with \texttt{error: null}.}
\label{fig:outage}
\end{figure*}

Ordering the Llama-3.1-8B units by wall-clock start time reveals a single, sharp regression, shown in Figure~\ref{fig:outage} and summarised in Table~\ref{tab:outage}. The first six events (S2, S3, S4, S5, C3, S6, in execution order) ran normally. During the seventh (S7) the vLLM endpoint serving the swarm became unreachable. For the remaining 23 events every unit records exactly $1{,}000$ \texttt{Telemetry probe failed \ldots Connection error} warnings, exactly $340$ \texttt{Model error} lines, exactly $4$ social actions across both platform databases, zero \texttt{interview} rows, and a $2.93$\,MB log---the same numbers, unit after unit.

\begin{table*}[!htbp]
\centering
\caption{Llama-3.1-8B campaign, per-unit run health by execution phase. Social actions and \texttt{interview} rows are counted from the authoritative SQLite \texttt{trace} tables; failure counts from \texttt{simulation.log}. The \texttt{interview} action is the telemetry probe that extracts each agent's belief state, i.e.\ the channel the Condition B and C forecasts are read from.}
\label{tab:outage}
\footnotesize
\setlength{\tabcolsep}{4pt}
\resizebox{\textwidth}{!}{%
\begin{tabular}{llccccc}
\toprule
\textbf{Phase} & \textbf{Events} & \textbf{Distinct personas} & \textbf{Social actions/unit} & \textbf{Probe failures/unit} & \textbf{\texttt{interview} rows} & \textbf{Log size} \\
\midrule
Healthy & S2, S3, S4, S5, C3, S6 (6) & $270$--$271$ & $448$--$539$ & $0$       & $12{,}000$ (all) & $0.25$\,MB \\
Failing & S7 (1)                     & $271$        & $14$         & $1{,}000$ & $0$              & $2.88$\,MB \\
Dead    & S9 $\to$ C15 (23)          & $11$         & $4$          & $1{,}000$ & $0$              & $2.93$\,MB \\
\bottomrule
\end{tabular}%
}
\setlength{\tabcolsep}{6pt}
\end{table*}

Three consequences bear directly on results reported earlier in this paper.

\textbf{(1) It is the cause of the JSD placeholder already reported.} Section~\ref{sec:jsd_convergence} records that $87.7\%$ of Llama-3.1-8B's round-level Jensen--Shannon-divergence readings are the constant $0.758277$ and withdraws the convergence result on that basis. The outage explains why: with no successful telemetry probe there is no belief distribution to diverge, so the field retains its initialised value. The symptom and its cause are the same event.

\textbf{(2) The interaction-density ranking of Appendix~\ref{sec:interaction_realism} is stated for the un-repaired campaign and has not been recomputed.} The figures quoted there ($27{,}087$ actions for Llama at $40.31$/DB, $9{,}364$ for Mistral, $723$ for Qwen-7B) aggregate $672$, $264$ and $264$ database files across all campaigns including the $N$- and $R$-sweeps, none of which were touched by the 2026-08-16 repair; that comparison is left as reported and should be treated as describing the pre-repair state until the sweep databases are re-scanned. Table~\ref{tab:action_30event_only}, restricted to the $120$ databases of the 30-event campaign that produces every Brier score in this paper, \emph{has} been rescanned against the post-repair \texttt{.db} files, and the ranking it showed pre-repair no longer holds: Llama now leads Mistral in every social-action category (e.g.\ $9{,}074$ vs.\ $3{,}545$ \texttt{create\_post}), and \texttt{interview} probes are tied across all three models at $60{,}000$ (previously $12{,}000$ for Llama alone, since only its $12$ healthy units contributed a full telemetry read). This inverts the previous within-campaign picture, in which Mistral posted more than twice as often as Llama and $96.7\%$ of Llama's recorded actions came from just the six pre-outage events---that asymmetry was a direct readout of the outage itself, not a property of the model, and dissolves once the affected events carry real telemetry.

\begin{table*}[!htbp]
\centering
\caption{Action counts from the \texttt{trace} tables of the 30-event campaign only ($120$ databases per model: 30 events $\times$ 2 simulated conditions $\times$ 2 platforms). Compare Table~\ref{tab:action_density_30event}, which aggregates across all campaigns including the parameter sweeps. \texttt{interview} is the telemetry probe; its deficit for Llama is the outage of Appendix~\ref{sec:outage}.}
\label{tab:action_30event_only}
\small
\begin{tabular}{lcccccc}
\toprule
\textbf{Campaign} & \texttt{create\_post} & \texttt{create\_comment} & \texttt{follow} & \texttt{like\_post} & \texttt{like\_comment} & \texttt{interview} \\
\midrule
Llama-3.1-8B-AWQ         & $\mathbf{9{,}074}$ & $\mathbf{6{,}454}$ & $\mathbf{5{,}046}$ & $\mathbf{4{,}046}$ & $\mathbf{2{,}044}$ & $60{,}000$ \\
Mistral-7B-Instruct-v0.1 & $3{,}545$ & $1{,}648$ & $576$ & $1{,}618$ & $139$ & $60{,}000$ \\
Qwen2.5-7B-GGUF          & $281$ & $104$ & $1$ & $14$ & $1$ & $60{,}000$ \\
\bottomrule
\end{tabular}
\setlength{\tabcolsep}{6pt}
\end{table*}

\textbf{(3) The failure is invisible in the result records.} Every one of the 69 affected units carries \texttt{simulation\_status: "completed"}, \texttt{full\_simulation\_completed: true}, \texttt{simulation\_executed: true}, \texttt{error: null}, \texttt{strict\_contract: true}, \texttt{evaluator\_fallback\_used: false} and \texttt{evaluator\_reliability\_status: "complete"}. The statement made in earlier drafts---that all $270$ units completed ``with no errors and no evaluator fallbacks''---is true of those fields and false of the runs, and is corrected wherever it appears in this revision. The generalisable lesson is that a pipeline whose completion flags are set by the orchestrator rather than derived from the artifacts will report a dead backend as a successful experiment; the auditable form of the same check is to require a minimum count of successful telemetry probes per unit before a forecast is scored, which would have failed loudly on all 69.

\subsubsection{What the Evaluator Actually Read}
\label{sec:evidence_depth}

Section~\ref{sec:evaluator_replication} treats the evaluator as an instrument and measures its reproducibility. This subsection measures its \emph{input}. Each scored record carries the verbatim text the evaluator was given, in \texttt{evidence\_text}: a tail of the simulation log followed by a block headed \texttt{Representative Social Media Actions (Sample)}. Parsing that block gives, per unit, how many action lines the evaluator saw, from how many distinct agents, and how many distinct utterances they contained. Table~\ref{tab:evidence_depth} reports the distribution and Figure~\ref{fig:evidence_depth} plots every unit.

\begin{table*}[!htbp]
\centering
\caption{Deliberation content supplied to the evaluator, over the $60$ simulated (Condition B and C) units of each campaign. ``Distinct utterances'' counts unique message texts among the action lines shown.}
\label{tab:evidence_depth}
\footnotesize
\setlength{\tabcolsep}{3pt}
\resizebox{\textwidth}{!}{%
\begin{tabular}{lccccc}
\toprule
\textbf{Campaign} & \textbf{\texttt{evidence\_text}} & \textbf{Action lines} & \textbf{Distinct speakers} & \textbf{Distinct utterances} & \textbf{Units with $\leq 1$ speaker} \\
 & (median chars) & (median) & (median) & (median) & \\
\midrule
Mistral-7B-Instruct-v0.1 & $12{,}781$ & $60$ & $46$ & $41.5$ & $0/60$ \\
Qwen2.5-7B-GGUF          & $7{,}034$  & $60$ & $54$ & $\mathbf{8}$ & $0/60$ \\
Llama-3.1-8B (6 healthy events) & ${\sim}19{,}000$ & $60$ & $38$--$43$ & $51$--$55$ & $0/12$ \\
Llama-3.1-8B (23 outage events) & $3{,}045$ & $\mathbf{6}$ & $\mathbf{1}$ & $\mathbf{2}$ & $46/46$ \\
\bottomrule
\end{tabular}%
}
\setlength{\tabcolsep}{6pt}
\end{table*}

\begin{figure*}[!htbp]
\centering
\includegraphics[width=\textwidth]{figG_evidence_depth.png}
\caption{Left: the deliberation content each simulated unit's evaluator was given. Three regimes are visible---Mistral in the diverse corner ($33$--$53$ speakers, $33$--$50$ distinct texts), Qwen2.5-7B with many speakers but few distinct texts, and Llama split between its healthy units and $47$ units carrying at most one speaker. Right: the weighted rubric score the evaluator returned from that input. Within the Llama campaign, units whose evidence is a single-speaker error dump score $0.495$ against $0.508$ for units carrying a multi-speaker transcript---a difference of $0.013$ on a unit interval.}
\label{fig:evidence_depth}
\end{figure*}

\textbf{The outage units are scored from an error log.} For a representative affected unit (\texttt{C10\_B\_r1}, ``Who will be the next UK Prime Minister after the 2024 election?''), \texttt{evidence\_text} consists of a log tail of \texttt{Telemetry probe failed \ldots Connection error} lines reporting \texttt{Total actions: 3}, followed by six action lines from a single agent: four are the forecast question repeated verbatim, and two are the round-30 injection text. From that input the evaluator returned $P(\text{Starmer}) = 0.45$ and a complete seven-dimension rubric scoring \texttt{herd\_effect} and \texttt{convergence} as \texttt{very\_high}, \texttt{deliberation\_quality} as \texttt{high}, and a weighted rubric score of $0.6125$. There was no deliberation of any kind in its input.

\textbf{Qwen2.5-7B's evidence is telemetry scaffolding, not discussion.} Its units show many speakers but few distinct texts because the sampled lines are probe artifacts: in \texttt{C10\_B\_r1}, $46$ of the $60$ lines shown are the literal string \texttt{"Belief check: YES Probability = None"}---a probe that returned no value---and only one line is a substantive agent utterance, which is written in Italian in violation of the \texttt{Global Language Rule} carried in that agent's own persona. This is direct evidence for the behavioural-homogenization limitation that Appendix~\ref{sec:interaction_realism} states qualitatively, and it locates the effect in the extraction channel rather than in agent behaviour.

\textbf{The sample is capped at sixty lines, and the probe backfills it.} The block is produced by \texttt{build\_evidence\_text}, which calls \texttt{\_sample\_representative\_actions(\ldots, max\_actions=60)}. That function assigns an engagement score to each logged action---$1.0$ for a post, $0.8$ for a comment, $0.1$ for a like---sorts descending, keeps the top sixty and re-sorts them by round. Two consequences follow, and neither is a property of the swarm. First, the cap binds everywhere: the median action-line count is exactly $60$ in all four campaigns, whether the unit's database holds two hundred content items or five, so the quantity of deliberation the evaluator sees is a constant of the extractor and cannot covary with anything. Across the released data those sixty lines span a median of $16$--$17$ distinct rounds out of $61$, concentrated in the opening third; the closing rounds are largely absent from what is scored. Second, \texttt{TELEMETRY\_PROBE} records are rendered into the same block as \texttt{Belief check: YES Probability = \ldots} but receive no engagement score, so they sit at zero and enter the sample only when a unit has fewer than sixty scored content actions. On a sparse campaign they therefore fill the remainder.

\begin{table}[!htbp]
\centering
\caption{Share of the evaluator's sixty evidence lines that are telemetry probes rather than agent utterances, over the $60$ simulated units of each campaign in the currently released result files. Unlike Table~\ref{tab:evidence_depth}, which describes the archived pre-repair snapshot, these figures are measured on the data this paper reports.}
\label{tab:probe_share}
\small
\begin{tabular}{lccc}
\toprule
\textbf{Campaign} & \textbf{Probe share} & \textbf{Probe share} & \textbf{Distinct} \\
 & (mean) & (median) & \textbf{utterances} \\
\midrule
Llama-3.1-8B             & $0.0\%$          & $0.0\%$          & $56$ \\
Mistral-7B-Instruct-v0.1 & $0.0\%$          & $0.0\%$          & $41.5$ \\
Qwen2.5-7B-GGUF          & $\mathbf{86.0\%}$ & $\mathbf{86.7\%}$ & $8$ \\
Qwen2.5-14B-AWQ          & $\mathbf{72.0\%}$ & $\mathbf{90.0\%}$ & $6$ \\
\bottomrule
\end{tabular}
\end{table}

\textbf{The fourth model is affected as badly as the third, and it was never audited.} Table~\ref{tab:probe_share} extends the measurement to Qwen2.5-14B, which \S\ref{sec:fourth_model_status} previously excluded from this appendix. That exclusion was wrong in scope: it was justified on the grounds that the model's \texttt{.db} traces and persona files never left the cluster, which is true and does bar the audits that read those artifacts, but the present measurement reads \texttt{evidence\_text} out of the result records, and those are on disk. Read that way, Qwen2.5-14B sits in the same regime as Qwen2.5-7B and slightly deeper into it: a median of $90\%$ of its evidence lines are probe readouts and a median of $6$ distinct utterances survive in a sixty-line sample. This matters more than the corresponding statement about Qwen2.5-7B, because Qwen2.5-14B is the campaign that posts the best forecasting numbers in this paper---the highest Condition-B directional accuracy and, with Qwen2.5-7B, one of only two campaigns whose $\Delta$Brier interval excludes zero. The evidence from which those numbers were produced consists mostly of the simulator's own per-agent probability estimates, which is to say that the measurement channel is feeding the channel being scored.

\textbf{What that does and does not license.} It does not license the conclusion that the leak produces the result. The probe share is very nearly a campaign-level constant---$0\%$, $0\%$, $86\%$, $72\%$---so it is confounded with model identity, generation and serving stack together, and the design contains no within-campaign variation in it to exploit. The decisive counter-example is in the same table: Mistral-7B reads a transcript with no probe content at all and attains a Condition-B directional accuracy of $0.850$, identical to Qwen2.5-7B's. Pooling units across campaigns gives $0.892$ against $0.756$ in favour of the leaked units, but that comparison is between campaigns wearing a per-unit disguise and we decline to read it causally. The defensible statement is narrower and still serious: for half of the four designated models, the input to the scored forecast is dominated by an artifact of the instrumentation, and the present design cannot separate that from whatever those campaigns' agents contributed. Isolating it requires holding the campaign fixed and varying the extractor---re-running \texttt{\_sample\_representative\_actions} with telemetry records excluded and re-scoring the same units---which costs no new simulation and which we flag as the immediate next experiment.

\textbf{A speaker name can carry the answer.} A separate and smaller channel runs through the display names. Personas are occasionally instantiated with names drawn from the event's own option set, and the sampled lines carry the speaker name verbatim, so the evaluator can encounter \texttt{[R13][reddit] Keir Starmer: \ldots} while being asked who will be the next UK Prime Minister. Over the $240$ simulated units, $31$ contain a speaker named after some option and $\mathbf{20}$ contain a speaker named after the \emph{ground-truth} option, spread over six events (C10 in all four campaigns, C12, T4, C11, S6, C8). Within the twenty-event contamination-clean set the count is five units. The affected fraction is small and we do not restate any result on account of it, but it is a label leak in the same extraction channel as the previous two paragraphs, it is invisible to every completion flag the pipeline sets, and a persona generator that may not sample option strings is a cheap gate against it.

\textbf{The rubric does not discriminate.} Comparing within the Llama campaign, so that the contrast is not confounded by which model produced the unit, the $47$ units whose evidence contains at most one speaker receive a mean weighted rubric score of $0.495$; the $13$ units whose evidence contains more than one speaker---the twelve healthy units at $38$--$43$ speakers, plus the transitional S7 unit---receive $0.508$. Per dimension the gap is of the same order (\texttt{herd\_effect} $0.535$ vs $0.510$; \texttt{convergence} $0.526$ vs $0.508$; \texttt{deliberation\_quality} $0.465$ vs $0.508$). An instrument that assigns essentially the same deliberation-quality score to a $40$-speaker transcript and to a connection-error dump is not measuring deliberation quality. This is a stronger statement than the low inter-run $\kappa$ of Section~\ref{sec:evaluator_reliability}, which shows the rubric is unreliable; the present measurement shows it is also insensitive to the presence or absence of its object, and it is the reason we regard the decision to exclude all seven dimensions from the composite score as necessary rather than conservative.


\textbf{(4) Status of the fourth model.}
\label{sec:fourth_model_status}
ECN-BENCH designates four primary benchmark models: Llama-3.1-8B, Mistral-7B-Instruct-v0.1, Qwen2.5-7B and Qwen2.5-14B. The Qwen2.5-14B campaign's original 24 scorable events were executed as two jobs---11381327 (29 events, 25 July 2026) and 11409272 (event C15 catch-up, 27 July 2026)---with outputs consolidated on the cluster at \texttt{simulation\_logs/\allowbreak qwen2.5-14b-awq/\allowbreak ecnbench\_\allowbreak 20260727\_\allowbreak qwen14b\_\allowbreak complete/}. Earlier drafts stated that its \texttt{event\_results.json} had not been transferred off the cluster and, once retrieved, that $87$ of its $90$ records carried an empty \texttt{ground\_truth} field and were therefore unscorable. Both are now superseded: the field has been backfilled from the authoritative catalogue, and $6$ of the $30$ events (T3, T4, T5, T6, C6, T8) were found, on backfilling, to have been simulated against the wrong topic entirely---a deeper reach of the same stale-catalogue collision documented in Section~\ref{sec:replication_design}, which for the other three campaigns corrupted only a downstream micro-question field rather than the swarm's own deliberation. That left a $24$-event, $72$-unit campaign, compared against the other three models on the $17$ of those events that also passed the contamination check of Appendix~\ref{sec:cutoff}, in the version of this section prior to 2026-08-21.

\textbf{2026-08-21: the six events recovered, Qwen2.5-14B now complete at 30/30.} The six events were re-simulated end-to-end against the correct root catalogue---300 agents, 60 rounds, all three conditions, split across two SLURM jobs (T3/T4/T5 and T6/C6/T8; the cluster's \texttt{medium} QOS caps a user at 2 GPUs, so the two jobs ran sequentially rather than in parallel) after a first re-simulation attempt reproduced the identical stale-catalogue bug (its launch script still resolved \texttt{--events-raw} to the wrong catalogue by a relative path). The 18 new units were rescored under all five evaluator passes used elsewhere in this section (primary \texttt{deepseek/deepseek-r1}, three Evaluator~A repeats, Evaluator~C) and merged into all five result files, verified against the required post-merge invariants (no empty evidence, question or ground truth; Condition~A correctly flagged \texttt{simulation\_executed=False}; no duplicate unit ids). Qwen2.5-14B is now a $30$-event, $90$-unit campaign scored the same way as the other three models, and the contamination-clean comparison of Section~\ref{sec:four_model_extension} and Table~\ref{tab:four_model_clean} now covers the full $20$-event globally-clean set of Appendix~\ref{sec:cutoff} rather than the $17$-event subset that excluded T6, C6 and T8. Every ``three-model'' statement elsewhere in this paper predates the original 2026-08-14 recovery and should still be read as ``three of the four designated models''; only Section~\ref{sec:four_model_extension} and Table~\ref{tab:four_model_clean} include the fourth.

\textbf{What remains a gap.} Most of the Qwen2.5-14B simulation artifacts (persona files, \texttt{.db} traces, logs) are unavailable, so this model remains largely absent from the run-integrity audits that read those artifacts directly. The six events re-simulated 2026-08-21 were transferred off the cluster on 2026-08-26 and are released ($18$ units), together with a three-unit fragment of the original campaign (event~C15); the remaining $21$ events' transcripts had by then been removed by the cluster's scratch-filesystem retention policy and are unrecoverable. Those audits are therefore not re-runnable for this model at campaign scale, whatever the release contains. Earlier revisions gave that range as Appendices~\ref{sec:persona_audit}--\ref{sec:evidence_depth} and thereby excluded the model from the evidence-depth audit as well; that was too wide. Section~\ref{sec:evidence_depth} reads \texttt{evidence\_text} out of the result records rather than off the cluster, so it applies to Qwen2.5-14B unchanged, and Table~\ref{tab:probe_share} now reports it---at a median $90\%$ telemetry-probe share, the deepest of the four campaigns. Its interaction counts in the action-density tables elsewhere in this section (2,407 actions, 180 DBs, 13.38/DB) come from a cluster-side SQL scan taken 2026-07-27, over what was at the time a full 30-event campaign on disk---which includes the six events later found to be simulated against the wrong topic (\S\ref{sec:replication_design}), so those counts describe real social-media actions produced under the \emph{wrong-topic} dossiers, not the correct ones re-simulated 2026-08-21. The re-simulated \texttt{.db} traces for those six events have since been transferred and released, but were never folded into a fresh scan---the 180 DBs the 2026-07-27 scan read no longer exist, so the scan cannot be repeated as taken. Tables~\ref{tab:action_density_30event}--\ref{tab:action_breakdown} therefore carry a topic mismatch on 6 of Qwen2.5-14B's 180 DBs that the other three models' rows do not have, and a reader wishing to correct it must rescan the released 18 units against a fresh baseline rather than reproduce the published counts. We note that earlier drafts of this paper, based on the undercounted \texttt{actions.jsonl} summary, incorrectly characterized Qwen-14B as ``near-zero'' interaction; the authoritative \texttt{.db} counts correct this. All interaction counts in Tables~\ref{tab:action_density_30event}--\ref{tab:action_breakdown} are from direct SQL queries on the authoritative \texttt{trace} tables of the \texttt{.db} files. The \texttt{actions.jsonl} summary files, used in earlier drafts, systematically undercount organic actions due to a cursor-tracking bug in the telemetry-probe function and should not be used for interaction analysis.

\textbf{(5) The outage was repaired, and the first repair attempt was itself defective.} The 24 affected events (S7, the transitional unit, plus the 23 events after it) were re-simulated end-to-end on 2026-08-16 against the original endpoint's cloud equivalent (\texttt{meta-llama/llama-3.1-8b-instruct} via OpenRouter, in place of the air-gapped node's local AWQ-INT4 serving) with real Twitter/Reddit social simulation and Neo4j-backed memory. This is a genuine model-serving substitution, not merely a re-score, and should be read as such when comparing Llama-3.1-8B's post-fix numbers against the other three models, which ran entirely on local HPC serving. The pre-fix data this section and Section~\ref{sec:evidence_depth} analyze is preserved unmodified in the campaign directory's \texttt{archive\_pre\_session\_repairs/} folder; the per-unit figures quoted above (e.g.\ \texttt{C10\_B\_r1}) describe that archived snapshot. A single Qwen2.5-14B unit (\texttt{C9\_C\_r1}) carries the identical corruption signature but was left unrepaired, since \texttt{qwen2.5:14b} is not served by either OpenRouter or Ollama Cloud.

The first attempt to merge that re-simulation into the campaign was defective in three independent ways, all traced to one script that rebuilt the result rows from each unit's \texttt{simulation\_config.json}. (i) It read the keys \texttt{question} and \texttt{evidence\_text}, neither of which that file contains (it stores \texttt{event\_question}, and carries no evidence field at all), so both came back empty and \emph{all 72 units were scored against an empty evidence string}---the evaluator saw only the question, the letter A, B or C, and nothing else, making the three conditions differ by a single character. (ii) It resolved event metadata against a stale catalogue copy in which seven identifiers (C6, T1, T3, T4, T5, T6, T8) denote different events; two of those collisions (T1, T5) share option set and ground truth with the intended events and so pass every automated consistency check while simulating an unrelated question. (iii) The merge filtered on \texttt{event\_id} without filtering on condition, overwriting Condition A---which has no simulation to re-read---and stamping \texttt{simulation\_executed=true} on it, and it fed one source file into both the primary and the replication result files, making two nominally independent evaluator passes identical on all 72 units. The negative post-release lift and the reversed susceptibility that an earlier version of this paper reported as unexplained findings were artifacts of (i) and (iii): restoring Condition A alone returns the post-release lift from $-0.161$ (interval excluding zero) to $-0.028$ (spanning it).

The corrected repair, reported throughout this paper, rebuilds \texttt{evidence\_text} with the pipeline's own \texttt{build\_evidence\_text} from each unit's \texttt{actions.jsonl} and \texttt{simulation.log} (mean $19{,}706$ characters, against $19{,}412$ for the units the outage never touched), re-simulates the seven catalogue-collision events against the correct catalogue and seed dossiers, re-scores all 48 simulated units under three independent evaluators plus three repeat passes, and restores Condition A from the shared No-Sim record rather than re-eliciting it---Condition A is a single record shared byte-for-byte across all three campaigns (Section~\ref{sec:condition_a}), it has no transcript to score, and the first repair overwrote it only as a side effect of the missing condition filter. Independence of the evaluator passes is restored: the primary and replication files now agree on $4$ of $90$ units, against $74$ of $90$ under the defective merge. One deviation is required and is not cosmetic: the reasoning evaluators exhaust the pipeline's default $1500$-token completion budget on real $\approx19$k-character evidence and return a null completion, which killed $15$ of the first $34$ units; the budget was raised to $6{,}000$--$16{,}000$ tokens for this re-scoring. The defective merge never encountered this failure precisely because its empty evidence field kept the evaluator's reasoning short---the first defect concealed the second.
\section{System Architecture and Engineering (MiroFish-Offline)}
\label{sec:architecture}

The ECN-BENCH execution stack is optimized for air-gapped HPC deployment. The simulation pipeline operates in six sequential steps, outlined in Figure~\ref{fig:pipeline}.

\begin{figure*}[!tbp]
\centering
\begin{tabular}{|c|}
\hline
\textbf{Step 1: GraphRAG Build} \\
$\downarrow$ \\
\textbf{Step 2: Persona Generation (4-Axis DNA, Neo4j Graph Memory)} \\
$\downarrow$ \\
\textbf{Step 3: 60-Step OASIS Simulation (MiroFish-Offline Engine)} \\
$\rightarrow$ \textit{Step 30: Perturbation Injection (Relevant vs. Null control)} \\
$\downarrow$ \\
\textbf{Step 4: Output Extraction (Telemetry Logging \& Actions DB)} \\
$\downarrow$ \\
\textbf{Step 5: ReportAgent MCQ Interview (Fixed Local Evaluator)} \\
$\downarrow$ \\
\textbf{Step 6: Scoring \& Statistical Inference (Brier, Esteban-Ray, DeGroot, DQI, DM Test)} \\
\hline
\end{tabular}
\caption{The ECN-BENCH Multi-Agent Social Simulation Pipeline.}
\label{fig:pipeline}
\end{figure*}

\subsection{The Local Serving Stack}
Because university supercomputers block outbound internet access during compute jobs, all dependencies must run locally. The offline stack consists of:
\begin{itemize}
    \item \textbf{Primary Inference Server:} vLLM~\cite{vllm} is utilized to host open-weight frontier models (e.g., Qwen2.5-72B-Instruct, Llama-3-70B-Instruct). vLLM's continuous batching and PagedAttention~\cite{vllm} maximize throughput.
    \item \textbf{Embedding Server:} An Infinity or Ollama background server hosting \texttt{nomic-embed-text} handles embedding queries for vector similarity.
    \item \textbf{Graph Engine:} A self-hosted Neo4j Community Edition instance coupled with the Graphiti context library manages the relational and temporal state of the swarm memory.
\end{itemize}

\subsection{HPC Scalability Migration (Architecture v5.0 - v5.10)}
Scaling multi-agent simulations to thousands of concurrent agents introduced severe system-level bottlenecks. ECN-BENCH implemented a series of engineering upgrades to achieve absolute stability:

\subsubsection{Hybrid LLM Routing (v5.0)}
Ollama-only serving failed under high-concurrency telemetry probes, leading to fatal 3-hour timeouts. Ollama lacks continuous batching. ECN-BENCH decoupled the inference stack: primary chat and generation are routed to a local \textbf{vLLM server} (Port 8000), while embedding pipelines are offloaded to a background \textbf{Ollama instance} (Port 11434). This preserves vLLM KV-cache capacity for agent conversation.

\subsubsection{Multi-GPU Server Isolation (v5.5)}
To support 100-agent simultaneous telemetry probes without stalling the social simulation, a 4-GPU isolation protocol was established:
\begin{itemize}
    \item \textbf{Benchmark Server (GPUs 0,1):} Handles active agent interactions (Post/Reply/Like) on Port 8000.
    \item \textbf{Evaluator/Telemetry Server (GPUs 2,3):} Dedicated to probability probing and final scoring on Port 8001.
\end{itemize}
This ensures that analytical compute loads do not cause global timeouts in the active simulation loop.

\subsubsection{Apptainer Containerization (v5.6 - v5.7)}
Host-level virtual environments on HPC systems frequently trigger driver mismatches (e.g., \texttt{NVIDIA driver too old}). ECN-BENCH migrated to \textbf{Apptainer (Singularity)} containers, building a self-contained environment from \texttt{vllm-openai:v0.6.3.post1}. To support state-of-the-art models (like Gemma 4 MoE architectures) which are unrecognized by the container's built-in \texttt{transformers} library, we implemented a \textbf{Surgical Bind-Mount Override}:
\begin{equation}
    \begin{split}
    \text{Bind Mount Target: } -B \; &\text{\texttt{/scratch/waves/transformers}} \\
    :\,&\text{\texttt{/usr/local/lib/python3.12/}} \\
     &\text{\texttt{dist-packages/transformers}}
    \end{split}
\end{equation}
This replaces the immutable container libraries at runtime with development versions, adjusting for Python 3.12 container environments while the host runs Python 3.11.

\subsubsection{NFS Database Concurrency and Waves NVMe Migration}
SQLite database files hosted on standard HPC Network File Systems (NFS) experienced slow file locking, causing fatal \texttt{database is locked} errors during parallel telemetry probes. Two protocols resolved this:
\begin{enumerate}
    \item \textbf{The Connection Lock Protocol:} All SQLite connectors enforce a strict 60-second queue timeout:
    \begin{verbatim}
    sqlite3.connect(database_path, timeout=60)
    \end{verbatim}
    \item \textbf{Waves NVMe Relocation:} Neo4j graph data and active \texttt{.db} files are symlinked directly onto local high-I/O scratch spaces (\texttt{/scratch/waves/} NVMe drives), completely eliminating network latency lock contention.
\end{enumerate}


\section{Experimental Results: The July 2026 Scaling Campaign}
\label{sec:results}

\begin{figure*}[!tbp]
\centering
\includegraphics[width=\textwidth]{n_sweep_calibration.png}
\caption{Agent population scaling vs.~forecasting calibration for Llama-3.1-8B-AWQ ($R=60$ rounds). Grouped bars show Conditions A (No-Sim), B (With-Sim), and C (Null Control) across agent counts $N \in \{30, 100, 300, 1000, 3000\}$ on Events C1 (US Election, left) and C9 (UK Trade, right). The shaded region marks the empirically identified optimal zone N=100--300.}
\label{fig:n_sweep_calibration}
\end{figure*}

\begin{figure*}[!tbp]
\centering
\includegraphics[width=\textwidth]{r_sweep_calibration.png}
\caption{Deliberation round scaling vs.~calibration for Llama-3.1-8B-AWQ ($N=300$ agents). Bars compare With-Sim (Condition B) and Null Control (Condition C) at $R \in \{30, 60, 90\}$ against the No-Sim baseline (Condition A). The green shading marks the selected $R=60$; the red arrow on Event C9 marks the $R=90$ degradation, which on that event is only $+0.022$ Brier (Appendix~\ref{sec:r_scaling}).}
\label{fig:r_sweep_calibration}
\end{figure*}

\begin{figure*}[!tbp]
\centering
\includegraphics[width=\textwidth]{cross_model_comparison.png}
\caption{Cross-model calibration validation at the derived optimal configuration ($N=300$ agents, $R=60$ rounds). Four model sweeps -- Llama-3.1-8B, Qwen2.5-7B, Qwen2.5-14B (C1 only), and Mistral-7B -- are compared under three conditions on Events C1 and C9. The annotation $\Delta$ denotes the With-Sim lift over the No-Sim baseline; negative lift (red) indicates simulation degradation under the evaluated conditions.}
\label{fig:cross_model_comparison}
\end{figure*}

We present the empirical findings from our multi-GPU scaling campaign executed on a university HPC cluster (A100 GPUs, July 2026). This section and Section~\ref{sec:full_campaign} report the benchmark comparison that motivated this paper and are the context in which the central finding of Section~\ref{sec:degenerate}---loss of sharpness and degenerate consensus under simulation---was discovered; per the caveats established in Sections~\ref{sec:condition_a} and~\ref{sec:serving_note}, the configuration-selection and cross-model conclusions below should be read as engineering and descriptive findings rather than as validated claims about architecture. The experimental design follows a two-stage protocol:

\begin{enumerate}
    \item \textbf{Stage 1 --- Reference Model Scaling Campaign (Llama-3.1-8B-AWQ):} A comprehensive sweep across both the agent population axis ($N \in \{30, 100, 300, 1{,}000, 3{,}000\}$) and the deliberation round axis ($R \in \{30, 60, 90\}$) is conducted with Llama-3.1-8B-AWQ as the designated reference model. This stage characterizes the scaling laws of social simulation calibration and identifies the Pareto-optimal configuration on the cost-performance frontier.
    \item \textbf{Stage 2 --- Cross-Model Validation (Qwen2.5-7B, Qwen2.5-14B, Mistral-7B):} At the derived optimal configuration from Stage 1, three additional model families are evaluated on Events C1 and C9 to assess generalizability of findings across architectures. This validates whether the identified optimal configuration represents a model-agnostic sweet spot or is architecture-dependent.
\end{enumerate}

All calibration metrics use the Multi-Category Brier Score ($BS$); lower values indicate better probabilistic calibration.

\textbf{A statistical caveat that governs this entire section.} Every cell of the N-sweep and R-sweep matrices below is a \emph{single run of a single event}: one configuration, one condition, one repeat ($\texttt{repeats}=1$ in the run manifest). There are no replicates, so no cell carries a variance estimate and no difference between two cells can be tested. The two anchor events C1 and C9 were chosen for compute reasons, not sampled. Consequently the configuration derived here---$N=300$, $R=60$---should be read as a \emph{defensible engineering choice supported by four independent criteria} (diversity, statistical validity, compute budget, echo-chamber avoidance), not as an empirically established optimum. Where earlier drafts described sweep-to-sweep changes as ``dramatic,'' ``catastrophic'' or ``skyrocketing,'' we have replaced that language with the measured deltas. Establishing a genuine scaling law would require replicated runs across a sampled event set, which we identify as future work.

\subsection{Stage 1: Agent Population Scaling (Llama-3.1-8B Reference Model)}
\label{sec:n_scaling}

Figure~\ref{fig:n_sweep_calibration} plots the full N-sweep calibration profile of Llama-3.1-8B-AWQ at $R=60$ rounds on both Events C1 and C9, and Table~\ref{tab:n_sweep_results} gives the underlying Brier scores. The results reveal a characteristic non-monotonic pattern:

\begin{itemize}
    \item \textbf{Cognitive Under-representation ($N = 30$):} Small agent counts fail to adequately sample the high-dimensional opinion space $\Omega$. The With-Sim condition achieves Brier scores of $0.530$ (C1) and $0.595$ (C9)---worse than the $N=100$ configuration on both events, consistent with stochastic sampling bias at low population densities. Mathematically, when $N$ is small, individual agent biases dominate the aggregated forecast rather than statistical cancellation.
    \item \textbf{Optimal Zone ($N = 100$--$300$):} Performance peaks in this range. On Event C1, the With-Sim condition achieves its best Brier score at $N=100$ ($BS_B = 0.396$), while on Event C9 the gradient is shallower and performance stabilizes across $N \in [100, 300]$ ($BS_B \approx 0.697$--$0.763$). We therefore identify \textbf{$N=100$--$300$} as the empirical optimal zone: it is the smallest range that simultaneously satisfies (i) sufficient cognitive diversity, (ii) statistical validity ($\epsilon \leq 0.10$ for $N \geq 100$), and (iii) avoidance of the conversational noise saturation described below. We standardize on $N=300$ as the production configuration because it provides the stronger margin of error ($\epsilon \approx 0.056$) required for multi-event significance testing.
    \item \textbf{Degradation at large $N$ ($N \geq 1{,}000$), non-monotonic:} Both large-$N$ configurations score worse than the optimal zone, but the pattern is not monotone in $N$ and we caution against reading it as one. On Event C1, $BS_B$ rises to $0.953$ at $N=1{,}000$ and then \emph{falls back} to $0.613$ at $N=3{,}000$; on C9 it moves $0.726 \to 0.789$. With one run per cell, the $N=1{,}000$ C1 value could equally be an outlier as a point on a curve. What the data supports is the weaker statement that neither large-$N$ configuration outperforms $N=100$--$300$ on either event. Two mechanisms are plausible candidates for such an effect, and we state them as hypotheses rather than as formalisations of a measured curve:
        \begin{enumerate}
            \item \textbf{Quadratic Edge Expansion and SNR Collapse:} In a multi-agent social network of size $N$, let $G = (V,E)$ be the directed interaction graph. The number of potential communicative links grows quadratically:
            \begin{equation}
                |E| = N(N - 1) \approx \mathcal{O}(N^2)
            \end{equation}
            At $N=300$, $|E| \approx 8.97 \times 10^4$; at $N=3{,}000$, $|E| \approx 8.99 \times 10^6$. This quadratic explosion floods the Graphiti temporal context engine with semantic redundancies, saturating individual agent memory with circular arguments. The Signal-to-Noise Ratio (SNR) of the forecasting signal collapses.
            \item \textbf{DeGroot Belief Convergence Decay:} Under the DeGroot model $x_i(t+1) = \sum_j W_{ij} x_j(t)$, at extreme $N$ the social weights $W_{ij}$ converge to a low-entropy uniform distribution, causing individual private analytical signals to evaporate into uncritical social mimicry.
        \end{enumerate}
\end{itemize}

\begin{table}[!htbp]
\centering
\caption{Llama-3.1-8B-AWQ: N-Scaling Calibration Matrix at $R=60$ (Brier Scores, lower is better)}
\label{tab:n_sweep_results}
\begin{tabular}{lcccccc}
\toprule
 & \multicolumn{3}{c}{\textbf{Event C1: US Election}} & \multicolumn{3}{c}{\textbf{Event C9: UK Trade}} \\
\cmidrule(r){2-4} \cmidrule(l){5-7}
\textbf{$N$ (agents)} & \textbf{Cond A} & \textbf{Cond B} & \textbf{Cond C} & \textbf{Cond A} & \textbf{Cond B} & \textbf{Cond C} \\
\midrule
30      & 0.3368 & 0.5300 & 0.4200 & 0.2150 & 0.5952 & 0.6069 \\
\textbf{100}     & 0.3368 & \textbf{0.3958} & 0.5958 & 0.2150 & 0.6971 & 0.6329 \\
\textbf{300}     & 0.3368 & 0.5350 & 0.6800 & 0.2150 & 0.7630 & 0.8197 \\
1000    & 0.3368 & 0.9533 & 0.9000 & 0.2150 & 0.7264 & 0.6417 \\
3000    & 0.3368 & 0.6133 & 0.7911 & 0.2150 & 0.7885 & 0.6095 \\
\bottomrule
\end{tabular}
\setlength{\tabcolsep}{6pt}
\end{table}

\subsection{Stage 1: Deliberation Round Scaling (Llama-3.1-8B Reference Model)}
\label{sec:r_scaling}

Figure~\ref{fig:r_sweep_calibration} presents the R-sweep calibration profile at $N=300$, with the underlying scores in Table~\ref{tab:r_sweep_results}. The results identify a clear deliberation optimum and a pathological collapse boundary:

\begin{itemize}
    \item \textbf{Under-deliberation ($R=30$):} At 30 rounds, the swarm has not yet reached belief convergence. The With-Sim condition achieves Brier scores of $0.720$ (C1) and $0.772$ (C9)---substantially worse than the direct No-Sim baseline ($0.337$ on C1). Empirical Jensen-Shannon Divergence (JSD) telemetry confirms that belief distributions at $R=30$ still exhibit high inter-agent variance, indicating incomplete consensus formation.
    \item \textbf{Optimal Deliberation ($R=60$):} At 60 rounds, calibration improves substantially vs. $R=30$ on both events. On C1, $BS_B$ drops from $0.720$ to $0.535$ (26\% improvement); the With-Sim condition achieves lift over the No-Sim baseline ($\Delta = -0.198$). This confirms that $R=60$ is sufficient for frontier-class models ($\geq$7B parameters) to refine collective forecasts through directed social deliberation.
    \item \textbf{Degradation at $R=90$:} Both events score worse at $R=90$ than at $R=60$, but by very different margins. On Event C1 the change is large ($BS_B$: $0.535 \to 0.998$); on Event C9 it is $0.763 \to 0.785$, a shift of $0.022$ that is well within the range we would expect from run-to-run noise and which earlier drafts described as ``skyrocketing.'' The Null Control on C9 moves $0.820 \to 0.840$. A plausible mechanism is that at $R=90$ agents take 50\% more conversational turns without new external information, so the social feedback loop amplifies existing bias into overconfident consensus. We stress that a single large change on one event does not locate a boundary: the ``echo-chamber boundary at $R \approx 70$--$80$'' proposed in earlier drafts was interpolated between two measured points on two events with no replicates, and we withdraw it as a quantitative claim. What the sweep supports is that $R=90$ was not better than $R=60$ on either event tested.
\end{itemize}

\begin{table}[!htbp]
\centering
\caption{Llama-3.1-8B-AWQ: R-Scaling Calibration Matrix at $N=300$ (Brier Scores, lower is better)}
\label{tab:r_sweep_results}
\begin{tabular}{lcccccc}
\toprule
 & \multicolumn{3}{c}{\textbf{Event C1: US Election}} & \multicolumn{3}{c}{\textbf{Event C9: UK Trade}} \\
\cmidrule(r){2-4} \cmidrule(l){5-7}
\textbf{$R$ (rounds)} & \textbf{Cond A} & \textbf{Cond B} & \textbf{Cond C} & \textbf{Cond A} & \textbf{Cond B} & \textbf{Cond C} \\
\midrule
30  & 0.3368 & 0.7200 & 0.6283 & 0.2150 & 0.7721 & 0.7386 \\
\textbf{60}  & 0.3368 & \textbf{0.5350} & 0.6800 & 0.2150 & \textbf{0.7630} & 0.8197 \\
90  & 0.3368 & 0.9983 & 0.6958 & 0.2150 & 0.7850 & 0.8401 \\
\bottomrule
\end{tabular}
\setlength{\tabcolsep}{6pt}
\end{table}

\subsection{Optimal Configuration: $N=300$, $R=60$}
\label{sec:optimal_config}

Combining the N-sweep and R-sweep findings, we select $\mathbf{N=300}$, $\mathbf{R=60}$ as the standardized production configuration for ECN-BENCH. While N=100 yields slightly better raw Brier performance on C1 ($0.396$ vs.~$0.535$), the $N=300$ configuration is selected for three reasons:
(1) it provides the statistical validity threshold ($\epsilon = 0.056 < 0.10$) required for multi-event Diebold-Mariano significance testing;
(2) it was expected to achieve superior cognitive diversity coverage ($\geq 75$ agents per persona axis dimension)---a design intent that Appendix~\ref{sec:persona_audit} shows the pipeline realised in only one of the three campaigns, and which should not be read as a measured property of these runs; and
(3) it delivers more stable performance across heterogeneous events (C1 and C9 show less divergent behavior at $N=300$ than at $N=100$).

The $R=60$ setting is supported more weakly: it improves on $R=30$ by $26\%$ on C1 and is not worse than $R=90$ on either anchor event. With one run per cell we do not claim it is an optimum, nor that a collapse boundary has been located.

\subsection{Stage 2: Cross-Model Validation at $N=300$, $R=60$}
\label{sec:cross_model}

\confoundnote{The four models compared here differ simultaneously in release
generation, numerical precision and inference server
(Section~\ref{sec:serving_note}, Table~\ref{tab:provenance}), and this sweep
predates the evaluator audit of Part~I, so its scores come from the single
Evaluator~J pass. It is reported as an engineering record of the configuration
sweep, not as a model ranking; the architectural interpretations offered in
earlier drafts of this appendix are withdrawn in
Appendix~\ref{sec:errata_appendix}.}

Figure~\ref{fig:cross_model_comparison} presents the cross-model calibration comparison at the derived optimal configuration, with the per-condition Brier scores tabulated in Table~\ref{tab:cross_model_results}. Three additional model families are evaluated alongside the reference model:

\begin{itemize}
    \item \textbf{Qwen2.5-7B (Positive Lift on C1):} At $N=300$, $R=60$, the Qwen2.5-7B swarm achieves $BS_B = 0.445$ on Event C1, below both the No-Sim baseline ($BS_A = 0.560$, $\Delta = +0.115$) and the Null Control ($BS_C = 0.622$). This demonstrates meaningful with-sim lift. On Event C9, $BS_B = 0.608$ vs.~$BS_A = 0.215$, yielding a negative lift ($\Delta = -0.393$), consistent with the event's high information complexity.
    \item \textbf{Qwen2.5-14B (Marginal Lift on C1):} On Event C1, $BS_B = 0.598$ vs.~$BS_A = 0.560$ ($\Delta = -0.038$), a marginal negative lift. The larger parameter count does not translate to improved social calibration at $R=60$, consistent with the hypothesis that instruction-tuned models larger than 7B may suffer from stronger conformity priors that resist group deliberation. A complete 30-event campaign for Qwen2.5-14B exists on the HPC cluster (Jobs 11381327 and 11409272, July 25--27, 2026). A direct .db scan (180 databases) confirms 2,407 social actions (13.38/DB: 1,215 posts, 585 comments, 461 likes, 11 follows)---moderate, analysis-focused interaction, higher than Qwen-7B (723) but well below Llama-8B (27,087) and Mistral-7B (9,364). Its full-campaign Brier statistics are pending local download and are excluded from the aggregate analysis reported here.
    \item \textbf{Mistral-7B-Instruct-v0.1 (Negative Lift at $R=60$; architecture-dependent optimum):} On Event C1, $BS_B = 0.660$ vs.~$BS_A = 0.560$ ($\Delta = -0.100$)---the With-Sim condition degrades calibration relative to direct elicitation. On Event C9, the degradation is more severe ($BS_B = 0.829$, $\Delta = -0.614$). Earlier drafts read this as an \textbf{architecture-dependent optimum}---the claim that Mistral-7B saturates its deliberation budget before $R=60$ and therefore requires different production settings. We withdraw that reading. As Section~\ref{sec:serving_note} establishes, this model is \texttt{Mistral-7B-Instruct-v0.1}, released roughly ten months before the other two and served at full precision rather than AWQ INT4; generation, precision and serving stack are all confounded with architecture. The measured fact is that this model configuration performs worst under the production settings. Whether that is an architectural property, a consequence of model age, or an artifact of the precision and serving differences cannot be determined from these data. A dedicated R-sweep on matched-generation, matched-precision models would be required, and we no longer recommend different production settings for Mistral on this evidence.
\end{itemize}

\begin{table*}[!htbp]
\centering
\caption{Cross-Model Calibration Validation at Optimal Configuration ($N=300$, $R=60$, Brier Scores)}
\label{tab:cross_model_results}
\begin{tabular}{llcccccc}
\toprule
 & & \multicolumn{3}{c}{\textbf{Event C1: US Election}} & \multicolumn{3}{c}{\textbf{Event C9: UK Trade}} \\
\cmidrule(r){3-5} \cmidrule(l){6-8}
\textbf{Model} & \textbf{Role} & \textbf{Cond A} & \textbf{Cond B} & \textbf{Cond C} & \textbf{Cond A} & \textbf{Cond B} & \textbf{Cond C} \\
\midrule
Llama-3.1-8B-AWQ  & Reference  & 0.3368 & 0.5350 & 0.6800 & 0.2150 & 0.7630 & 0.8197 \\
Qwen2.5-7B    & Validation & 0.5600 & \textbf{0.4450} & 0.6217 & 0.2150 & 0.6080 & 0.7209 \\
Qwen2.5-14B   & Validation & 0.5600 & 0.5978 & 0.5958 & ---   & ---   & ---   \\
Mistral-7B-Instruct-v0.1    & Validation & 0.5600 & 0.6600 & 0.6333 & 0.2150 & 0.8286 & 1.1037 \\
\bottomrule
\end{tabular}
\setlength{\tabcolsep}{6pt}
\end{table*}

\subsection{Optimal Settings Rationale and Cognitive Boundaries}
\label{sec:optimal_bounds}

We now characterize the precise cognitive and computational boundaries established through the Stage~1 reference campaign, with cross-model validation from Stage~2.

Table~\ref{tab:deliberation_susceptibility_n300} presents the full 3-condition Brier matrix at the standardized configuration ($N=300$, $R=60$) for all evaluated models and events.

\begin{table*}[!htbp]
\centering
\caption{Empirical Deliberation Susceptibility Matrix at Production Configuration ($N=300$, $R=60$). Brier scores; conditions are A (No-Sim), B (With-Sim) and C (Null injection).}
\label{tab:deliberation_susceptibility_n300}
\footnotesize
\setlength{\tabcolsep}{4pt}
\begin{tabular}{llcccccc}
\toprule
 & & \multicolumn{3}{c}{\textbf{Event C1: US Election}} & \multicolumn{3}{c}{\textbf{Event C9: UK Trade}} \\
\cmidrule(r){3-5} \cmidrule(l){6-8}
\textbf{Model} & \textbf{Stage} & \textbf{A} & \textbf{B} & \textbf{C} & \textbf{A} & \textbf{B} & \textbf{C} \\
\midrule
\textbf{Llama-3.1-8B-AWQ}  & Reference  & 0.3368 & 0.5350 & 0.6800 & 0.2150 & 0.7630 & 0.8197 \\
\midrule
\textbf{Qwen2.5-7B}    & Validation & 0.5600 & \textbf{0.4450} & 0.6217 & 0.2150 & 0.6080 & 0.7209 \\
\textbf{Qwen2.5-14B}   & Validation & 0.5600 & 0.5978 & 0.5958 & --- & --- & --- \\
\textbf{Mistral-7B-Instruct-v0.1}    & Validation & 0.5600 & 0.6600$^*$ & 0.6333 & 0.2150 & 0.8286$^*$ & 1.1037 \\
\bottomrule
\end{tabular}
\small $^*$Mistral-7B exhibits negative With-Sim lift at $R=60$ (simulation degrades vs.~No-Sim baseline), indicating architecture-dependent saturation at $R < 60$.
\end{table*}

\begin{figure*}[!tbp]
\centering
\includegraphics[width=\textwidth]{compute_cost_analysis.png}
\caption{Computational complexity scaling vs.~agent population (Llama-3.1-8B-AWQ, $R=60$, Event C1). Stacked bars show Phase A (persona generation + evaluation, purple) and simulation phases B+C (blue) in minutes; the red dashed line shows total wall-clock time. The LLM persona generation component is approximately $0.8$\,min of the $\sim$3--4\,min Phase A. The shaded zone marks the optimal $N=100$--$300$ range, where calibration performance peaks at manageable GPU cost. Beyond $N=1000$, compute scales super-linearly while calibration degrades.}
\label{fig:compute_cost_analysis}
\end{figure*}

\subsubsection{Swarm Population Capacity Bound ($N = 100$--$300$ Optimal Zone)}
The choice of $N=300$ as the production standard (with $N=100$ as the computationally lighter alternative) is established through a multi-criteria decision framework balancing four competing objectives: (1) cognitive diversity coverage, (2) statistical validity, (3) computational feasibility, and (4) echo-chamber avoidance.

\textbf{Criterion 1: Cognitive Diversity Coverage.}
Small populations ($N < 100$) suffer from \textbf{cognitive under-representation}; they fail to adequately sample the high-dimensional opinion space $\Omega$ (e.g., Llama-3.1-8B: $BS_B = 0.530$ at $N=30$ vs.~$0.396$ at $N=100$ on Event C1). At $N=300$ the swarm has sufficient \emph{slots} to sample the full spectrum of persona viewpoints encoded in the 4-axis DNA space (Worldview, Motivation, Style, Biases) at $\geq 75$ distinct agents per dimension. Whether those slots are filled with distinct personas is a separate, empirical question, since Stage~1 persona generation is itself an LLM call: Appendix~\ref{sec:persona_audit} measures the realised pools and finds a median of $298$ distinct personas per $300$ agents for Mistral-7B but only $40$ for Qwen2.5-7B and $11$ for Llama-3.1-8B on 23 of its 30 events. The population argument of this subsubsection therefore bounds diversity from above and does not establish it.

\textbf{Criterion 2: Statistical Power and Validity.}
For the Diebold-Mariano test (Section~\ref{sec:statistics}) to achieve statistical significance at $\alpha=0.05$ with power $1-\beta=0.80$, we require a minimum sample size. Let $d_e = L_A(e) - L_B(e)$ denote the loss differential for event $e$. The test statistic is $DM = \bar{d} / \sqrt{\hat{\sigma}_d^2 / N}$, where $N$ is the number of events. For a two-sided test with $\alpha=0.05$ and $\beta=0.20$, the critical value is $z_{1-\alpha/2} + z_{1-\beta} = 1.96 + 0.84 = 2.80$. To detect a medium effect size (Cohen's $d = 0.5$), we require $N \geq (2.80 / 0.5)^2 = 31.36$ events. However, within each event, the swarm comprises $N_{agents}$ agents whose beliefs are aggregated. For the aggregated swarm forecast to achieve a margin of error $\epsilon = 0.05$ at 95\% confidence, we require $N_{agents} \geq (1.96 / 0.05)^2 \times p(1-p)$, where $p$ is the true probability. For $p=0.5$ (maximum variance), this yields $N_{agents} \geq 384.16$. Thus, $N=300$ agents provides a margin of error $\epsilon \approx 0.056$, while $N=100$ provides $\epsilon \approx 0.098$---both are acceptable for benchmarking, but $N=300$ is preferred for multi-event campaigns requiring significance testing.

\textbf{Criterion 3: Computational Feasibility.}
From a systems engineering perspective, $N=300$ represents the practical ceiling for high-throughput, multi-model evaluation on standard university-grade HPC clusters. As documented in Figure~\ref{fig:compute_cost_analysis} and detailed in Table~\ref{tab:runtime_breakdown}, a single condition-run at $N=300$ completes in approximately 75--115 minutes wall-clock (where persona LLM generation alone takes $\sim$0.8 min, the full Phase A---including persona evaluation/scoring---takes $\sim$2.7--3.8 min, swarm simulation B+C takes $\sim$60--70 min, e.g., 35.1+31.4 min for Llama-3.1-8B on Event C1, and evaluator scoring takes $\sim$1.2 min). A full 3-condition event evaluates in approximately 1.5--2 hours, and a 10-event campaign under 20 hours.

\textbf{Criterion 4: Echo-Chamber Avoidance.}
Conversely, neither large-population configuration ($N \ge 1{,}000$) improved calibration, while costing $10$--$100\times$ the compute. A candidate mechanism is that the number of potential conversational channels grows quadratically ($\mathcal{O}(N^2)$), flooding the Graphiti temporal context engine and individual agent timelines with redundant, homogeneous interactions, degrading the signal-to-noise ratio of the forecasting signal. We did not instrument message redundancy or agent-level signal retention, so this mechanism is proposed rather than demonstrated; the measured fact is the absence of benefit, not its cause.

\textbf{Decision:} $N=300$ satisfies all four criteria simultaneously, providing the optimal trade-off between statistical validity ($\epsilon \approx 0.056$), headroom for cognitive diversity ($\geq 75$ agent slots per persona dimension; realised diversity is measured in Appendix~\ref{sec:persona_audit} and falls well short of it for two campaigns), computational feasibility ($<$20 hours for 10-event campaign), and echo-chamber avoidance (SNR remains above critical threshold). $N=100$ is recommended for rapid prototyping or limited-compute contexts where statistical marginal error $\epsilon < 0.10$ is acceptable.

\subsubsection{Deliberation Round Horizons ($R = 60$)}
The choice of $R=60$ rounds is determined by three empirical constraints derived from the Llama-3.1-8B reference campaign: (1) convergence to stable beliefs, (2) avoidance of over-deliberation, and (3) computational cost-effectiveness.

\textbf{Constraint 1: Convergence to Stable Beliefs.}
The swarm must execute sufficient rounds for individual agent beliefs to converge toward a stable collective forecast. Empirical telemetry data (Jensen-Shannon Divergence, JSD) shows that belief distributions stabilize after approximately 40--50 rounds for frontier models ($\geq$7B parameters). At $R=30$, the swarm has not yet reached convergence: Llama-3.1-8B achieves Brier score $0.720$ on Event C1 at $R=30$, but improves to $0.535$ at $R=60$---a 26\% reduction in calibration error. This demonstrates that 30 rounds is insufficient for high-capacity models to fully process the injected information and refine their forecasts through social deliberation.

\textbf{Constraint 2: Echo-Chamber Boundary.}
A plausible concern is that beyond some round threshold, closed-loop social interaction without external grounding drives belief structures into overconfident, miscalibrated consensus. Our data are consistent with degradation somewhere above $R=60$ but do not locate a threshold: at $R=90$, Llama-3.1-8B's $BS_B$ rises from $0.535$ to $0.998$ on Event C1 but only from $0.763$ to $0.785$ on Event C9, with the Null Control moving $0.820 \to 0.840$. Two events, one run each, cannot distinguish a threshold from noise. We adopt $R=60$ as the conservative choice below the region where degradation was observed.

\textbf{Constraint 3: Computational Cost-Effectiveness.}
The simulation subprocess time scales linearly with $R$ (Appendix~\ref{sec:runtime_breakdown}). Thus:
\begin{itemize}
    \item $R=30$: fast, but insufficient convergence for reflective models (Llama-3.1-8B)
    \item $R=60$: \textbf{optimal balance}---convergence achieved, safely below echo-chamber boundary
    \item $R=90$: no better than $R=60$ on either anchor event; 48\% slower
\end{itemize}

\textbf{Architecture-Dependent Considerations.}
The empirical cross-model validation reveals important model-specific deliberation dynamics:
\begin{itemize}
    \item \textbf{Reflective frontier models (Llama-3.1-8B):} Require $R=60$ to fully converge. The marginal calibration gain from $R=30 \to R=60$ is substantial: $BS_B$ improves from $0.720$ to $0.535$ on C1 (26\% reduction).
    \item \textbf{Standard instruction models (Qwen2.5-7B):} Show positive lift over No-Sim baseline at $R=60$ ($BS_B = 0.445 < BS_A = 0.560$ on C1). Standardizing on $R=60$ is appropriate for cross-model comparability.
    \item \textbf{Mistral-7B-Instruct-v0.1:} shows negative With-Sim lift at $R=60$ ($BS_B = 0.660 > BS_A = 0.560$ on C1). We report this without an architectural interpretation, for the provenance reasons given in Section~\ref{sec:serving_note}.
\end{itemize}

\textbf{Decision:} $R=60$ rounds is standardized as the ECN-BENCH production configuration: it satisfies convergence requirements for frontier models (40--50 rounds required), operates safely below the echo-chamber boundary (70--80 rounds), and is computationally cost-effective. Results for Mistral-7B-Instruct-v0.1 are reported with the caveat that its generation, precision and serving stack all differ from the other two models (Section~\ref{sec:serving_note}), so no architecture-specific conclusion should be drawn from them.

\subsubsection{Susceptibility Margins (Condition B vs. Condition C)}
To confirm that the swarm's updates represent rational Bayesian deliberation rather than simple stimulus-reactivity, we isolate the susceptibility delta ($d = BS_C - BS_B$). A positive delta indicates the swarm remains resistant to irrelevant null control perturbations while updating productively on relevant signal. From Table~\ref{tab:deliberation_susceptibility_n300}, at $N=300$, $R=60$: Llama-3.1-8B shows positive susceptibility on C1 ($d = 0.680 - 0.535 = +0.145$) and C9 ($d = 0.820 - 0.763 = +0.057$). Qwen2.5-7B shows strong positive susceptibility on C1 ($d = 0.622 - 0.445 = +0.177$) and C9 ($d = 0.721 - 0.608 = +0.113$). Mistral-7B-Instruct-v0.1 shows a negative susceptibility margin on C1 at $R=60$ ($d = 0.633 - 0.660 = -0.027$). Note that on the full 30-event campaign this same model records the \emph{largest} and only statistically significant positive susceptibility margin ($+0.115$, Table~\ref{tab:cross_model_stats}), which illustrates how little weight single-event sweep cells can bear.

\subsubsection{Empirical Computational Budget}
\label{sec:runtime_breakdown}

Each ECN-BENCH unit decomposes into four sequential phases, instrumented via \texttt{traces/execution.jsonl} wall-clock timestamps.

\begin{table}[!htbp]
\centering
\caption{Measured Per-Event Wall-Clock Budget (single-event, Event C1) at $N=300$, $R=60$ on A100 GPUs (from Slurm sacct)}
\label{tab:runtime_breakdown}
\footnotesize
\setlength{\tabcolsep}{4pt}
\begin{tabular}{lcccc}
\toprule
\textbf{Model} & \textbf{Phase A (Persona Gen+Eval)} & \textbf{Sim B (With-Sim)} & \textbf{Sim C (Null)} & \textbf{Wall-Clock Total} \\
\midrule
Llama-3.1-8B-AWQ  & 3.8 min (229s)   & 35.1 min (2104s) & 31.4 min (1883s) & 77.6 min (4655s) \\
Qwen2.5-7B    & 2.7 min (161s)   & 47.9 min (2874s) & 18.4 min (1103s) & 77.1 min (4624s) \\
Qwen2.5-14B   & 2.7 min (160s)   & 69.5 min (4168s) & 31.5 min (1890s) & 114.0 min (6842s) \\
Mistral-7B-Instruct-v0.1    & 2.7 min (160s)   & 33.7 min (2024s) & 30.1 min (1806s) & 75.4 min (4526s) \\
\bottomrule
\end{tabular}
\smallskip \\
\footnotesize Each row reflects a single-event (C1) three-condition run from the July 2026 v523 optimization sweep at $N=300$, $R=60$ on \texttt{A100}. ``Wall-Clock Total'' is the Slurm job wall time and includes setup/teardown overhead, so phase times need not sum exactly to it. LLM persona generation alone is approximately $49$\,s ($0.8$\,min); the remainder of Phase A ($\sim$110--180\,s) is persona evaluation and scoring.
\end{table}

\textbf{Round Budget Optimization.} For Llama-3.1-8B on Event C1: scaling from $R=30$ ($BS_B=0.720$) to $R=60$ ($BS_B=0.535$) yields a substantial positive gain (+0.185 Brier improvement). Scaling further to $R=90$ yields $BS_B = 0.998$, the worst value in the sweep. On this single event $R=60$ is Pareto-dominant over both alternatives; whether that generalises is untested.

\subsubsection{Configuration Decision Matrix}
Table~\ref{tab:decision_matrix} synthesizes the multi-criteria evaluation that led to the selection of $N=300$ agents and $R=60$ rounds.

\begin{table*}[!htbp]
\centering
\caption{Configuration Decision Matrix: Multi-Criteria Evaluation of Agent Count ($N$) and Round Count ($R$). The Brier column reports the single measured $BS_B$ on Event C1; each is one run with no replicate, so the column ranks configurations but does not establish differences. $N=300$ is selected on the four criteria jointly, not on Brier score, on which $N=100$ is better.}
\label{tab:decision_matrix}
\small
\begin{tabular}{lcccc>{\raggedright\arraybackslash}p{6.5cm}}
\toprule
\textbf{Config} & \textbf{Brier} & \textbf{Stat.} & \textbf{Compute} & \textbf{Echo} & \textbf{Rationale} \\
 & \textbf{Score} & \textbf{Validity} & \textbf{Feasibility} & \textbf{Chamber} & \\
\midrule
$N=30, R=60$  & 0.530 & $\epsilon=0.18$ & $\checkmark\checkmark$ & Safe & Cognitive under-representation \\
$N=100, R=60$ & \textbf{0.396} & $\epsilon=0.098$ & $\checkmark\checkmark$ & Safe & \textbf{Best raw Brier on C1}; marginal stat.\ validity \\
\textbf{$N=300, R=60$} & 0.535 & $\epsilon=0.056$ & $\checkmark\checkmark$ & \textbf{Safe} & \textbf{Selected}: best joint satisfaction of all four criteria, not best Brier \\
$N=300, R=30$ & 0.720 & $\epsilon=0.056$ & $\checkmark\checkmark$ & Safe & Under-converged on C1 \\
$N=300, R=90$ & 0.998 & $\epsilon=0.056$ & $\checkmark$ & Risk & Worst C1 score in the sweep; 48\% slower \\
$N=1000, R=60$ & 0.953 & $\epsilon=0.031$ & $\times$ & Risk & 10$\times$ compute for worse calibration \\
$N=3000, R=60$ & 0.613 & $\epsilon=0.018$ & $\times$ & Risk & 100$\times$ compute for no calibration gain \\
\bottomrule
\end{tabular}
\setlength{\tabcolsep}{6pt}
\end{table*}

\subsection{Hardware Compatibility and bfloat16 Routing}
During the campaign, we identified a critical hardware constraint. Attempting to host reasoning-distilled weights stored in native \texttt{bfloat16} on older \textbf{NVIDIA V100} nodes (compute capability 7.0) triggered immediate vLLM initialization failures. We successfully resolved this by enforcing explicit Slurm routing policies:
\begin{verbatim}
sbatch --constraint=gpu_a100 run_agent_scaling.slurm
\end{verbatim}
This guaranteed execution strictly on Ampere architectures.

\subsection{Small Model Limitations: The Sub-15B Parameter Boundary}
\label{sec:small_model_limits}

In addition to the main benchmark sweeps, we attempted to deploy two sub-15B models on the HPC cluster: \textbf{Gemma 1.1 2B-IT} and \textbf{CodeGemma 7B-IT}. These experiments surfaced three distinct, compounding failure modes that collectively define a practical sub-15B boundary for the unmodified ECN-BENCH pipeline under local vLLM execution.

\subsubsection{Failure Mode I: Context Window Overflow (\texttt{HTTP 400})}

Both Gemma 1.1 2B and CodeGemma 7B natively support a maximum context length of 8,192 tokens. However, the batch Slurm scripts were initially configured with \texttt{--max-model-len 4096}. When the pipeline processes events with dense prompt chains (e.g., events T1 and C1), the concatenated prompt length exceeds this limit, causing the vLLM API server to reject requests with a fatal \texttt{HTTP 400 Bad Request} error. This failure was resolved by patching all Slurm templates to enforce \texttt{--max-model-len 8192}.

\subsubsection{Failure Mode II: KV-Cache VRAM Saturation (OOM)}

Extending the context window to 8,192 tokens dramatically increases the Key-Value (KV) cache size. With initial \texttt{--gpu-memory-utilization 0.20}, the vLLM server crashed with OOM errors during KV cache initialization. Minimum stable utilization thresholds are:
\begin{itemize}
    \item \textbf{Gemma 1.1 2B-IT:} \texttt{--gpu-memory-utilization 0.40}
    \item \textbf{CodeGemma 7B-IT:} \texttt{--gpu-memory-utilization 0.45}
\end{itemize}
These values provide sufficient VRAM headroom for the KV cache while avoiding the full-utilization allocation that can itself trigger allocation failures in the Apptainer container environment.

\subsubsection{Failure Mode III: Guided-Decoding JSON Schema Incompatibility}
\label{sec:schema_bypass}

The most fundamental and irrecoverable failure mode is a structural incompatibility between local vLLM's guided-decoding backend and sub-15B AWQ-quantized models. The ECN-BENCH pipeline is architecturally dependent on \textbf{Stage 1 persona generation}, in which the benchmark calls the model with a strict \texttt{json\_schema} response format to extract a JSON list of 30 diverse stakeholders from the event's seed context.

In the standard configuration, this guided-decoding schema is enforced via the \textbf{outlines}~\cite{vllm_guided} backend within vLLM. However, on the HPC cluster, loading and compiling the outlines grammar for AWQ-quantized models (which involve lower-precision weight layouts) causes one of two failure outcomes:
\begin{enumerate}
    \item \textbf{Compilation hang:} The outlines grammar compilation allocates excessive CPU and VRAM resources, causing the Slurm job to stall indefinitely during server initialization.
    \item \textbf{Silent degrade:} Even when the server initializes, small-parameter models lack sufficient instruction-following capacity to reliably produce schema-valid JSON arrays when guided decoding is effectively disabled or malformed. The model returns raw text, partial arrays, or structurally invalid objects.
\end{enumerate}

In either case, the Stage 1 extraction fails to return a parseable list of stakeholder objects. The ECN-BENCH protocol runner raises a terminal \texttt{ValueError: Intelligent persona generation failed}, immediately aborting the job before any simulation begins. Both Gemma 1.1 2B-IT and CodeGemma 7B-IT exhibited this failure mode consistently across all 30 events under the original unmodified system code.

\subsubsection{The Sub-15B Exclusion Criterion}

The three failure modes above interact multiplicatively: resolving the context-window overflow does not prevent the VRAM saturation, and resolving both does not prevent the guided-decoding failure. The root cause is architectural: sub-15B models in the AWQ quantization format, when served locally via vLLM on the HPC cluster, cannot satisfy the structured output requirements that ECN-BENCH's two-stage persona extraction pipeline demands.

We therefore formalize the following \textbf{Sub-15B Exclusion Criterion} for the ECN-BENCH HPC benchmark:

\begin{quote}
\textit{Models with fewer than 15B active parameters are excluded from local vLLM HPC execution unless one of the following holds: (a) the system code is patched to support non-strict JSON fallback parsing at Stage 1, or (b) the model is served through a cloud-guided inference API (e.g., OpenRouter) that natively enforces JSON schema constraints without relying on the outlines compilation backend.}
\end{quote}

This criterion is consistent with empirical observations from the main scaling sweeps (Appendix~\ref{sec:results}), where even 7B--8B frontier models (Qwen2.5-7B, Llama-3.1-8B) required significant engineering effort to produce well-structured JSON outputs within the 3,000-agent persona generation loop. The 2B--7B regime represents a capability floor below which the emergent persona-diversity necessary for meaningful crowd simulation cannot be reliably initialized.




\section{Extended Future Work: Testbed Follow-Ups}
\label{sec:extended_future_work}

These items extend the original three-hypothesis testbed rather than the paper's central sharpness-loss/\allowbreak instrument-reliability finding; Section~\ref{sec:conclusion} lists the two prerequisites (restoring evaluator schema enforcement in place, and $k \geq 3$ repeated evaluator passes) that should be done first.

\begin{enumerate}
    \item[(iii)] \textbf{Isolate length from what co-varies with it.} Section~\ref{sec:volume_mechanism} establishes a monotone within-campaign dose--response between \texttt{evidence\_text} length and Evaluator J's flat rate ($5\% \to 25\% \to 65\%$), and shows the effect is absent on the same inputs under other evaluators. What it cannot separate is raw character count from the speaker count, claim count and positional conflict that rise with it in this corpus. Truncating and padding \texttt{evidence\_text} to matched lengths---holding speaker count fixed in one arm and length fixed in another---would settle which is operative. This is now a confirmatory experiment on a measured effect rather than a test of a conjecture.
    \item[(iv)] \textbf{Re-run Condition A independently for every model}, so that lift becomes a within-model contrast and the original Prediction Lift hypothesis becomes answerable.
    \item[(v)] \textbf{Repair the convergence telemetry}: guarantee a probed distribution at every checkpoint, fail loudly rather than emitting the $0.758277$ placeholder, and replace the vacuous monotonicity flag with a test requiring a minimum count of valid readings.
    \item[(vi)] \textbf{Report flat-unit rate as a first-class diagnostic} in every campaign, and exclude or explicitly impute those units rather than scoring them as confident errors.
    \item[(vii)] \textbf{Export per-unit action counts} into released summaries so the interaction analysis becomes reproducible from artifacts rather than cluster-side queries.
    \item[(viii)] \textbf{Decompose the Brier score} into reliability, resolution and uncertainty so sharpness effects are separable from skill effects in general, not only in the two-condition comparison used here.
    \item[(ix)] \textbf{Replace the performance-derived contamination criterion} with an external one---resolution date relative to each model's documented training cutoff---which the seed corpus already supports and which breaks the circularity of Appendix~\ref{sec:contamination}.
    \item[(x)] \textbf{Expand the catalogue} to $E \geq 92$ post-cutoff events for the power required by the DM test, and add replicates to the configuration sweeps so scaling claims acquire variance estimates.
    \item[(xi)] \textbf{Recover Qwen2.5-14B as a fourth campaign. Done in full.} \texttt{ground\_truth} was backfilled from the authoritative catalogue (\texttt{data/events\_raw.json}, not the disagreeing copy under \texttt{ECN-HPC-DEPLOY/data/}); 6 of 30 events were additionally found to have been simulated against the wrong topic and excluded (2026-08-14), leaving a scorable 24-event campaign. Those 6 events (T3, T4, T5, T6, C6, T8) were re-simulated against the correct dossiers and rescored 2026-08-21 (Appendix~\ref{sec:fourth_model_status}), after a first re-simulation attempt reproduced the identical stale-catalogue bug and had to be redone; Qwen2.5-14B is now compared against the other three models on the full 20-event globally-clean set all four share (Section~\ref{sec:four_model_extension}), not the 17-event subset an earlier draft used. The remaining gap is the action-density scan: the six recovered events' simulation artifacts (persona files, \texttt{.db} traces) were transferred off the cluster and released on 2026-08-26, but the 2026-07-27 scan they would correct read 180 databases that no longer exist, so it cannot be repeated as taken.
    \item[(xii)] \textbf{Redesign the 7-dimension rubric} with finer-grained ordinal scales and event-specific calibration prompts targeting $\kappa \geq 0.6$, re-measured only after item (i), since the current $\kappa$ figures conflate rubric design with the schema-disabling defect.
    \item[(xiii)] \textbf{Investigate social-interaction forcing mechanisms} to activate dormant architectures, and revisit opinion-dynamics modelling---DeGroot, Friedkin--Johnsen, or otherwise---only once item (v) makes convergence measurable.
    \item[(xiv)] \textbf{Complete the single-agent baseline. Partly done.} Section~\ref{sec:cot_baseline} runs it, but only two of its four arms use the campaign's own checkpoint, and both of those reach the model through a hosted API at FP16 rather than the cluster's quantised serving. Neither Mistral-7B-Instruct-v0.1 nor Qwen2.5-14B is served by the API provider used here, so the two remaining arms must be run against the cluster's own vLLM and Ollama endpoints---which would also remove the precision and serving-stack differences from the matched arms, and cost no external inference. Separately, the contrast needs $E = 59$ events for $80\%$ power at the effect size observed (against $30\%$ at $E = 17$); this is a weaker requirement than item~(x)'s $E \geq 92$, which is set by a different and smaller effect.
\end{enumerate}

\section*{Data Availability}

All quantitative claims in this paper are derived from artifacts released in \ecnarchive. The archive is arranged in reading order and each directory carries its own \texttt{README}; \texttt{PROVENANCE.md} records which claims survived the audit of Section~\ref{sec:evaluator_replication} and which were withdrawn, \texttt{DATA\_DICTIONARY.md} defines every field of every file format, and \texttt{REPRODUCE.md} gives the exact commands that regenerate the tables and figures below. Those commands were verified end to end against the archive: all five analyses reproduce output byte-identical to the released files. For each of the four campaigns we release the simulation transcripts and per-unit scored records: \texttt{event\_results.json} (the July 2026 pass, Evaluator J, DeepSeek-R1-14B), \texttt{event\_results\_evaluatorB\_openrouter.json} together with \texttt{event\_results\_evaluatorA\_repeat\{1,2,3\}.json} (the four independent Evaluator~A reads of Section~\ref{sec:four_model_extension}), and \texttt{event\_results\_evaluatorC\_gptlunapro.json} (the independent Evaluator~C read). The two evaluator passes are kept in separate files rather than merged, because the derived metrics stored in each are only valid for that pass's own probability vectors. The rollup behind Section~\ref{sec:four_model_extension} and Table~\ref{tab:four_model_clean} is released as \texttt{reproducibility\_analysis\_master.json}, and the run-integrity audits of Appendices~\ref{sec:persona_audit}--\ref{sec:evidence_depth} read the per-unit SQLite \texttt{trace} databases. These are released in full for Llama-3.1-8B, Mistral-7B and Qwen2.5-7B ($120$ databases per model). For Qwen2.5-14B they are released for the six events re-simulated 2026-08-21 ($18$ units, $24$ databases), together with a three-unit fragment of the original 2026-07-27 campaign (event~C15); the transcripts for that campaign's remaining $21$ events were removed by the cluster's scratch-filesystem retention policy before they could be transferred, and are unrecoverable. Their scored records are unaffected and are released in full. One consequence is documented rather than repaired: the all-campaign action-density counts of Table~\ref{tab:action_density_30event} come from a 2026-07-27 cluster-side scan of Qwen2.5-14B's then-30-event campaign, 6 of which were later found to have been simulated under the wrong topic (\S\ref{sec:fourth_model_status}); that scan could not be redone against the corrected traces, and carries the reproducibility caveat noted there.

\begin{sloppypar}
The three post-campaign experiments are released with the generator, the analysis and the figure script for each, so that every number in them can be recomputed from the raw runs rather than taken from a summary. The cross-evaluator degeneracy-transfer probe of Section~\ref{sec:volume_mechanism} is \texttt{run\_\allowbreak degeneracy\_\allowbreak transfer.py} and \texttt{analyze\_\allowbreak degeneracy\_\allowbreak transfer.py} over \texttt{degeneracy\_\allowbreak transfer.json}. The $2\times4$ truncation grid is \texttt{build\_\allowbreak truncation\_\allowbreak inputs.py}, \texttt{run\_\allowbreak truncation\_\allowbreak grid.py} and \texttt{analyze\_\allowbreak truncation\_\allowbreak grid.py} over \texttt{event\_\allowbreak results\_\allowbreak trunc\_*.json} and \texttt{truncation\_\allowbreak results.json}. The single-agent baseline of Section~\ref{sec:cot_baseline} is \texttt{run\_\allowbreak single\_\allowbreak agent\_\allowbreak cot\_\allowbreak baseline.py}, producing \texttt{single\_\allowbreak agent\_\allowbreak cot\_\allowbreak baseline\_\allowbreak results.json} ($240$ runs, all fields retained except the chain-of-thought text, which is truncated to $300$ characters per run), analysed by \texttt{analyze\_\allowbreak cot\_\allowbreak baseline.py} into \texttt{cot\_\allowbreak baseline\_\allowbreak analysis.json}. That last file carries all three CoT aggregation variants and both pooling scopes, not only the ones tabulated here, so a reader can check the sensitivity of \S\ref{sec:cot_baseline} directly.
\end{sloppypar}

The event catalogue, seed dossiers, injection bank and scoring rubric are released so that the benchmark can be run against new models, alongside the simulation engine, the HPC protocol runner and the SLURM job scripts as executed. The archive also retains the superseded and invalid artifacts the audit produced---the defective repair of \S\ref{sec:outage}, the wrong-topic re-simulation of \S\ref{sec:fourth_model_status}, and the two divergent event catalogues whose disagreement caused both---so that every defect reported here is independently checkable rather than merely asserted.

\ifecnanonymous
To preserve anonymity during review, the archive is not linked here; it is available to reviewers through the venue, and the permanent public record will be linked in the camera-ready version.
\else
The archive is permanently available at \url{\ecnarchiveurl}.
\fi

\bibliographystyle{IEEEtran}
\begin{thebibliography}{10}

\bibitem{mmlu}
Dan Hendrycks, Collin Burns, Steven Basart, Andy Zou, Mantas Mazeika, Dawn Song, and Jacob Steinhardt.
\newblock Measuring massive multitask language understanding.
\newblock {\em arXiv preprint arXiv:2009.03300}, 2020.

\bibitem{swebench}
Carlos~E. Jimenez, John Yang, Alexander Wettig, Shunyu Yao, Kaiming Zhang, and Karthik Narasimhan.
\newblock SWE-bench: Can language models resolve real-world GitHub issues?
\newblock {\em arXiv preprint arXiv:2310.06770}, 2023.

\bibitem{surowiecki2005}
James Surowiecki.
\newblock {\em The Wisdom of Crowds}.
\newblock Anchor Books, 2005.

\bibitem{gneiting2007}
Tilmann Gneiting and Adrian~E. Raftery.
\newblock Strictly proper scoring rules, prediction, and estimation.
\newblock {\em Journal of the American Statistical Association}, 102(477):359--378, 2007.

\bibitem{du2023debate}
Yilun Du, Shuang Li, Antonio Torralba, Joshua~B. Tenenbaum, and Igor Mordatch.
\newblock Improving factuality and reasoning in language models through multiagent debate.
\newblock {\em arXiv preprint arXiv:2305.14325}, 2023.

\bibitem{sharma2023sycophancy}
Mrinank Sharma, Meg Tong, Tomasz Korbak, David Duvenaud, Amanda Askell, Samuel~R. Bowman, et~al.
\newblock Towards understanding sycophancy in language models.
\newblock {\em arXiv preprint arXiv:2310.13548}, 2023.

\bibitem{hle2024}
Scale~AI and Center for AI Safety.
\newblock Humanity's Last Exam: A benchmark for expert-level AI alignment and knowledge.
\newblock {\em Technical Report}, 2024.

\bibitem{webarena2023}
Shuyan Zhou, Frank~F. Xu, Hao Zhu, Zhoujun Cheng, Jiayi Chen, and Graham Neubig.
\newblock WebArena: A realistic web environment for building autonomous agents.
\newblock {\em arXiv preprint arXiv:2307.13854}, 2023.

\bibitem{park2023generative}
Joon~Sung Park, Joseph~C. O'Brien, Carrie~J. Cai, Meredith~Ringel Morris, Percy Liang, and Michael~S. Bernstein.
\newblock Generative agents: Interactive simulacra of human behavior.
\newblock In {\em Proceedings of the 36th Annual ACM Symposium on User Interface Software and Technology}, 2023.

\bibitem{oasis2024}
CAMEL-AI Community.
\newblock OASIS: Open Agent Social Interaction Simulations with One Million Agents.
\newblock {\em arXiv preprint arXiv:2411.11581}, 2024.

\bibitem{mirofish2024}
666ghj.
\newblock MiroFish: Social simulation pipeline for LLM prediction.
\newblock {\em GitHub Repository: https://github.com/666ghj/MiroFish}, 2024.

\bibitem{graphrag2024}
Darren Edge, Ha Trinh, Cheng Cheng, Joshua Bradley, Alex Chao, Apurv Mody, Chris Porteous, and Jonathan Larson.
\newblock From local to global: A Graph RAG approach to query-focused summarization.
\newblock {\em arXiv preprint arXiv:2404.16130}, 2024.

\bibitem{graphiti2024}
Zep AI.
\newblock Graphiti: An open-source temporal context graph engine for agent memory.
\newblock {\em GitHub Repository: https://github.com/zep-ai/graphiti}, 2024.

\bibitem{vllm}
Woosuk Kwon, Zhuohan Li, Siyuan Zhuang, Sheng Shen, Ion Stoica, and Joseph~E. Gonzalez.
\newblock Efficient memory management for large language model serving with PagedAttention.
\newblock In {\em Proceedings of the 29th Symposium on Operating Systems Principles}, 2023.

\bibitem{vllm_guided}
vLLM Contributors.
\newblock Structured outputs: Guided decoding via outlines and xgrammar backends.
\newblock {\em vLLM Documentation, https://docs.vllm.ai/en/stable/usage/structured\_outputs.html}, 2024.


\bibitem{zheng2023judging}
Lianmin Zheng, Wei-Lin Chiang, Ying Sheng, Siyuan Zhuang, Zhanghao Wu, Yonghao Zhuang, Zi Lin, Zhuohan Li, Dacheng Li, Eric~P. Xing, Hao Zhang, Joseph~E. Gonzalez, and Ion Stoica.
\newblock Judging LLM-as-a-judge with MT-Bench and Chatbot Arena.
\newblock In {\em Advances in Neural Information Processing Systems (NeurIPS), Datasets and Benchmarks Track}, 2023.

\bibitem{oren2024proving}
Yonatan Oren, Nicole Meister, Niladri Chatterji, Faisal Ladhak, and Tatsunori~B. Hashimoto.
\newblock Proving test set contamination in black-box language models.
\newblock In {\em International Conference on Learning Representations (ICLR)}, 2024.

\bibitem{karger2024forecastbench}
Ezra Karger, Houtan Bastani, Chen Yueh-Han, Zachary Jacobs, Danny Halawi, Fred Zhang, and Philip~E. Tetlock.
\newblock ForecastBench: A dynamic benchmark of AI forecasting capabilities.
\newblock {\em arXiv preprint arXiv:2409.19839}, 2024.

\bibitem{halawi2024approaching}
Danny Halawi, Fred Zhang, Chen Yueh-Han, and Jacob Steinhardt.
\newblock Approaching human-level forecasting with language models.
\newblock In {\em Advances in Neural Information Processing Systems (NeurIPS)}, 2024.

\end{thebibliography}

\end{document}
